\documentclass[11pt]{article}

    \usepackage[T2A]{fontenc}
\usepackage[utf8]{inputenc}
\usepackage[english, russian]{babel}

\usepackage[breakable]{tcolorbox}
    \usepackage{parskip} % Stop auto-indenting (to mimic markdown behaviour)
    
    \usepackage{iftex}
    \ifPDFTeX
    	\usepackage[T1]{fontenc}
    	\usepackage{mathpazo}
    \else
    	\usepackage{fontspec}
    \fi

    % Basic figure setup, for now with no caption control since it's done
    % automatically by Pandoc (which extracts ![](path) syntax from Markdown).
    \usepackage{graphicx}
    % Maintain compatibility with old templates. Remove in nbconvert 6.0
    \let\Oldincludegraphics\includegraphics
    % Ensure that by default, figures have no caption (until we provide a
    % proper Figure object with a Caption API and a way to capture that
    % in the conversion process - todo).
    \usepackage{caption}
    \DeclareCaptionFormat{nocaption}{}
    \captionsetup{format=nocaption,aboveskip=0pt,belowskip=0pt}

    \usepackage{float}
    \floatplacement{figure}{H} % forces figures to be placed at the correct location
    \usepackage{xcolor} % Allow colors to be defined
    \usepackage{enumerate} % Needed for markdown enumerations to work
    \usepackage{geometry} % Used to adjust the document margins
    \usepackage{amsmath} % Equations
    \usepackage{amssymb} % Equations
    \usepackage{textcomp} % defines textquotesingle
    % Hack from http://tex.stackexchange.com/a/47451/13684:
    \AtBeginDocument{%
        \def\PYZsq{\textquotesingle}% Upright quotes in Pygmentized code
    }
    \usepackage{upquote} % Upright quotes for verbatim code
    \usepackage{eurosym} % defines \euro
    \usepackage[mathletters]{ucs} % Extended unicode (utf-8) support
    \usepackage{fancyvrb} % verbatim replacement that allows latex
    \usepackage{grffile} % extends the file name processing of package graphics 
                         % to support a larger range
    \makeatletter % fix for old versions of grffile with XeLaTeX
    \@ifpackagelater{grffile}{2019/11/01}
    {
      % Do nothing on new versions
    }
    {
      \def\Gread@@xetex#1{%
        \IfFileExists{"\Gin@base".bb}%
        {\Gread@eps{\Gin@base.bb}}%
        {\Gread@@xetex@aux#1}%
      }
    }
    \makeatother
    \usepackage[Export]{adjustbox} % Used to constrain images to a maximum size
    \adjustboxset{max size={0.9\linewidth}{0.9\paperheight}}

    % The hyperref package gives us a pdf with properly built
    % internal navigation ('pdf bookmarks' for the table of contents,
    % internal cross-reference links, web links for URLs, etc.)
    \usepackage{hyperref}
    % The default LaTeX title has an obnoxious amount of whitespace. By default,
    % titling removes some of it. It also provides customization options.
    \usepackage{titling}
    \usepackage{longtable} % longtable support required by pandoc >1.10
    \usepackage{booktabs}  % table support for pandoc > 1.12.2
    \usepackage[inline]{enumitem} % IRkernel/repr support (it uses the enumerate* environment)
    \usepackage[normalem]{ulem} % ulem is needed to support strikethroughs (\sout)
                                % normalem makes italics be italics, not underlines
    \usepackage{mathrsfs}
    
\usepackage{wrapfig}
\usepackage[rightcaption]{sidecap}
\providecommand{\keywords}[1]{\textbf{\textit{Keywords:}} #1}

\author{A. Yu. Drozdov}

    
    % Colors for the hyperref package
    \definecolor{urlcolor}{rgb}{0,.145,.698}
    \definecolor{linkcolor}{rgb}{.71,0.21,0.01}
    \definecolor{citecolor}{rgb}{.12,.54,.11}

    % ANSI colors
    \definecolor{ansi-black}{HTML}{3E424D}
    \definecolor{ansi-black-intense}{HTML}{282C36}
    \definecolor{ansi-red}{HTML}{E75C58}
    \definecolor{ansi-red-intense}{HTML}{B22B31}
    \definecolor{ansi-green}{HTML}{00A250}
    \definecolor{ansi-green-intense}{HTML}{007427}
    \definecolor{ansi-yellow}{HTML}{DDB62B}
    \definecolor{ansi-yellow-intense}{HTML}{B27D12}
    \definecolor{ansi-blue}{HTML}{208FFB}
    \definecolor{ansi-blue-intense}{HTML}{0065CA}
    \definecolor{ansi-magenta}{HTML}{D160C4}
    \definecolor{ansi-magenta-intense}{HTML}{A03196}
    \definecolor{ansi-cyan}{HTML}{60C6C8}
    \definecolor{ansi-cyan-intense}{HTML}{258F8F}
    \definecolor{ansi-white}{HTML}{C5C1B4}
    \definecolor{ansi-white-intense}{HTML}{A1A6B2}
    \definecolor{ansi-default-inverse-fg}{HTML}{FFFFFF}
    \definecolor{ansi-default-inverse-bg}{HTML}{000000}

    % common color for the border for error outputs.
    \definecolor{outerrorbackground}{HTML}{FFDFDF}

    % commands and environments needed by pandoc snippets
    % extracted from the output of `pandoc -s`
    \providecommand{\tightlist}{%
      \setlength{\itemsep}{0pt}\setlength{\parskip}{0pt}}
    \DefineVerbatimEnvironment{Highlighting}{Verbatim}{commandchars=\\\{\}}
    % Add ',fontsize=\small' for more characters per line
    \newenvironment{Shaded}{}{}
    \newcommand{\KeywordTok}[1]{\textcolor[rgb]{0.00,0.44,0.13}{\textbf{{#1}}}}
    \newcommand{\DataTypeTok}[1]{\textcolor[rgb]{0.56,0.13,0.00}{{#1}}}
    \newcommand{\DecValTok}[1]{\textcolor[rgb]{0.25,0.63,0.44}{{#1}}}
    \newcommand{\BaseNTok}[1]{\textcolor[rgb]{0.25,0.63,0.44}{{#1}}}
    \newcommand{\FloatTok}[1]{\textcolor[rgb]{0.25,0.63,0.44}{{#1}}}
    \newcommand{\CharTok}[1]{\textcolor[rgb]{0.25,0.44,0.63}{{#1}}}
    \newcommand{\StringTok}[1]{\textcolor[rgb]{0.25,0.44,0.63}{{#1}}}
    \newcommand{\CommentTok}[1]{\textcolor[rgb]{0.38,0.63,0.69}{\textit{{#1}}}}
    \newcommand{\OtherTok}[1]{\textcolor[rgb]{0.00,0.44,0.13}{{#1}}}
    \newcommand{\AlertTok}[1]{\textcolor[rgb]{1.00,0.00,0.00}{\textbf{{#1}}}}
    \newcommand{\FunctionTok}[1]{\textcolor[rgb]{0.02,0.16,0.49}{{#1}}}
    \newcommand{\RegionMarkerTok}[1]{{#1}}
    \newcommand{\ErrorTok}[1]{\textcolor[rgb]{1.00,0.00,0.00}{\textbf{{#1}}}}
    \newcommand{\NormalTok}[1]{{#1}}
    
    % Additional commands for more recent versions of Pandoc
    \newcommand{\ConstantTok}[1]{\textcolor[rgb]{0.53,0.00,0.00}{{#1}}}
    \newcommand{\SpecialCharTok}[1]{\textcolor[rgb]{0.25,0.44,0.63}{{#1}}}
    \newcommand{\VerbatimStringTok}[1]{\textcolor[rgb]{0.25,0.44,0.63}{{#1}}}
    \newcommand{\SpecialStringTok}[1]{\textcolor[rgb]{0.73,0.40,0.53}{{#1}}}
    \newcommand{\ImportTok}[1]{{#1}}
    \newcommand{\DocumentationTok}[1]{\textcolor[rgb]{0.73,0.13,0.13}{\textit{{#1}}}}
    \newcommand{\AnnotationTok}[1]{\textcolor[rgb]{0.38,0.63,0.69}{\textbf{\textit{{#1}}}}}
    \newcommand{\CommentVarTok}[1]{\textcolor[rgb]{0.38,0.63,0.69}{\textbf{\textit{{#1}}}}}
    \newcommand{\VariableTok}[1]{\textcolor[rgb]{0.10,0.09,0.49}{{#1}}}
    \newcommand{\ControlFlowTok}[1]{\textcolor[rgb]{0.00,0.44,0.13}{\textbf{{#1}}}}
    \newcommand{\OperatorTok}[1]{\textcolor[rgb]{0.40,0.40,0.40}{{#1}}}
    \newcommand{\BuiltInTok}[1]{{#1}}
    \newcommand{\ExtensionTok}[1]{{#1}}
    \newcommand{\PreprocessorTok}[1]{\textcolor[rgb]{0.74,0.48,0.00}{{#1}}}
    \newcommand{\AttributeTok}[1]{\textcolor[rgb]{0.49,0.56,0.16}{{#1}}}
    \newcommand{\InformationTok}[1]{\textcolor[rgb]{0.38,0.63,0.69}{\textbf{\textit{{#1}}}}}
    \newcommand{\WarningTok}[1]{\textcolor[rgb]{0.38,0.63,0.69}{\textbf{\textit{{#1}}}}}
    
    
    % Define a nice break command that doesn't care if a line doesn't already
    % exist.
    \def\br{\hspace*{\fill} \\* }
    % Math Jax compatibility definitions
    \def\gt{>}
    \def\lt{<}
    \let\Oldtex\TeX
    \let\Oldlatex\LaTeX
    \renewcommand{\TeX}{\textrm{\Oldtex}}
    \renewcommand{\LaTeX}{\textrm{\Oldlatex}}
    % Document parameters
    % Document title
    \title{On\_the\_applicability\_of\_Lorentz\_calibration\_to\_Nikolaev-s\_electrodynamics}
    
    
    
    
    
% Pygments definitions
\makeatletter
\def\PY@reset{\let\PY@it=\relax \let\PY@bf=\relax%
    \let\PY@ul=\relax \let\PY@tc=\relax%
    \let\PY@bc=\relax \let\PY@ff=\relax}
\def\PY@tok#1{\csname PY@tok@#1\endcsname}
\def\PY@toks#1+{\ifx\relax#1\empty\else%
    \PY@tok{#1}\expandafter\PY@toks\fi}
\def\PY@do#1{\PY@bc{\PY@tc{\PY@ul{%
    \PY@it{\PY@bf{\PY@ff{#1}}}}}}}
\def\PY#1#2{\PY@reset\PY@toks#1+\relax+\PY@do{#2}}

\@namedef{PY@tok@w}{\def\PY@tc##1{\textcolor[rgb]{0.73,0.73,0.73}{##1}}}
\@namedef{PY@tok@c}{\let\PY@it=\textit\def\PY@tc##1{\textcolor[rgb]{0.25,0.50,0.50}{##1}}}
\@namedef{PY@tok@cp}{\def\PY@tc##1{\textcolor[rgb]{0.74,0.48,0.00}{##1}}}
\@namedef{PY@tok@k}{\let\PY@bf=\textbf\def\PY@tc##1{\textcolor[rgb]{0.00,0.50,0.00}{##1}}}
\@namedef{PY@tok@kp}{\def\PY@tc##1{\textcolor[rgb]{0.00,0.50,0.00}{##1}}}
\@namedef{PY@tok@kt}{\def\PY@tc##1{\textcolor[rgb]{0.69,0.00,0.25}{##1}}}
\@namedef{PY@tok@o}{\def\PY@tc##1{\textcolor[rgb]{0.40,0.40,0.40}{##1}}}
\@namedef{PY@tok@ow}{\let\PY@bf=\textbf\def\PY@tc##1{\textcolor[rgb]{0.67,0.13,1.00}{##1}}}
\@namedef{PY@tok@nb}{\def\PY@tc##1{\textcolor[rgb]{0.00,0.50,0.00}{##1}}}
\@namedef{PY@tok@nf}{\def\PY@tc##1{\textcolor[rgb]{0.00,0.00,1.00}{##1}}}
\@namedef{PY@tok@nc}{\let\PY@bf=\textbf\def\PY@tc##1{\textcolor[rgb]{0.00,0.00,1.00}{##1}}}
\@namedef{PY@tok@nn}{\let\PY@bf=\textbf\def\PY@tc##1{\textcolor[rgb]{0.00,0.00,1.00}{##1}}}
\@namedef{PY@tok@ne}{\let\PY@bf=\textbf\def\PY@tc##1{\textcolor[rgb]{0.82,0.25,0.23}{##1}}}
\@namedef{PY@tok@nv}{\def\PY@tc##1{\textcolor[rgb]{0.10,0.09,0.49}{##1}}}
\@namedef{PY@tok@no}{\def\PY@tc##1{\textcolor[rgb]{0.53,0.00,0.00}{##1}}}
\@namedef{PY@tok@nl}{\def\PY@tc##1{\textcolor[rgb]{0.63,0.63,0.00}{##1}}}
\@namedef{PY@tok@ni}{\let\PY@bf=\textbf\def\PY@tc##1{\textcolor[rgb]{0.60,0.60,0.60}{##1}}}
\@namedef{PY@tok@na}{\def\PY@tc##1{\textcolor[rgb]{0.49,0.56,0.16}{##1}}}
\@namedef{PY@tok@nt}{\let\PY@bf=\textbf\def\PY@tc##1{\textcolor[rgb]{0.00,0.50,0.00}{##1}}}
\@namedef{PY@tok@nd}{\def\PY@tc##1{\textcolor[rgb]{0.67,0.13,1.00}{##1}}}
\@namedef{PY@tok@s}{\def\PY@tc##1{\textcolor[rgb]{0.73,0.13,0.13}{##1}}}
\@namedef{PY@tok@sd}{\let\PY@it=\textit\def\PY@tc##1{\textcolor[rgb]{0.73,0.13,0.13}{##1}}}
\@namedef{PY@tok@si}{\let\PY@bf=\textbf\def\PY@tc##1{\textcolor[rgb]{0.73,0.40,0.53}{##1}}}
\@namedef{PY@tok@se}{\let\PY@bf=\textbf\def\PY@tc##1{\textcolor[rgb]{0.73,0.40,0.13}{##1}}}
\@namedef{PY@tok@sr}{\def\PY@tc##1{\textcolor[rgb]{0.73,0.40,0.53}{##1}}}
\@namedef{PY@tok@ss}{\def\PY@tc##1{\textcolor[rgb]{0.10,0.09,0.49}{##1}}}
\@namedef{PY@tok@sx}{\def\PY@tc##1{\textcolor[rgb]{0.00,0.50,0.00}{##1}}}
\@namedef{PY@tok@m}{\def\PY@tc##1{\textcolor[rgb]{0.40,0.40,0.40}{##1}}}
\@namedef{PY@tok@gh}{\let\PY@bf=\textbf\def\PY@tc##1{\textcolor[rgb]{0.00,0.00,0.50}{##1}}}
\@namedef{PY@tok@gu}{\let\PY@bf=\textbf\def\PY@tc##1{\textcolor[rgb]{0.50,0.00,0.50}{##1}}}
\@namedef{PY@tok@gd}{\def\PY@tc##1{\textcolor[rgb]{0.63,0.00,0.00}{##1}}}
\@namedef{PY@tok@gi}{\def\PY@tc##1{\textcolor[rgb]{0.00,0.63,0.00}{##1}}}
\@namedef{PY@tok@gr}{\def\PY@tc##1{\textcolor[rgb]{1.00,0.00,0.00}{##1}}}
\@namedef{PY@tok@ge}{\let\PY@it=\textit}
\@namedef{PY@tok@gs}{\let\PY@bf=\textbf}
\@namedef{PY@tok@gp}{\let\PY@bf=\textbf\def\PY@tc##1{\textcolor[rgb]{0.00,0.00,0.50}{##1}}}
\@namedef{PY@tok@go}{\def\PY@tc##1{\textcolor[rgb]{0.53,0.53,0.53}{##1}}}
\@namedef{PY@tok@gt}{\def\PY@tc##1{\textcolor[rgb]{0.00,0.27,0.87}{##1}}}
\@namedef{PY@tok@err}{\def\PY@bc##1{{\setlength{\fboxsep}{\string -\fboxrule}\fcolorbox[rgb]{1.00,0.00,0.00}{1,1,1}{\strut ##1}}}}
\@namedef{PY@tok@kc}{\let\PY@bf=\textbf\def\PY@tc##1{\textcolor[rgb]{0.00,0.50,0.00}{##1}}}
\@namedef{PY@tok@kd}{\let\PY@bf=\textbf\def\PY@tc##1{\textcolor[rgb]{0.00,0.50,0.00}{##1}}}
\@namedef{PY@tok@kn}{\let\PY@bf=\textbf\def\PY@tc##1{\textcolor[rgb]{0.00,0.50,0.00}{##1}}}
\@namedef{PY@tok@kr}{\let\PY@bf=\textbf\def\PY@tc##1{\textcolor[rgb]{0.00,0.50,0.00}{##1}}}
\@namedef{PY@tok@bp}{\def\PY@tc##1{\textcolor[rgb]{0.00,0.50,0.00}{##1}}}
\@namedef{PY@tok@fm}{\def\PY@tc##1{\textcolor[rgb]{0.00,0.00,1.00}{##1}}}
\@namedef{PY@tok@vc}{\def\PY@tc##1{\textcolor[rgb]{0.10,0.09,0.49}{##1}}}
\@namedef{PY@tok@vg}{\def\PY@tc##1{\textcolor[rgb]{0.10,0.09,0.49}{##1}}}
\@namedef{PY@tok@vi}{\def\PY@tc##1{\textcolor[rgb]{0.10,0.09,0.49}{##1}}}
\@namedef{PY@tok@vm}{\def\PY@tc##1{\textcolor[rgb]{0.10,0.09,0.49}{##1}}}
\@namedef{PY@tok@sa}{\def\PY@tc##1{\textcolor[rgb]{0.73,0.13,0.13}{##1}}}
\@namedef{PY@tok@sb}{\def\PY@tc##1{\textcolor[rgb]{0.73,0.13,0.13}{##1}}}
\@namedef{PY@tok@sc}{\def\PY@tc##1{\textcolor[rgb]{0.73,0.13,0.13}{##1}}}
\@namedef{PY@tok@dl}{\def\PY@tc##1{\textcolor[rgb]{0.73,0.13,0.13}{##1}}}
\@namedef{PY@tok@s2}{\def\PY@tc##1{\textcolor[rgb]{0.73,0.13,0.13}{##1}}}
\@namedef{PY@tok@sh}{\def\PY@tc##1{\textcolor[rgb]{0.73,0.13,0.13}{##1}}}
\@namedef{PY@tok@s1}{\def\PY@tc##1{\textcolor[rgb]{0.73,0.13,0.13}{##1}}}
\@namedef{PY@tok@mb}{\def\PY@tc##1{\textcolor[rgb]{0.40,0.40,0.40}{##1}}}
\@namedef{PY@tok@mf}{\def\PY@tc##1{\textcolor[rgb]{0.40,0.40,0.40}{##1}}}
\@namedef{PY@tok@mh}{\def\PY@tc##1{\textcolor[rgb]{0.40,0.40,0.40}{##1}}}
\@namedef{PY@tok@mi}{\def\PY@tc##1{\textcolor[rgb]{0.40,0.40,0.40}{##1}}}
\@namedef{PY@tok@il}{\def\PY@tc##1{\textcolor[rgb]{0.40,0.40,0.40}{##1}}}
\@namedef{PY@tok@mo}{\def\PY@tc##1{\textcolor[rgb]{0.40,0.40,0.40}{##1}}}
\@namedef{PY@tok@ch}{\let\PY@it=\textit\def\PY@tc##1{\textcolor[rgb]{0.25,0.50,0.50}{##1}}}
\@namedef{PY@tok@cm}{\let\PY@it=\textit\def\PY@tc##1{\textcolor[rgb]{0.25,0.50,0.50}{##1}}}
\@namedef{PY@tok@cpf}{\let\PY@it=\textit\def\PY@tc##1{\textcolor[rgb]{0.25,0.50,0.50}{##1}}}
\@namedef{PY@tok@c1}{\let\PY@it=\textit\def\PY@tc##1{\textcolor[rgb]{0.25,0.50,0.50}{##1}}}
\@namedef{PY@tok@cs}{\let\PY@it=\textit\def\PY@tc##1{\textcolor[rgb]{0.25,0.50,0.50}{##1}}}

\def\PYZbs{\char`\\}
\def\PYZus{\char`\_}
\def\PYZob{\char`\{}
\def\PYZcb{\char`\}}
\def\PYZca{\char`\^}
\def\PYZam{\char`\&}
\def\PYZlt{\char`\<}
\def\PYZgt{\char`\>}
\def\PYZsh{\char`\#}
\def\PYZpc{\char`\%}
\def\PYZdl{\char`\$}
\def\PYZhy{\char`\-}
\def\PYZsq{\char`\'}
\def\PYZdq{\char`\"}
\def\PYZti{\char`\~}
% for compatibility with earlier versions
\def\PYZat{@}
\def\PYZlb{[}
\def\PYZrb{]}
\makeatother


    % For linebreaks inside Verbatim environment from package fancyvrb. 
    \makeatletter
        \newbox\Wrappedcontinuationbox 
        \newbox\Wrappedvisiblespacebox 
        \newcommand*\Wrappedvisiblespace {\textcolor{red}{\textvisiblespace}} 
        \newcommand*\Wrappedcontinuationsymbol {\textcolor{red}{\llap{\tiny$\m@th\hookrightarrow$}}} 
        \newcommand*\Wrappedcontinuationindent {3ex } 
        \newcommand*\Wrappedafterbreak {\kern\Wrappedcontinuationindent\copy\Wrappedcontinuationbox} 
        % Take advantage of the already applied Pygments mark-up to insert 
        % potential linebreaks for TeX processing. 
        %        {, <, #, %, $, ' and ": go to next line. 
        %        _, }, ^, &, >, - and ~: stay at end of broken line. 
        % Use of \textquotesingle for straight quote. 
        \newcommand*\Wrappedbreaksatspecials {% 
            \def\PYGZus{\discretionary{\char`\_}{\Wrappedafterbreak}{\char`\_}}% 
            \def\PYGZob{\discretionary{}{\Wrappedafterbreak\char`\{}{\char`\{}}% 
            \def\PYGZcb{\discretionary{\char`\}}{\Wrappedafterbreak}{\char`\}}}% 
            \def\PYGZca{\discretionary{\char`\^}{\Wrappedafterbreak}{\char`\^}}% 
            \def\PYGZam{\discretionary{\char`\&}{\Wrappedafterbreak}{\char`\&}}% 
            \def\PYGZlt{\discretionary{}{\Wrappedafterbreak\char`\<}{\char`\<}}% 
            \def\PYGZgt{\discretionary{\char`\>}{\Wrappedafterbreak}{\char`\>}}% 
            \def\PYGZsh{\discretionary{}{\Wrappedafterbreak\char`\#}{\char`\#}}% 
            \def\PYGZpc{\discretionary{}{\Wrappedafterbreak\char`\%}{\char`\%}}% 
            \def\PYGZdl{\discretionary{}{\Wrappedafterbreak\char`\$}{\char`\$}}% 
            \def\PYGZhy{\discretionary{\char`\-}{\Wrappedafterbreak}{\char`\-}}% 
            \def\PYGZsq{\discretionary{}{\Wrappedafterbreak\textquotesingle}{\textquotesingle}}% 
            \def\PYGZdq{\discretionary{}{\Wrappedafterbreak\char`\"}{\char`\"}}% 
            \def\PYGZti{\discretionary{\char`\~}{\Wrappedafterbreak}{\char`\~}}% 
        } 
        % Some characters . , ; ? ! / are not pygmentized. 
        % This macro makes them "active" and they will insert potential linebreaks 
        \newcommand*\Wrappedbreaksatpunct {% 
            \lccode`\~`\.\lowercase{\def~}{\discretionary{\hbox{\char`\.}}{\Wrappedafterbreak}{\hbox{\char`\.}}}% 
            \lccode`\~`\,\lowercase{\def~}{\discretionary{\hbox{\char`\,}}{\Wrappedafterbreak}{\hbox{\char`\,}}}% 
            \lccode`\~`\;\lowercase{\def~}{\discretionary{\hbox{\char`\;}}{\Wrappedafterbreak}{\hbox{\char`\;}}}% 
            \lccode`\~`\:\lowercase{\def~}{\discretionary{\hbox{\char`\:}}{\Wrappedafterbreak}{\hbox{\char`\:}}}% 
            \lccode`\~`\?\lowercase{\def~}{\discretionary{\hbox{\char`\?}}{\Wrappedafterbreak}{\hbox{\char`\?}}}% 
            \lccode`\~`\!\lowercase{\def~}{\discretionary{\hbox{\char`\!}}{\Wrappedafterbreak}{\hbox{\char`\!}}}% 
            \lccode`\~`\/\lowercase{\def~}{\discretionary{\hbox{\char`\/}}{\Wrappedafterbreak}{\hbox{\char`\/}}}% 
            \catcode`\.\active
            \catcode`\,\active 
            \catcode`\;\active
            \catcode`\:\active
            \catcode`\?\active
            \catcode`\!\active
            \catcode`\/\active 
            \lccode`\~`\~ 	
        }
    \makeatother

    \let\OriginalVerbatim=\Verbatim
    \makeatletter
    \renewcommand{\Verbatim}[1][1]{%
        %\parskip\z@skip
        \sbox\Wrappedcontinuationbox {\Wrappedcontinuationsymbol}%
        \sbox\Wrappedvisiblespacebox {\FV@SetupFont\Wrappedvisiblespace}%
        \def\FancyVerbFormatLine ##1{\hsize\linewidth
            \vtop{\raggedright\hyphenpenalty\z@\exhyphenpenalty\z@
                \doublehyphendemerits\z@\finalhyphendemerits\z@
                \strut ##1\strut}%
        }%
        % If the linebreak is at a space, the latter will be displayed as visible
        % space at end of first line, and a continuation symbol starts next line.
        % Stretch/shrink are however usually zero for typewriter font.
        \def\FV@Space {%
            \nobreak\hskip\z@ plus\fontdimen3\font minus\fontdimen4\font
            \discretionary{\copy\Wrappedvisiblespacebox}{\Wrappedafterbreak}
            {\kern\fontdimen2\font}%
        }%
        
        % Allow breaks at special characters using \PYG... macros.
        \Wrappedbreaksatspecials
        % Breaks at punctuation characters . , ; ? ! and / need catcode=\active 	
        \OriginalVerbatim[#1,codes*=\Wrappedbreaksatpunct]%
    }
    \makeatother

    % Exact colors from NB
    \definecolor{incolor}{HTML}{303F9F}
    \definecolor{outcolor}{HTML}{D84315}
    \definecolor{cellborder}{HTML}{CFCFCF}
    \definecolor{cellbackground}{HTML}{F7F7F7}
    
    % prompt
    \makeatletter
    \newcommand{\boxspacing}{\kern\kvtcb@left@rule\kern\kvtcb@boxsep}
    \makeatother
    \newcommand{\prompt}[4]{
        {\ttfamily\llap{{\color{#2}[#3]:\hspace{3pt}#4}}\vspace{-\baselineskip}}
    }
    

    
    % Prevent overflowing lines due to hard-to-break entities
    \sloppy 
    % Setup hyperref package
    \hypersetup{
      breaklinks=true,  % so long urls are correctly broken across lines
      colorlinks=true,
      urlcolor=urlcolor,
      linkcolor=linkcolor,
      citecolor=citecolor,
      }
    % Slightly bigger margins than the latex defaults
    
    \geometry{verbose,tmargin=1in,bmargin=1in,lmargin=1in,rmargin=1in}
    
    

\begin{document}
    
    \maketitle
    
    

    
    Электрический импульс центрально симметричного взрыва плазмы А.Ю.
Дроздов Одной из нерешённых проблем современной электродинамики является
проблема электромагнитного импульса ядерного взрыва. В журнале
Инженерная физика появилась статья {[} CITATION Мен13 \l 1049 {]} в
которой была сделана попытка объяснить это явление в рамках концепции
скалярно-векторного потенциала. Эта концепция по мнению ее автора
предполагает зависимость скалярного потенциала заряда от скорости. Такая
зависимость была получена из анализа законов индукции электрического
поля магнитным и магнитного поля электрическим, записанных с
использованием субстанциональной производной полевых функций в форме,
инвариантной относительно преобразований координат классической физики,
включающих преобразования Галилея. Приводятся также косвенные
экспериментальные данные в пользу справедливости концепции скалярно
векторного потенциала, которые заключаются в появлении электрического
потенциала на сверхпроводящих обмотках и торах при введении в них
постоянного тока {[}CITATION WFE76 \l 1033 {]} {[}CITATION WGV62 \l 1033
{]} {[} CITATION Bak64 \l 1033 {]} {[} CITATION Men93 \l 1033 {]}. В
2015 году появилась публикация {[} CITATION ФФМ \l 1049 {]} с описанием
модельного эксперимента моделирующего возникновение ЭМИ посредством
регистрации импульса возникающего при электрическом взрыве медной
проволочки. Анализу результатов данного модельного эксперимента и
посвящена данная работа. Для начала следует отметить, что попытка
объяснить электромагнитный импульс центрально симметричного взрыва
плазмы в рамках концепции скалярно-векторного потенциала является
ошибочной. Чтобы в этом убедиться достаточно привести рассуждения Ф.Ф.
Менде, применённые при выводе скалярно-векторного потенциала, но для
случая случае центрально-симметричного движения зарядов. Действительно,
пусть мы имеем три ИСО: первая неподвижная, вторая движется со скоростью
\(dv\), а третья со скоростью{[}2dv{]}. В третьей ИСО расположен
заряженный стержень. Во второй появляется прибавка магнитного поля
\(dB\). А в первой прибавка электрического поля \(dE\). Однако в случае
разогретой плазмы мы имеем не единственный движущийся заряженный
стержень. Таких "стержней", движущихся во всех направлениях с различными
скоростями очень много. В случае разогретой плазмы, имеющей сферическую
симметрию, о распределении "стержней" по скоростям можно утверждать
следующее: сколько стержней движется со скоростью {[}2dv{]}, столько же
движется со скоростью {[}-2dv{]}. Таким образом для разбора ситуации нам
достаточно ввести в рассмотрение ещё две ИСО: четвёртая, движущаяся со
скоростью \(-dv\). И пятая, движущаяся со скоростью {[}-2dv{]}. При этом
в пятой ИСО имеется ещё один заряженный стержень. В четвёртой ИСО
появляется прибавка магнитного поля \(-dB\). А в первой ИСО появляется
прибавка электрического поля\(-dE\). Таким образом мы видим, что
согласно предложенного Менде при выводе скалярно-векторного потенциала
механизма в конфигурации сферически симметричной разогретой плазмы
суммарная прибавка электрического поля равна нулю. Причина такого вывода
заключается в том, концепция скалярно-векторного потенциала Менде
основана на «перпендикулярном» механизме образования прибавки
электрического поля. И поэтому его применение может быть распространено
на токовые системы обычного типа, в которых движущиеся электроны имеют
преимущественное направления движения в какую-либо сторону, например,
упомянутые выше сверхпроводящие обмотки и торы. Но эта концепция не
применима к сферически симметрическому электрическому вибратору. Кроме
того нужно иметь ввиду, что при выводе формулы скалярно-векторного
потенциала Менде была допущена принципиальная ошибка, заключающаяся в
том, что формула, полученная для потенциала движущегося заряженного
длинного стержня не может быть преобразована в формулу потенциала
движущегося точечного заряда одним лишь увеличением степени в
знаменателе. При детальном повторении рассуждений Менде, но не с
движущимся заряженным стержнем, а с движущимся точечным зарядом
возникает необходимость учёта того обстоятельства, что угол между
вектором скорости движущегося заряда и радиус вектором от заряда к точке
наблюдения может отличаться от 90 градусов. При учете этого
обстоятельства получается, что формула для скалярно-векторного
потенциала движущегося точечного заряда неверна. Желающий уточнить
формулу скалярно-векторного потенциала, учитывая лишь только
перпендикулярный механизм возникновения прибавки электрического поля
столкнется с необходимостью использования формулы очень сложного вида.
Но если кроме перпендикулярного механизма образования прибавки
электрического поля учесть также и продольный согласно {[} CITATION ГНи
\l 1049 {]} механизм образования прибавки Е, то в результате формула
скалярно-векторного потенциала для движущегося точечного заряда будет
отличаться от формулы Менде тем, при букве \(v\) (скорость) отсутствует
значок перпендикуляр. Вместо перпендикулярной проекции вектора скорости
в формулу входит полная скорость. Расчёт задачи центрально-симметричного
электрического вибратора11Equation Section (Next) Взятая за основу
данного расчёта идея происхождения электромагнитного импульса
центрально-симметричного взрыва была сформулирована В.Н. Фефеловым:
«Поток релятивистских электронов из первичной неравновесной плазмы резко
обгоняет тяжелые ядра и образует сферический конденсатор с меняющейся
емкостью. Это порождает по формулам Николаева переменное СМП (скалярное
магнитное поле) и далее - продольную волну». Для решения задачи
вычисления поля центрально-симметричного взрыва, был использован метод,
применённый Таммом в § 98 «Осциллятор. Запаздывающие потенциалы поля
осциллятора» {[}CITATION Там57 \l 1049 {]} Условия задачи следующие.
Электронный газ, а также газ положительных ионов в результате взрыва
имеют центрально сферически симметричное распределение плотности заряда
и скорости. Предположим, что нам дано распределение плотностей заряда
для электронного газа \({{\rho }_{e}}\left( {{R}_{O}} \right)\) их для
газа положительно заряженных ионов
\({{\rho }_{p}}\left( {{R}_{O}} \right)\) , в зависимости от расстояния
\({{R}_{O}}\) от центра взрыва. Кроме того, исходя из распределения
радиальных проекций скоростей \({{v}_{e}}^{R}\left( {{R}_{O}} \right)\)
и \({{v}_{p}}^{R}\left( {{R}_{O}} \right)\) мы знаем также и радиальные
компоненты плотности тока \({{j}_{e}}^{R}\left( {{R}_{O}} \right)\) и
\({{j}_{p}}^{R}\left( {{R}_{O}} \right)\) Применим скалярный и векторный
запаздывающие потенциалы
{[}\varphi =\frac{1}{\varepsilon }\int{\frac{\rho \left( t-\frac{R}{v} \right)}{R}}~dV'{]}
{[}\{\{A\}\emph{\{x\}\}=\frac{\mu }{c}\int{\frac{{{j}_{x}}\left( t-\frac{R}{v} \right)}{R}}~dV'{]}
Которые в сферической системе координат запишутся следующим образом
{[}\varphi \left( \theta ,\phi ,r,t
\right)=\frac{1}{\varepsilon }\int{\frac{\rho \left( \theta ',\phi ',r',t' \right)}{R}}~dV'{]}
{[}\overrightarrow{{{A}_{B}}}\left( \theta ,\phi ,r,t
\right)=\frac{\mu }{c}\int{\frac{\overrightarrow{j}\left( \theta ',\phi ',r',t' \right)}{R}}~dV'{]}
В данной работе при обозначения векторного потенциала используется
индексация для различения следующих потенциалов
\(\overrightarrow{B}=rot\ \overrightarrow{{{A}_{B}}}\) и
\(\overrightarrow{H}=rot\ \overrightarrow{{{A}_{H}}}\) отличающихся друг
от друга множителем \$\mu \$ : поскольку
{[}\overrightarrow{B}=\mu ~\overrightarrow{H}{]} постольку и
\(\overrightarrow{{{A}_{B}}}=\mu \overrightarrow{{{A}_{H}}}\) Пусть
точка \(O\) -- центр взрыва. Пусть \({{R}_{O}}\) - радиус вектор,
проведённый из \(O\) в точку наблюдения \(P\). Пусть \(R\) - радиус
вектор проведённый из произвольной точки взрыва
\(Q\left( {{\theta }^{\grave{\ }}},{{\phi }^{\grave{\ }}},{{r}^{\grave{\ }}} \right)\)
в точку наблюдения \(P\).
{[}\overrightarrow{R}=\overrightarrow{{{R}_{O}}}-\overrightarrow{R'}{]}
Где {[}\overrightarrow{R'}{]} - есть расстояние от точки взрыва \(Q\) до
центра взрыва \(O\). Раскладывая выражение для {[}\left\textbar{}
\overrightarrow{R} \right\textbar{}{]} {[}R=\left\textbar{}
\overrightarrow{R}
\right\textbar{}=\sqrt{{{R}_{O}}^{2}-2\overrightarrow{{{R}_{O}}}\overrightarrow{R'}+R{{'}^{2}}}{]}
в ряд Тейлора и пренебрегая вторыми и старшими степенями
{[}\{\{R\}\^{}\{'\}\}{]}, получаем
{[}R\approx {{R}_{O}}-\frac{\overrightarrow{{{R}_{O}}}\overrightarrow{R'}}{{{R}_{O}}}{]}
Далее, согласно {[} CITATION Там57 \l 1033 {]} выражая в подынтегральных
выражениях для скалярного и векторного потенциалов \(R\)через
{[}\{\{R\}}\{O\}\}-\frac{\overrightarrow{{{R}_{O}}}\overrightarrow{R'}}{{{R}_{O}}}{]},
разлагая в ряд Тейлора и ограничиваясь двумя членами разложения (что
допустимо для дальней, волновой, зоны) мы получаем выражение для
скалярного потенциала как сумму двух интегралов, первых из которых равен
полному заряду системы
{[}\int{\rho \left( \theta ',\phi ',r',t-\frac{{{R}_{O}}}{c} \right)}~dV'{]}
с множителем \({}^{1}/{}_{{{R}_{O}}}\) , вынесенным за знак интеграла, а
второй интеграл
{[}-\frac{\overrightarrow{{{R}_{O}}}}{{{R}_{O}}}\frac{\partial }{\partial {{R}_{O}}}\left\{
\frac{1}{{{R}_{O}}}\int{\overrightarrow{R'}\cdot \rho \left( \theta ',\phi ',r',t-\frac{{{R}_{O}}}{c} \right)}~dV'
\right\}=-\frac{\overrightarrow{{{R}_{O}}}}{{{R}_{O}}}\frac{\partial }{\partial {{R}_{O}}}\left\{
\frac{\overrightarrow{p}\left( t-\frac{{{R}_{O}}}{c} \right)}{{{R}_{O}}}
\right\}{]} равен дивергенции, взятой по точке наблюдения вектора Герца
системы {[}-di\{\{v\}\emph{\{a\}\}\left\{
\frac{\overrightarrow{p}\left( t-\frac{{{R}_{O}}}{c} \right)}{{{R}_{O}}}
\right\}{]}. В контексте решаемой задачи оба этих интеграла равны нулю.
Первый -- в виду изначальной электронейтральности системы. Второй -- в
виду того, что дипольный момент сферически симметричной системы зарядов
равен нулю. Что, впрочем, справедливо также и для квадрупольного,
октупольного моментов, а также мультипольных моментов более высоких
порядков. Это обстоятельство позволяет распространить полученные выводы
не только на дальнюю волновую зону, но также и на ближнюю зону.
Рассматривая разложение в ряд выражения для векторного потенциала мы в
контексте решения данной задачи отбрасываем равное нулю первое
слагаемое{[}\frac{1}{c{{R}_{O}}}\int{\overrightarrow{j}\left( \theta ',\phi ',r',t-\frac{{{R}_{O}}}{c} \right)}~dV'{]},
являющееся производной по времени вектора Герца системы
{[}\frac{1}{c{{R}_{O}}}\frac{\partial }{\partial t}p\left(
t-\frac{{{R}_{O}}}{c} \right){]}. Но вот второе слагаемое, отброшенное
Таммом при рассмотрении поля линейного осциллятора, в контексте
рассмотрения поля центрально-симметричного осциллятора следует
рассмотреть. Итак, интеграл
{[}\int{\frac{\overrightarrow{j}\left( \theta ',\phi ',r',t-\frac{R}{c} \right)}{R}}~dV'\approx -\int{\frac{\overrightarrow{R'}\overrightarrow{{{R}_{O}}}}{{{R}_{O}}}\frac{\partial }{\partial {{R}_{O}}}\frac{\overrightarrow{j}\left( \theta ',\phi ',r',t-\frac{{{R}_{O}}}{c} \right)}{{{R}_{O}}}\ dV'}{]}
Ввиду независимости \({{R}_{O}}\)от {[}\theta ',\phi ',r'{]} можно
записать в следующей форме
{[}-\frac{\overrightarrow{{{R}_{O}}}}{{{R}_{O}}}\frac{\partial }{\partial {{R}_{O}}}\left\{
\frac{1}{{{R}_{O}}}\int{\overrightarrow{R'}\cdot \overrightarrow{j}\left( \theta ',\phi ',r',t-\frac{{{R}_{O}}}{c} \right)}~dV'
\right\}{]} Интеграл
{[}\int{\overrightarrow{R'}\cdot \overrightarrow{j}\left( \theta ',\phi ',r',t-\frac{{{R}_{O}}}{c} \right)}~dV'=s\left(
t-\frac{{{R}_{O}}}{c} \right){]} Представляет собой величину, которую мы
назовём скалярным магнитным моментом центрально-симметричной системы
зарядов в момент {[}t-\frac{{{R}_{O}}}{c}{]} Таким образом векторный
потенциал равен
{[}\overrightarrow{{{A}_{H}}}=-\frac{1}{c}\frac{\overrightarrow{{{R}_{O}}}}{{{R}_{O}}}\frac{\partial }{\partial {{R}_{O}}}\left\{
\frac{s\left( t-\frac{{{R}_{O}}}{c} \right)}{{{R}_{O}}} \right\}{]} Это
выражение можно в окончательной форме записать как дивергенцию в точке
наблюдения
{[}\overrightarrow{{{A}_{H}}}=-\frac{1}{c}\frac{\overrightarrow{{{R}_{O}}}}{{{R}_{O}}}di\{\{v\}}\{a\}\}\left\{
\frac{s\left( t-\frac{{{R}_{O}}}{c} \right)}{{{R}_{O}}} \right\}{]}
Выражение стоящее под знаком дифференцирования
{[}\frac{s\left( t-\frac{{{R}_{O}}}{c} \right)}{{{R}_{O}}}{]}математически
аналогично такому понятию как вектор Герца. Поэтому для удобства
восприятия я буду в дальнейших выкладках именовать его скаляром Герца.
Более поздняя редакция. Соблазн переписать интеграл 19 в форме 110 ,
получив по аналогии с дипольным моментом так называемый скалярный момент
велик. Но математически эта операция не корректна. Потому что
произведение скалярного
произведения{[}\frac{\overrightarrow{R'}\overrightarrow{{{R}_{O}}}}{{{R}_{O}}}{]}
на вектор
{[}\frac{\partial }{\partial {{R}_{O}}}\frac{\overrightarrow{j}}{{{R}_{O}}}{]}
отличается от произведения вектора
{[}\frac{\overrightarrow{{{R}_{O}}}}{{{R}_{O}}}{]} на скалярное
произведение {[}\frac{\partial }{\partial {{R}_{O}}}\left\{
\frac{1}{{{R}_{O}}}\overrightarrow{R'}\cdot \overrightarrow{j}
\right\}{]} по меньшей мере направлением. Вот если бы
{[}\overrightarrow{j}{]} был бы не вектором, а скаляром, тогда такая
операция была бы математически допустима. А поскольку мы исследуем
сферически симметричную задачу, то мы можем ее переформулировать
следующим образом. Пусть точка наблюдения находится на оси \(z\) в
положительном направлении, тогда для вычисления векторного потенциала
нам достаточно знать его проекцию на \({{R}_{0}}\) Уравнение 15 примет
теперь вид {[}\{\{A\}\emph{\{B,\{\{R\}}\{0\}\}\}\}\left( r,t
\right)=\frac{\mu }{c}\int{\frac{{{j}_{{{R}_{0}}}}\left( \theta ',\phi ',r',t' \right)}{R}}~dV'{]}
А уравнение 19 становится
{[}\int{\frac{{{j}_{{{R}_{0}}}}\left( \theta ',\phi ',r',t-\frac{R}{c} \right)}{R}}~dV'\approx -\int{\frac{\overrightarrow{R'}\overrightarrow{{{R}_{O}}}}{{{R}_{O}}}\frac{\partial }{\partial {{R}_{O}}}\frac{{{j}_{{{R}_{0}}}}\left( \theta ',\phi ',r',t-\frac{{{R}_{O}}}{c} \right)}{{{R}_{O}}}\ dV'}{]}
где {[}\{\{j\}\emph{\{\{\{R\}}\{0\}\}\}\}{]}по сути это проекция тока
переноса на ось \(z\) {[}\{\{j\}\emph{\{\{\{R\}}\{0\}\}\}\}\left(
\theta ',\phi ',r',t-\frac{{{R}_{O}}}{c} \right)=j\left(
\theta ',\phi ',r',t-\frac{{{R}_{O}}}{c} \right)\cdot \cos \left(
\theta ' \right){]} Теперь уравнение 110 будет выглядеть так
{[}-\frac{\overrightarrow{{{R}_{O}}}}{{{R}_{O}}}\frac{\partial }{\partial {{R}_{O}}}\left\{
\frac{1}{{{R}_{O}}}\int{\overrightarrow{R'}\cdot {{j}_{{{R}_{0}}}}\left( \theta ',\phi ',r',t-\frac{{{R}_{O}}}{c} \right)}~dV'
\right\}{]}

Расчёт ЭМИ центрально-симметричного взрыва на основе скалярного
магнитного момента Выделим условно четыре периода развития взрыва:
период испарения и ионизации вещества -- образование плазмы, период
разогрева плазмы, период разлёта частиц из центра взрыва,
сопровождающийся ростом дебаевского радиуса системы, и период
схлопывания, сопровождающийся уменьшением дебаевского радиуса системы.
Для периода разлёта частиц рассмотрим приближение, в котором на
расстоянии r от центра в каждый момент периода разлёта частиц скорость и
ускорение частиц одного и того же типа одинаковы. Для расчёта скалярного
магнитного момента выразим плотность тока следующим образом
{[}\overrightarrow{j}\left( r',\varphi ',\theta ',t'
\right)=e\cdot n\left( r',t' \right)\cdot v\left( r',t' \right){]}
учитывая, что плотность и скорость частиц не зависит от полярных
координат. Скалярный магнитный момент {[}s\left( t'
\right)=\int{\overrightarrow{R'}\cdot \overrightarrow{j}\left( \theta ',\phi ',r',t' \right)}~dV'=\int{r'\cdot {{j}_{r}}\left( r',t' \right)}~4\pi r\{\{'\}\^{}\{2\}\}dr'{]}
Для вычисления связанной со скалярным магнитным моментом компоненты
электрического поля в первое слагаемое формулы 268 подставляем 112.
Получаем
{[}\overrightarrow{E}=-\frac{\mu }{c}\frac{\partial \overrightarrow{{{A}_{H}}}}{\partial t}=\frac{\mu }{{{c}^{2}}}\frac{\overrightarrow{{{R}_{O}}}}{{{R}_{O}}}\frac{\partial }{\partial t}\frac{\partial }{\partial {{R}_{O}}}\left\{
\frac{s\left( t-\frac{{{R}_{O}}}{c} \right)}{{{R}_{O}}} \right\}{]} Эта
компонента электрического поля возникает благодаря закону
электромагнитной индукции. Поскольку \({{R}_{O}}\)не зависит от времени
мы имеем право поменять порядок дифференцирования и отдельно рассмотреть
производную скалярного магнитного момента по времени
{[}\frac{\partial }{\partial t}s\left( t-\frac{{{R}_{O}}}{c} \right){]}.
{[}\overrightarrow{E}=\frac{\mu }{{{c}^{2}}}\frac{\overrightarrow{{{R}_{O}}}}{{{R}_{O}}}\frac{\partial }{\partial {{R}_{O}}}\left\{
\frac{1}{{{R}_{O}}}\frac{\partial }{\partial t}s\left(
t-\frac{{{R}_{O}}}{c} \right) \right\}{]} Поверхностный интеграл
нормальной компоненты электрического поля по поверхности сферы радиуса
будет иметь вид
{[}\int{{{E}_{n}}dS=4\pi {{R}_{O}}^{2}}\frac{\mu }{{{c}^{2}}}\frac{\partial }{\partial {{R}_{O}}}\left\{
\frac{1}{{{R}_{O}}}\frac{\partial }{\partial t}s\left(
t-\frac{{{R}_{O}}}{c} \right) \right\}{]} где выражение для скалярного
магнитного момента определено в следующем виде {[}

\begin{align}
  & s\left( t-\frac{{{R}_{O}}}{c} \right)=\int\limits_{0}^{\infty }{r'\cdot {{j}_{r}}\left( r',t-\frac{{{R}_{O}}}{c} \right)}\ 4\pi r{{'}^{2}}dr' \\ 
 & s\left( t-\frac{{{R}_{O}}}{c} \right)=e\int\limits_{0}^{\infty }{r'\cdot n\left( r',t-\frac{{{R}_{O}}}{c} \right)}\ \cdot {{v}_{r}}\left( r',t-\frac{{{R}_{O}}}{c} \right)4\pi r{{'}^{2}}dr' \\ 
\end{align}

{]} Рассмотрим движение заряженных частиц (электронов) заключенных в
объёме, ограниченном поверхностями сферы радиусов \(r'\) и \(r'+dr'\).
За время \(dt\) эти частицы переместятся таким образом, что они будут
заключены в объёме между поверхностями радиуса
\(r'+v'\left( r' \right)dt\) и \(r'+dr'+v'\left( r'+dr' \right)dt\).
Если во время разлёта частиц осуществляется подвод энергии, тогда
существенный вклад в {[}\frac{\partial }{\partial t}s\left(
t-\frac{{{R}_{O}}}{c} \right){]} вносит рост скорости частиц. Теперь
предположим, что подвод энергии на исходе и в какой-то момент происходит
стабилизация скорости частиц
\(v'\left( r' \right)=v'\left( r'+v'\left( r' \right)dt \right)\). В
этот момент положительный вклад в производную по времени скалярного
магнитного момента вносит процесс расширения облака заряженных частиц.
При расширении облака частиц концентрация {[}n\left(
r',t-\frac{{{R}_{O}}}{c} \right){]}убывает с квадратом {[}r'{]}. Но в
подынтегральное выражение входит {[}r\{\{'\}\^{}\{3\}\}{]}, чем
обусловлена положительная величина
{[}\frac{\partial }{\partial t}s\left( t-\frac{{{R}_{O}}}{c} \right){]}.
Положительная величина производной по времени скалярного магнитного
момента будет сохраняться некоторое время и на стадии замедления
(вследствие кулоновского взаимодействия) скорости разлетающихся
электронов вплоть до момента когда будет выполняться равенство
\(v'\left( r' \right)>const\cdot \frac{1}{r'}\). В момент когда для
основной массы разлетающихся электронов вследствие их замедления в
кулоновском поле станет выполняться условие
\(v'\left( r' \right)<const\cdot \frac{1}{r'}\)производная по времени
скалярного магнитного момента поменяет знак. Рассмотрим производную
{[}\frac{\partial }{\partial {{R}_{O}}}{]}в выражении для скорости
изменения векторного потенциала
{[}-~c\frac{{{\left( \partial {{A}_{H}} \right)}_{R}}}{\partial t}=\frac{\partial }{\partial {{R}_{O}}}\left\{
\frac{1}{{{R}_{O}}}\frac{\partial }{\partial t}s\left(
t-\frac{{{R}_{O}}}{c} \right) \right\}{]} Дифференцируем {[}

\begin{align}
  & -c\frac{{{\left( \partial {{A}_{H}} \right)}_{R}}}{\partial t}=\frac{\partial }{\partial t}s\left( t-\frac{{{R}_{O}}}{c} \right)\frac{\partial }{\partial {{R}_{O}}}\left\{ \frac{1}{{{R}_{O}}} \right\}+\frac{1}{{{R}_{O}}}\frac{\partial }{\partial {{R}_{O}}}\left\{ \frac{\partial }{\partial t}s\left( t-\frac{{{R}_{O}}}{c} \right) \right\} \\ 
 & =-\frac{1}{{{R}_{O}}^{2}}\frac{\partial }{\partial t}s\left( t-\frac{{{R}_{O}}}{c} \right)+\frac{1}{{{R}_{O}}}\frac{\partial }{\partial t}\left\{ \frac{\partial }{\partial {{R}_{O}}}s\left( t-\frac{{{R}_{O}}}{c} \right) \right\} \\ 
\end{align}

{]} Где {[}

\begin{align}
  & \frac{\partial }{\partial {{R}_{O}}}s\left( t-\frac{{{R}_{O}}}{c} \right)=\frac{\partial }{\partial {{R}_{O}}}\int\limits_{0}^{\infty }{r'\cdot {{j}_{r}}\left( r',t-\frac{{{R}_{O}}}{c} \right)}\ 4\pi r{{'}^{2}}dr' \\ 
 & =-\frac{4\pi }{c}\int\limits_{0}^{\infty }{r{{'}^{3}}\left( \frac{\partial }{\partial t'}{{j}_{r}}\left( r',t' \right),\left\{ t'=t-\frac{{{R}_{O}}}{c} \right\} \right)}\ \text{d}r' \\ 
\end{align}

{]} Интеграл 122 возникает благодаря учёту эффектов запаздывания
потенциала. Благодаря множителю
{[}-\{\}\textsuperscript{\{1\}/\{\}\emph{\{c\}{]}этот интеграл будет не
существенным в практических расчётах. В итоге получаем
{[}-c\frac{{{\left( \partial {{A}_{H}} \right)}_{R}}}{\partial t}=-\frac{1}{{{R}_{O}}^{2}}\frac{\partial }{\partial t}\int\limits}\{0\}}\{\infty \}\{r'\cdot {{j}_{r}}\left(
r',t'
\right)\}~4\pi r\{\{'\}\^{}\{2\}\}dr'-\frac{4\pi }{c}\frac{1}{{{R}_{O}}}\frac{\partial }{\partial t}\left\{
\int\limits\emph{\{0\}\textsuperscript{\{\infty \}\{r\{\{'\}}\{3\}\}\left(
\frac{\partial }{\partial t'}{{j}_{r}}\left( r',t' \right)
\right)\}~\text{d}r' \right\}{]} Пренебрегая эффектами запаздывания
потенциала, получаем
{[}-~c\frac{{{\left( \partial {{A}_{H}} \right)}_{R}}}{\partial t}=\frac{\partial }{\partial t}s\left(
t-\frac{{{R}_{O}}}{c} \right)\frac{\partial }{\partial {{R}_{O}}}\left\{
\frac{1}{{{R}_{O}}}
\right\}=-\frac{1}{{{R}_{O}}^{2}}\frac{\partial }{\partial t}s\left(
t-\frac{{{R}_{O}}}{c} \right){]} Таким образом, слагающая электрического
поля, создаваемая эффектом разлёта заряженных частиц
{[}\overrightarrow{E}=-\frac{\mu }{c}\frac{\partial \overrightarrow{{{A}_{H}}}}{\partial t}=-\frac{\mu }{{{c}^{2}}}\frac{\overrightarrow{{{R}_{O}}}}{{{R}_{O}}^{3}}\frac{\partial }{\partial t}s\left(
\{\{R\}}\{O\}\},t-\frac{{{R}_{O}}}{c} \right){]} Эта формула по своей
форме аналогична закону Кулона{[} CITATION Там57 \l 1049 {]}, \$3
{[}\overrightarrow{E}=\frac{q}{{{R}^{3}}}\overrightarrow{R}{]} 26126*
MERGEFORMAT (.) В котором роль дополнительного «заряда» играет величина
{[}\{\{q\}\emph{\{add\}\}=-\frac{\mu }{{{c}^{2}}}\frac{\partial }{\partial t}s\left(
\{\{R\}}\{O\}\},t-\frac{{{R}_{O}}}{c} \right){]} 27127* MERGEFORMAT (.)
Слово «заряд» применительно к формуле 127 я беру в кавычки поскольку
фактически здесь имеет место лишь возникновение электрического поля как
производной по времени векторного потенциала как следствие закона
электромагнитной индукции.

Сравнение с результатами эксперимента Пусть в конце фазы разогрева
первоначальная плазма имеет радиус \({{R}_{1}}\) а изначальная плотность
частиц равна
{[}\{\{n\}\_\{i\}\}=\frac{{{N}_{i}}}{\int\limits_{0}^{{{R}_{i}}}{4\pi {{r}^{2}}dr}}=\frac{3{{N}_{i}}}{4\pi {{R}_{1}}^{3}}{]}
Для распределения радиальной компоненты скорости частиц примем линейную
аппроксимацию в виде:

\begin{verbatim}
${{v}_{R}}=\alpha \frac{R}{{{R}_{1}}}{{v}_{T}}$     28128\* MERGEFORMAT (.)
\end{verbatim}

Где \({{v}_{T}}\) средняя тепловая скорость частиц, \$\alpha \$
коэффициент пропорциональности. На возможность принятия линейного
аппроксимации зависимости скорости частиц от координаты в качестве
первого приближения для проведения приближённых оценочных расчётов
указывают результаты расчётной работы {[}CITATION Сас16 \l 1049 {]}.
Приведенная здесь методика расчёта может быть усовершенствована с
применением итерационных компьютерных алгоритмов {[}CITATION АСД15
\l 1049 {]} Принимая в качестве статистики распределения скорости частиц
по окончании периода разогрева плазмы распределение Максвелла

{[}\frac{dN}{d{{v}_{0}}}=N\{\{\left[ \frac{m}{2\pi kT} \right]\}\textsuperscript{\{\{\}}\{3\}/\{\}\_\{2\}\}\}\{\{v\}\^{}\{2\}\}\exp \left(
-\frac{m{{v}^{2}}}{2kT} \right){]}

Получаем среднюю тепловую скорость частиц
\({{v}_{T}}=\sqrt{\frac{2{{k}_{B}}T}{m}}\) откуда
\({{v}_{R}}=\alpha \frac{R}{{{R}_{1}}}\sqrt{\frac{2{{k}_{B}}T}{m}}\)
29129* MERGEFORMAT (.) В описании своего опыта в работе «Является ли
заряд инвариантом скорости?» {[} CITATION ФФМ \l 1049 {]} приводит
следующие данные

``Если посчитать энергию, необходимую для разогрева, плавления и
испарения медной проволоки диаметром 0.2 мм и длиной 5 мм то она
составит около 8 Дж. При этом температура паров меди уже будет
составлять около 2800 К. Энергия же конденсатора ёмкостью 3000 мкФ,
заряженного до напряжения 300 В составляет 135 Дж. Следовательно,
энергия порядка 125 Дж уйдёт на разогрев паров меди и окружающего
воздуха, на их ионизацию и то световое и другие виды излучения, которое
сопутствует нагреву газа и плазмы.'' Итак, объём меди превращаемой в
плазму равен
\({{V}_{Cu}}=\pi {{r}^{2}}l=\pi \cdot {{0.1}^{2}}\cdot 0.5=\text{0}\text{.157e-3}\ {{}^{3}}\)
масса этой меди равна
{[}\{\{M\}\emph{\{Cu\}\}=8.92\cdot {{V}_{Cu}}=~\text{0}\text{.14e-2}~{]}
количество вещества
{[}\{\{\nu \}}\{Cu\}\}=\{\{\{M\}\emph{\{Cu\}\}\}/\{\text{63}\text{.546}=\};~\text{0}\text{.22e-4 }{]}
число атомов
\({{N}_{i}}={{\nu }_{Cu}}\cdot 6.022\cdot {{10}^{23}}=\ 1.3278\cdot {{10}^{19}}\)
Энергия разряда конденсатора
\({{E}_{C}}={3000\cdot {{10}^{-6}}\cdot {{300}^{2}}}/{2=135\ }\;\) Для
вычисления затрат энергии для нагревания до температуры плавления были
взяты табличные данные теплоёмкости меди при постоянном давлении в
\({}/{\left( \cdot \right)}\;\) из\\
Источники: {[}CITATION Зин89 \l 1033 {]} {[}CITATION Чир68 \l 1049 {]} T
= {[}300, 400, 500, 600, 700, 800, 900, 1000, 1100, 1200, 1300,
1357.6{]} Cp = {[}385.0, 397.7, 408.0, 416.9, 425.1, 432.9, 441.7,
451.4, 464.3, 480.8, 506.6, 525.2{]} При интерполяции кубическими
сплайнами и интегрировании в пределах от 300 до 1357.6 К было получено
\(4.66\cdot {{10}^{5}}\) Дж/кг , при умножении на массу меди получено
0.653 Дж. Затраты энергии на плавление 0.287 Дж Далее, когда через
медную проволоку проходит разряд большой мощности, то она не просто
испаряется при температуре кипения меди, а за счёт инерционности и
малого времени разряда давление внутри меди возрастает до очень больших
значений. При достаточной мощности разряда есть вероятность достижения
критического давления внутри меди. Поэтому дальнейшие расчёты будем
производить в предположении, что жидкая медь достигает критической точки
Расчётные данные о термодинамических свойствах меди в критической точке
согласно {[}CITATION EMA09 \l 1049 {]} температура
\({{T}_{1}}=7093\ \)плотность
{[}\{\{\rho \}}\{Cu\}\}\textsuperscript{\{crit\}=1.95~\{\}/\{\{\{\}}\{3\}\}\};{]}
давление = 450 атм Объём меди в критическом состоянии
\(V_{Cu}^{crit}={{{M}_{Cu}}}/{{{\rho }_{Cu}}^{crit}=\text{0}\text{.7185e-3}\ {{}^{3}}}\;\)
Первоначальный радиус как радиус шара меди в критическом состоянии
объёма \(V_{Cu}^{crit}\)
\({{R}_{1}}=0.001\cdot \sqrt[3]{\frac{3V_{Cu}^{crit}}{4\pi }}=\text{0}\text{.5556e-4}\ \)
Затраты энергии на ионизацию исходя из энергии ионизации меди (первый
электрон) 745 кДж/моль в предположении 100 процентной степени ионизации
составляют 16.4268 Дж Затраты на нагрев жидкой меди до состояния
критического пара оценить достаточно сложно поэтому приблизительно они
были оценены как затраты на нагрев при постоянном давлении (4.998 Дж) и
плюс затраты на испарение при постоянном давлении (6.716 Дж) + энергию
на сжатие пара до расчётного критического давления (450 атмосфер) А =
12.7477 Дж
{[}V\_\{Cu\}\^{}\{1~\}=\frac{{{\nu }_{Cu}}R{{T}_{1}}}{{{p}^{1\ }}}{]}
\(A=\int\limits_{V_{Cu}^{crit}}^{V_{Cu}^{1\ }}{\frac{{{\nu }_{Cu}}}{V}R{{T}_{1}}dV}\)
Остаток энергии, который идёт на разогрев уже ионизированной плазмы
\(\Delta E={{E}_{C}}-{{Q}_{pl}}-{{Q}_{isp}}-{{Q}_{nagrev}}-{{Q}_{ioniz}}-A\)
равен 93.825 Дж Промежуток времени, в течение которого происходит нагрев
ионизированной плазмы, равен
\(\Delta t=0.5\cdot {{10}^{-3}}\frac{\Delta E}{{{E}_{C}}}=0.3475\cdot {{10}^{-3}}\)
сек Максимальная разность температур, которая может быть сообщена
ионизированной плазме равна
{[}\Delta T=\frac{\Delta E}{{{c}_{p\_extrapolation}}\cdot {{\nu }_{Cu}}}=127~548~K{]}
При этом ее температура может составить
\({{T}_{e}}=\Delta T+{{T}_{1}}=134\ 641\ K\)

В температурных единицах 1 эВ соответствует
\(1.16\cdot {{10}^{4}}\ \)откуда максимальная температура плазмы в
электронвольтах равна \({{{T}_{e}}}/{1.16\cdot {{10}^{4}}=11.6\ }\;\)
Для сравнения расчётных данных приведу цитату из литературы {[}CITATION
Адамьян99 \l 1033 {]} «Наилучшее совпадение с данными эксперимента даст
расчёт при начальной плотности 1 кг/м3 и начальной температуре 5 эВ»
Принимаем в распределении радиальной компоненты скорости частиц
коэффициент пропорциональности равным \(\alpha =0.00025\) Подставляя
формулу линейной аппроксимации радиальной скорости частиц в 114 и 115
для скалярного магнитного момента 119 получено выражение

\begin{verbatim}
\[s\left( {{R}_{1}},q,m,T \right)=\frac{3{{N}_{i}}}{4\pi {{R}_{1}}^{3}}\cdot \int\limits_{0}^{{{R}_{1}}}{Rq\alpha \frac{R}{{{R}_{1}}}\sqrt{\frac{2{{k}_{B}}T}{m}}\cdot 4\pi {{R}^{2}}dR}\]  30130\* MERGEFORMAT (.)
\end{verbatim}

Теперь, в соответствии с формулой 127 дополнительный «заряд» возникающий
благодаря производной по времени скалярного магнитного момента, как
следствие закона электромагнитной индукции (производная по времени
векторного потенциала в точке наблюдения). В системе СИ
{[}\{\{q\}\emph{\{add\}\}=-\{\{\mu \}}\{0\}\}\{\{\varepsilon \}\emph{\{0\}\}\frac{\partial }{\partial t}s\left(
\{\{R\}}\{O\}\},t-\frac{{{R}_{O}}}{c} \right){]} Учитывая только лишь
рост температуры методом конечных разностей в соответствии с формулой
{[}\{\{q\}\emph{\{add\}\}=-\{\{\mu \}}\{0\}\}\{\{\varepsilon \}\_\{0\}\}\frac{\left( s\left( {{R}_{1}},-{{q}_{e}},{{m}_{e}},{{T}_{e}} \right)+s\left( {{R}_{1}},+{{q}_{e}},{{m}_{Cu}},{{T}_{e}} \right) \right)-\left( s\left( {{R}_{1}},-{{q}_{e}},{{m}_{e}},{{T}_{1}} \right)+s\left( {{R}_{1}},+{{q}_{e}},{{m}_{Cu}},{{T}_{1}} \right) \right)}{\Delta t}{]}
Для дополнительного «заряда» была получена величина
\(+8.81\cdot {{10}^{-15}}\) кулон Применяя вместо дифференцирования по
времени \(\frac{\Delta T}{\Delta t}\frac{\partial }{\partial T}\)
\(-{{\mu }_{0}}{{\varepsilon }_{0}}{{\left\{ \frac{\Delta T}{\Delta t}\frac{\partial }{\partial T}\left( s\left( {{R}_{1}},-{{q}_{e}},{{m}_{e}},T \right)+s\left( {{R}_{1}},+{{q}_{e}},{{m}_{Cu}},T \right) \right) \right\}}_{T={{T}_{1}}}}\)
\(+2.36\cdot {{10}^{-14}}\) кулон результат вычисления слагаемого,
связанного с попыткой учесть запаздывание потенциала, равен
\(-1.673\cdot {{10}^{-22}}\) кулон В работе {[}CITATION ФФМ \l 1049 {]}
измеренный в процессе разогрева плазмы дополнительный заряд,
рассчитывается следующим образом. «Рассмотрим решение задачи для случая
сферической конфигурации этих элементов, когда разогретая плазма
находится в центре сфер. Будем также считать, что размеры сгустка плазмы
значительно меньше размеров экранов, т.е. будем рассматривать случай
точечного заряда. Будем считать, что радиус клетки Фарадея равен
\({{r}_{1}}\) , а радиус внешнего экрана равен \({{r}_{2}}\) . В
рассматриваемой установке максимальный радиус нижней части экрана клетки
Фарадея составляет 0.11 м, а радиус внешнего экрана равен 0.15 м. Эти
размеры примем для сферических поверхностей, рассматриваемых в задаче,
для ориентировочного расчёта величины эквивалентного заряда. Амплитуда
отрицательной части импульса, составляет 30 мВ. Для этого случая
максимальная величина эквивалентного заряда взрыва, образовавшегося в
процессе разогрева плазмы, будет равен»
\(d{{q}_{\exp eriment}}=\frac{4\pi {{\varepsilon }_{0}}U}{\frac{1}{{{r}_{1}}^{2}}-\frac{1}{{{r}_{2}}^{2}}}=4\pi {{\varepsilon }_{0}}\frac{-30\cdot {{10}^{-3}}}{\frac{1}{{{0.11}^{2}}}-\frac{1}{{{0.15}^{2}}}}=-8.738\cdot {{10}^{-14}}\ \)
Этот дополнительный заряд имеет отрицательное значение.

Таким образом расчёт основанный на предположении о том, что импульс
электрического поля при центрально-симметричном взрыве плазмы связан,
согласно представлений классической электродинамики, дополненных
понятием скалярного магнитного момента, - с производной по времени
векторного потенциала в удалённой от центра взрыва точке измерения
электрического поля оказывается не в состоянии правильно оценить знак
импульса. Поэтому возникает предположение о том, что кроме
рассмотренного выше явления имеет место также компонента импульса
электрического поля, связанная с градиентом в точке измерения скалярного
потенциала, который в свою очередь связан с возникновением
действительного дополнительного заряда смещения вакуумной среды в форме
изменения во времени дивергенции векторного потенциала (скалярного
магнитного поля, по {[}CITATION ГНи \l 1049 {]} ), но уже не в точке
наблюдения, а в области расширяющейся плазмы взрыва. Расчёт ЭМИ
центрально-симметричного взрыва на основе концепции скалярного
магнитного поля Заряд смещения вакуумной среды Подставляя выражение для
скалярного момента в форме 119 в 112 запишем выражение для векторного
потенциала
{[}\overrightarrow{{{A}_{H}}}=-\frac{1}{c}\frac{\overrightarrow{{{R}_{O}}}}{{{R}_{O}}}\frac{\partial }{\partial {{R}_{O}}}\left\{
\frac{1}{{{R}_{O}}}s\left( \{\{R\}\emph{\{O\}\},t' \right)
\right\}=-\frac{1}{c}\frac{\overrightarrow{{{R}_{O}}}}{{{R}_{O}}}\frac{\partial }{\partial {{R}_{O}}}\left\{
\frac{1}{{{R}_{O}}}\int\limits}\{0\}\^{}\{\infty \}\{r'\cdot {{j}_{r}}\left(
r',t' \right)\}~4\pi r\{\{'\}\^{}\{2\}\}dr' \right\}{]} Для радиальной
компоненты векторного потенциала
{[}\{\{A\}\emph{\{\{\{H\}}\{R\}\}\}\}=-\frac{1}{c}\frac{\partial }{\partial {{R}_{O}}}\left\{
\frac{1}{{{R}_{O}}}s\left( \{\{R\}\emph{\{O\}\},t' \right)
\right\}=-\frac{1}{c}\frac{\partial }{\partial {{R}_{O}}}\left\{
\frac{1}{{{R}_{O}}}\int\limits}\{0\}\^{}\{\infty \}\{r'\cdot {{j}_{r}}\left(
r',t' \right)\}~4\pi r\{\{'\}\^{}\{2\}\}dr' \right\}{]} Вычисляем
дивергенцию векторного потенциала, воспользовавшись формулой векторного
анализа
\(div\ \overrightarrow{a}=\frac{1}{{{R}^{2}}}\frac{\partial }{\partial R}\left( {{R}^{2}}{{a}_{R}} \right)\)
{[}di\{\{v\}\emph{\{q\}\}~\overrightarrow{{{A}_{H}}}=\frac{1}{{{R}_{O}}^{2}}\frac{\partial }{\partial {{R}_{O}}}\left(
\{\{R\}}\{O\}\}\^{}\{2\}\{\{A\}\emph{\{\{\{H\}}\{R\}\}\}\}
\right)=-\frac{1}{c\ {{R}_{O}}^{2}}\frac{\partial }{\partial {{R}_{O}}}\left(
\{\{R\}\emph{\{O\}\}\^{}\{2\}\frac{\partial }{\partial {{R}_{O}}}\left\{
\frac{1}{{{R}_{O}}}\int\limits}\{0\}\^{}\{\infty \}\{r'\cdot {{j}_{r}}\left(
r',t' \right)\}~4\pi r\{\{'\}\^{}\{2\}\}dr' \right\} \right){]}
Расписываем пренебрегая при этом эффектами запаздывания потенциала
\(\frac{\partial }{\partial {{R}_{O}}}\left\{ \frac{1}{{{R}_{O}}}\int\limits_{0}^{\infty }{r'\cdot {{j}_{r}}\left( r',t' \right)}\ 4\pi r{{'}^{2}}dr' \right\}=\frac{\partial }{\partial {{R}_{O}}}\left\{ \frac{1}{{{R}_{O}}} \right\}\int\limits_{0}^{\infty }{r'\cdot {{j}_{r}}\left( r',t' \right)}\ 4\pi r{{'}^{2}}dr'\)
{[}\frac{\partial }{\partial {{R}_{O}}}\left\{
\frac{1}{{{R}_{O}}}\int\limits\emph{\{0\}\^{}\{\infty \}\{r'\cdot {{j}_{r}}\left(
r',t' \right)\}~4\pi r\{\{'\}\^{}\{2\}\}dr'
\right\}=-\frac{1}{{{R}_{O}}^{2}}\int\limits}\{0\}\^{}\{\infty \}\{r'\cdot {{j}_{r}}\left(
r',t' \right)\}~4\pi r\{\{'\}\^{}\{2\}\}dr'{]} Полученное при
дифференцировании по {[}\frac{\partial }{\partial {{R}_{O}}}{]}скаляра
Герца значение имеет чисто геометрическую природу.
\(di{{v}_{q}}\ \overrightarrow{{{A}_{H}}}=-\frac{1}{c\ {{R}_{0}}^{2}}\frac{\partial }{\partial {{R}_{0}}}\left( {{R}_{0}}^{2}\left[ -\frac{1}{{{R}_{O}}^{2}}\int\limits_{0}^{\infty }{r'\cdot {{j}_{r}}\left( r',t' \right)}\ 4\pi r{{'}^{2}}dr' \right] \right)\)
\(di{{v}_{q}}\ \overrightarrow{{{A}_{H}}}=\frac{1}{c\ {{R}_{0}}^{2}}\frac{\partial }{\partial {{R}_{0}}}\left( \int\limits_{0}^{{{R}_{O}}}{r'\cdot {{j}_{r}}\left( r',t' \right)}\ 4\pi r{{'}^{2}}dr' \right)\)
Итак, промежуточный результат
\(di{{v}_{q}}\ \overrightarrow{{{A}_{H}}}=\frac{1}{c\ {{R}_{0}}^{2}}\left( 4\pi {{R}_{O}}^{3}{{j}_{r}}\left( {{R}_{O}},t' \right) \right)=\frac{4\pi }{c}{{R}_{O}}{{j}_{r}}\left( {{R}_{O}},t' \right)\)
31131* MERGEFORMAT (.) Дополнительная компонента электрического поля
{[}\overrightarrow{{{E}_{add}}}=-\frac{\mu }{c}\frac{\partial \overrightarrow{{{A}_{H}}}}{\partial t}{]}
{[}di\{\{v\}\emph{\{q\}\}\overrightarrow{{{E}_{add}}}=-\frac{\mu }{c}\frac{\partial }{\partial t}di\{\{v\}}\{q\}\}\overrightarrow{{{A}_{H}}}=-\frac{4\pi \mu }{{{c}^{2}}}{{R}_{O}}\frac{\partial }{\partial t}{{j}_{r}}\left(
\{\{R\}\emph{\{O\}\},t' \right){]} Откуда (учитывая уравнение
\(di{{v}_{q}}\ \varepsilon \overrightarrow{{{E}_{add}}}=4\pi {{\rho }_{}}\))
связанная с этим плотность дополнительного заряда смещения вакуумной
среды
\({{\rho }_{}}=\frac{di{{v}_{q}}\ \varepsilon \overrightarrow{{{E}_{add}}}}{4\pi }=-\frac{\varepsilon \mu }{4\pi c}\frac{\partial \ }{\partial t}di{{v}_{q}}\overrightarrow{{{A}_{H}}}\)
\({{\rho }_{}}=-\frac{\varepsilon \mu }{{{c}^{2}}}\frac{\partial \ }{\partial t}\left( {{R}_{O}}{{j}_{r}}\left( {{R}_{O}},t' \right) \right)\)
32132* MERGEFORMAT (.) Производя вычисления дополнительного заряда,
возникшего в эксперименте, подставляя 129 и 114 в 132 и применяя вместо
дифференцирования по времени
\(\frac{\partial }{\partial t}\approx \frac{\Delta T}{\Delta t}\frac{\partial }{\partial T}\)
в системе СИ по формулам {[}\{\{\rho \}}\{\}\}\left(
R,\{\{R\}\emph{\{1\}\},q,m,T
\right)=-\{\{\mu \}}\{0\}\}\{\{\varepsilon \}\emph{\{0\}\}q\frac{3{{N}_{i}}}{4\pi {{R}_{1}}^{3}}\frac{\Delta T}{\Delta t}\frac{\partial }{\partial T}\left\{
R\alpha \frac{R}{{{R}_{1}}}\sqrt{\frac{2{{k}_{B}}T}{m}} \right\}{]}
{[}\{\{q\}}\{\}\}\left( \{\{R\}\emph{\{1\}\},q,m,T
\right)=\int\limits}\{0\}\^{}\{\{\{R\}\emph{\{1\}\}\}\{\{\{\rho \}}\{\}\}\left(
R,\{\{R\}\emph{\{1\}\},q,m,T \right)\}~4\pi {{R}^{2}}dR{]} Была получена
величина заряда смещения вакуумной среды равная
\(+2.36\cdot {{10}^{-14}}\ \) что численно равно дополнительному
«заряду» обусловленному явлению электромагнитной индукции, но опять
оцененный знак заряда смещения вакуумной среды не соответствует
результатам эксперимента. Скалярно-векторный потенциал Менде Оценка
дополнительного заряда по формулам скалярно-векторного потенциала Менде
\(d{{q}_{memde}}(m,q,T,{{N}_{i}})=q{{N}_{i}}\cosh \left( \frac{{{v}_{T}}\left( m,T \right)}{c} \right)\)
\(d{{q}_{memde}}({{m}_{e}},-{{q}_{e}},{{T}_{1}},{{N}_{i}})+d{{q}_{memde}}({{m}_{Cu}},{{q}_{e}},{{T}_{1}},{{N}_{i}})=-0.2544\cdot {{10}^{-5}}\ \)
приводит к результату, завышенному на 8-9 порядков. Что и следовало
ожидать исходя из показанной выше неприменимости концепции
скалярно-векторного потенциала к оценке ЭМИ центрально симметричного
взрыва. Противоречия с электродинамической теорией Известное из
классической электродинамики {[} CITATION Там57 \l 1049 {]} выражение
электрического поля через потенциалы
\$E=-\frac{\mu }{c}\frac{\partial }{\partial t}{{A}_{H}}-\nabla \varphi \$
оказывается неприменимым к рассматриваемой задаче, потому как для
рассматриваемого в данной работе момента начала расширения плазмы в
первое слагаемое даёт положительный импульс, а второе слагаемое равно
нулю, что в итоге не может дать полученный в ходе эксперимента импульс
отрицательного знака. Попытка вычислить электрическое поле импульса
исходя из градиента скалярного потенциала создаваемого зарядами смещения
вакуумной среды {[} CITATION ГНи \l 1049 {]}
\(E=-\nabla \varphi -\nabla {{\varphi }_{}}\) где \({{\varphi }_{}}\)
скалярный потенциал, создаваемый зарядами смещения вакуумной среды
{[}\{\{\varphi \}}\{\}\}=\frac{1}{\varepsilon }\int{\frac{{{\rho }_{}}\left( t-\frac{R}{v} \right)}{R}}~dV'{]}
Привела к тому же не совпадающему с экспериментом результату. Дело в
том, что и в электродинамике Николаева заряд смещения вакуумной среды
вычисляется с использованием соотношения
\(E=-\frac{\mu }{c}\frac{\partial }{\partial t}{{A}_{H}}\) 33133*
MERGEFORMAT (.) отражающем закон электромагнитной индукции Фарадея.
очевидно , что запись закона электромагнитной индукции в виде 133
отражает далеко не все физически возможные ситуации, в которых возникает
индукция электрического поля. В частности, анализ электромагнитного
импульса центрально симметричного взрыва показывает, что индукция
электрического поля в направлении по оси движущегося с ускорением заряда
отличается от рассчитываемой по формуле 133 индукции по меньшей мере
знаком, и не исключено что также и значением коэффициента , которую в
данном случае можно было бы назвать продольная магнитная проницаемость
вакуума. таким образом правильная формула для заряда электрического
смещения вакуума выглядела бы
\({{\rho }_{}}=-\frac{\varepsilon {{\mu }_{\parallel }}}{4\pi c}\frac{\partial \ }{\partial t}di{{v}_{q}}\overrightarrow{{{A}_{H}}}\)
34134* MERGEFORMAT (.) где \({{\mu }_{\parallel }}<0\) Но прежде чем
искать правильную формулу индукции электрического поля в
электродинамике, следует отметить, что что формулы, по которым
рассчитывается в классической электродинамике векторный потенциал не
соответствуют физической реальности, что можно увидеть хотя бы из
противоречий, на которых следует остановиться более подробно.

О противоречии, возникающем при выводе формулы для векторного потенциала
в калибровке Кулона Вывод уравнения для векторного потенциала в
калибровке Кулона даётся, например, в ЛЛ2 §43 На основе уравнения
Максвелла {[}rot~\overrightarrow{H}=\frac{4\pi }{c}\overrightarrow{j}{]}
35135* MERGEFORMAT (.) Выражая напряжённость магнитного поля через
векторный потенциал
{[}\overrightarrow{H}=rot~\overrightarrow{{{A}_{H}}}{]} применяя
известную из векторного анализа формулу получаем
{[}grad~div~\overrightarrow{{{A}_{H}}}-\Delta \overrightarrow{{{A}_{H}}}=\frac{4\pi }{c}\overrightarrow{j}{]}
36136* MERGEFORMAT (.) Выбирая для потенциалов калибровку кулона
{[}div~\overrightarrow{{{A}_{H}}}=0{]} получаем уравнение для векторного
потенциала постоянного магнитного поля
{[}\Delta \overrightarrow{{{A}_{H}}}=-\frac{4\pi }{c}\overrightarrow{j}{]}
37137* MERGEFORMAT (.) Решение этого уравнения, аналогично решению
уравнения Пуассона
{[}\overrightarrow{{{A}_{H}}}=\frac{1}{c}\int{\frac{\overrightarrow{j}}{R}}~dV{]}
38138* MERGEFORMAT (.) Но теперь возникает вопрос допустимо ли
использование уравнения 138 для вычисления векторного потенциала
одиночного заряда движущегося прямолинейно и равномерно с
нерелятивистской скоростью? Если принять на веру слова Ландау, который
пишет, что «векторный потенциал определён неоднозначно, поэтому на него
можно наложить любое условие», то, казалось бы, а почему бы и нет? Но
тем не менее для решения этого вопроса попробуем вычислить дивергенцию
векторного потенциала, опираясь на рассуждения {[} CITATION Там57
\l 1049 {]} §46 п.5
{[}di\{\{v\}\emph{\{a\}\}~\overrightarrow{{{A}_{H}}}=\frac{1}{c}di\{\{v\}}\{a\}\}~\int{\frac{\overrightarrow{j}}{R}}~dV{]}
39139* MERGEFORMAT (.) Порядок интегрирования по объёму токов и
дифференцирования по точке наблюдения может быть заменён на обратный
{[}di\{\{v\}\emph{\{a\}\}~\overrightarrow{{{A}_{H}}}=\frac{1}{c}~\int{di{{v}_{a}}\ \left( \frac{\overrightarrow{j}}{R} \right)}~dV{]}
40140* MERGEFORMAT (.) Применяя уравнение векторного анализа
{[}div\left( \varphi \overrightarrow{a} \right)=\varphi ~div\left(
\overrightarrow{a} \right)+a~grad\left( \varphi  \right){]}
{[}di\{\{v\}}\{a\}\}\left( \frac{\overrightarrow{j}}{R}
\right)=\frac{1}{R}di\{\{v\}\emph{\{a\}\}\left( \overrightarrow{j}
\right)+\overrightarrow{j}~gra\{\{d\}}\{a\}\}\left( \frac{1}{R}
\right)=\overrightarrow{j}~gra\{\{d\}\emph{\{a\}\}\left( \frac{1}{R}
\right){]} 41141* MERGEFORMAT (.) Поскольку значение
\(\overrightarrow{j}\) от координат точки наблюдения не зависит
{[}di\{\{v\}}\{a\}\}\left( \overrightarrow{j} \right)=0{]}. С другой
стороны, воспользовавшись формулой
\(gra{{d}_{q}}\left( \frac{1}{R} \right)=\frac{\overrightarrow{R}}{{{R}^{3}}}=-gra{{d}_{a}}\left( \frac{1}{R} \right)\)
{[}di\{\{v\}\emph{\{a\}\}\left( \frac{\overrightarrow{j}}{R}
\right)=-\overrightarrow{j}~gra\{\{d\}}\{q\}\}\left( \frac{1}{R}
\right)=-\frac{\overrightarrow{j}\ \overrightarrow{R}}{{{R}^{3}}}{]}
42142* MERGEFORMAT (.) Итак,
{[}di\{\{v\}\emph{\{a\}\}~\overrightarrow{{{A}_{H}}}=-\frac{1}{c}\int{\frac{\overrightarrow{j}\ \overrightarrow{R}}{{{R}^{3}}}}~dV{]}
43143* MERGEFORMAT (.) Выражая плотность тока через
\(\overrightarrow{j}=ne\overrightarrow{v}\), принимая для точечного
заряда в качестве концентрации зарядов дельта-функцию
\(n=\delta \left( \overrightarrow{R}-\overrightarrow{v}t \right)\) и
интегрируя
{[}di\{\{v\}}\{a\}\}~\overrightarrow{{{A}_{H}}}=-\frac{1}{c}\int{\frac{e\overrightarrow{v}\ \overrightarrow{R}}{{{R}^{3}}}}~\delta \left(
\overrightarrow{R}-\overrightarrow{v}t \right)dV{]} для точечного заряда
получим
{[}di\{\{v\}\emph{\{a\}\}~\overrightarrow{{{A}_{H}}}=-\frac{q}{c}\frac{\overrightarrow{v}\ \overrightarrow{R}}{{{R}^{3}}}{]}
44144* MERGEFORMAT (.) Таким образом при движении отрицательного заряда
по ходу его движения перед зарядом образуется область
\(di{{v}_{a}}\ \overrightarrow{{{A}_{H}}}>0\) а позади заряда образуется
область \(di{{v}_{a}}\ \overrightarrow{{{A}_{H}}}<0\). {[} CITATION ГНи
\l 1049 {]} приводит следующие формулы для векторного и скалярного
магнитных полей движущегося заряда
{[}\{\{H\}}\{\bot \}\}=\frac{1}{c}\overrightarrow{v}\times \overrightarrow{E}{]}
45145* MERGEFORMAT (.)
{[}\{\{H\}\emph{\{\parallel \}\}=\frac{1}{c}\overrightarrow{v}\cdot \overrightarrow{E}{]}
46146* MERGEFORMAT (.) Учитывая выражение электрического поля из закона
Кулона {[}\overrightarrow{E}=q\frac{\overrightarrow{R}}{{{R}^{3}}}{]}
сравнивая 146 с 144 находим выражение, определяющее скалярное магнитное
поле по Николаеву
{[}\{\{H\}}\{\parallel \}\}=-~di\{\{v\}\emph{\{a\}\}~\overrightarrow{{{A}_{H}}}{]}
К этим же результатам можно прийти, если в формуле 141 после замены
\(gra{{d}_{a}}\left( \frac{1}{R} \right)=-gra{{d}_{q}}\left( \frac{1}{R} \right)\)применить
соотношение векторного анализа
\(\overrightarrow{a}\ grad\ \varphi =div\left( \varphi \overrightarrow{a} \right)-\varphi \ div\left( \overrightarrow{a} \right)\)
{[}di\{\{v\}}\{a\}\}\left( \frac{\overrightarrow{j}}{R}
\right)=-~di\{\{v\}\emph{\{q\}\}\left( \frac{\overrightarrow{j}}{R}
\right)+\frac{1}{R}di\{\{v\}}\{q\}\}\left( \overrightarrow{j} \right){]}
47147* MERGEFORMAT (.) Учитывая уравнение непрерывности
\$di\{\{v\}\emph{\{q\}\}~\overrightarrow{j}=-\frac{\partial }{\partial t}\rho \$
получаем {[}di\{\{v\}}\{a\}\}\left( \frac{\overrightarrow{j}}{R}
\right)=-~di\{\{v\}\emph{\{q\}\}\left( \frac{\overrightarrow{j}}{R}
\right)-\frac{1}{R}\frac{\partial }{\partial t}\rho {]}
{[}di\{\{v\}}\{a\}\}~\overrightarrow{A}=-~\frac{1}{c}\int{\left( di{{v}_{q}}\left( \frac{\overrightarrow{j}}{R} \right)+\frac{1}{R}\frac{\partial }{\partial t}\rho  \right)dV}{]}
48148* MERGEFORMAT (.) Первое слагаемое интеграла можно, как это делает
Тамм, преобразовать по теореме Гаусса, так как пространственное
дифференцирование под знаком проводится по тем же координатам точек
истока, как и интегрирование, следовательно
{[}di\{\{v\}\_\{a\}\}~\overrightarrow{A}=-\frac{1}{c}~\oint{\frac{{{j}_{n}}}{R}}~dS-\frac{1}{c}\int{\left( \frac{1}{R}\frac{\partial }{\partial t}\rho  \right)dV}{]}
49149* MERGEFORMAT (.) Причем интегрирование первого слагаемого должно
быть распространено «по поверхности всех обтекаемый током проводников»
(цитата из Тамма).

Пусть заряд представляет собой движущийся вдоль своей оси цилиндр с
объёмной плотностью заряда \(n\). Пусть размер заряда вдоль направления
его движения \(l\). Выделим условно три области объёма: центральную
объёмом \(d{{V}_{centr}}\) и длиной \(l-2dl\), включающую в себя
практически весь объём заряженного цилиндра (исключая торцы) и две
торцевые объёмом \(d{{V}_{left}}\), \(d{{V}_{right}}\) и длиной \(2dl\)
каждая, включающие в себя лишь малую при-торцевую часть объёма заряда.
Пусть в некоторый момент времени \(t\) граничные поверхности заряда
проходят через центры торцевых объёмов. В этом случае поток вектора
\(\overrightarrow{j}\) через общую для всех трёх объёмов
\(d{{V}_{left}}\), \(d{{V}_{centr}}\) и \(d{{V}_{right}}\) поверхность
равен нулю. А производная плотности заряда по времени отлична от нуля
только в \(d{{V}_{centr}}\) и \(d{{V}_{right}}\). При движении заряда
слева направо
{[}\frac{\partial }{\partial t}{{\rho }_{left}}=-\frac{nev}{2dl}{]} и
{[}\frac{\partial }{\partial t}{{\rho }_{right}}=\frac{nev}{2dl}{]}.
Дивергенция векторного потенциала
{[}di\{\{v\}\emph{\{a\}\}~\overrightarrow{A}=-\frac{1}{c}\int{\left( \frac{1}{{{R}_{left}}}\frac{\partial }{\partial t}{{\rho }_{left}} \right)d{{V}_{left}}}-\frac{1}{c}\int{\left( \frac{1}{{{R}_{right}}}\frac{\partial }{\partial t}{{\rho }_{right}} \right)d{{V}_{right}}}{]}
{[}di\{\{v\}}\{a\}\}~\overrightarrow{A}=+\frac{1}{c}\int{\left( \frac{1}{{{R}_{left}}}\frac{nev}{2dl} \right)d{{V}_{left}}}-\frac{1}{c}\int{\left( \frac{1}{{{R}_{right}}}\frac{nev}{2dl} \right)d{{V}_{right}}}=\frac{S}{c}\left(
\frac{1}{{{R}_{left}}}-\frac{1}{{{R}_{right}}} \right)nev{]} Таким
образом показывается, что дивергенция векторного потенциала, создаваемая
также и не точечным движущимся зарядом, равно как и отрезком
незамкнутого тока не равна нулю. Таким образом мы видим, что
использование формулы 138 для расчёта векторного потенциала, полученной
в калибровке Кулона приводит к противоречию с этой самой калибровкой.
Какова природа этого противоречия? Может быть причина этого противоречия
заключается в игнорировании возникающих при движении одиночного заряда
токов смещения? Действительно, с учётом токов смещения любой незамкнутый
ток становится замкнутым. Николаев в работе {[} CITATION ГНи \l 1049 {]}
показывает, что в случае равномерного движения одиночного заряда
(магнитостатика) вместо уравнения 135 учитывающего только лишь ток
переноса для устранения возникающих при этом противоречий следует
пользоваться уравнением учитывающим полный ток как сумму тока переноса и
тока смещения. {[}rot~\overrightarrow{H}=\frac{4\pi }{c}\left(
\overrightarrow{j}+\{\{\overrightarrow{j}\}\emph{\{\}\} \right){]} Тогда
уравнение 136 перепишется в виде\\
{[}grad~div~\overrightarrow{{{A}_{}}}-\Delta \overrightarrow{{{A}_{}}}=\frac{4\pi }{c}\left(
\overrightarrow{j}+\{\{\overrightarrow{j}\}}\{\}\} \right){]} В котором
я изменяю индекс в обозначении векторного потенциала на букву П --
подразумевая под {[}\overrightarrow{{{A}_{}}}{]}полный векторный
потенциал, создаваемый полным током, в отличие от векторного потенциала
{[}\overrightarrow{{{A}_{H}}}{]}, который согласно классической
электродинамической теории создаётся только токами переноса. Выбирая для
полного векторного потенциала калибровку кулона
{[}div~\overrightarrow{{{A}_{}}}=0{]} получаем уравнение для полного
векторного потенциала постоянного магнитного поля
{[}\Delta \overrightarrow{{{A}_{}}}=-\frac{4\pi }{c}\left(
\overrightarrow{j}+\{\{\overrightarrow{j}\}\emph{\{\}\} \right){]}
решение которого имеет вид
{[}\overrightarrow{{{A}_{}}}=\frac{1}{c}\int{\frac{\left( \overrightarrow{j}+{{\overrightarrow{j}}_{}} \right)}{R}}~dV{]}
Теперь полученное выражение для полного векторного потенциала применим
для вычисления дивергенции полного векторного потенциала одиночного
заряда, движущегося прямолинейно и равномерно с нерелятивистской
скоростью. Вышеприведенные рассуждения, начиная от формулы 139 и до
формулы 147 остаются справедливыми также и для вычисления дивергенции
полного векторного потенциала через полный ток, таким образом получаем
{[}di\{\{v\}}\{a\}\}~\overrightarrow{{{A}_{}}}=\frac{1}{c}~\int{di{{v}_{a}}\ \left( \frac{{{\overrightarrow{j}}_{}}}{R} \right)}~dV{]}
где {[}di\{\{v\}\emph{\{a\}\}\left( \frac{{{\overrightarrow{j}}_{}}}{R}
\right)=-~di\{\{v\}}\{q\}\}\left( \frac{{{\overrightarrow{j}}_{}}}{R}
\right)+\frac{1}{R}di\{\{v\}\emph{\{q\}\}\left(
\{\{\overrightarrow{j}\}}\{\}\} \right){]} дивергенция полного тока, как
показывает Николаев в {[} CITATION ГНи \l 1049 {]} равна нулю. Хотя это
верно для случая магнитостатики. Для электродинамических явлений
сопровождаемых, например, электрические взрывы проволочек или ЭМИ
ядерного взрыва, это утверждение вполне может оказаться неверным, но для
рассматриваемого здесь равномерного и прямолинейного движения одиночного
заряда подходы магнитостатики справедливы. Поэтому второе слагаемое
рассматриваемого интеграла равно нулю. Первое слагаемое можно
преобразовать по теореме Гаусса, так как пространственное
дифференцирование под знаком проводится по тем же координатам точек
истока, как и интегрирование, следовательно
{[}di\{\{v\}\emph{\{a\}\}~\overrightarrow{{{A}_{}}}=-\frac{1}{c}~\oint{\frac{{{j}_{}}_{n}}{R}}~dS{]}
Поскольку интегрирование первого слагаемого должно быть распространено
«по поверхности всех обтекаемый током проводников», а в случае полного
тока поверхность интегрирования должна быть распространена на все
бесконечное пространство, занимаемое током смещения. Итак, в случае
дивергенции полного векторного потенциала действительно очевидно, что
для равномерного движения одиночного заряда
{[}di\{\{v\}}\{a\}\}~\overrightarrow{{{A}_{}}}=0{]} Таким образом в
магнитостатике калибровка Кулона действительно может быть применяема,
однако только лишь для полного векторного потенциала, но никак не для
классического векторного потенциала. Потому как дивергенция
классического векторного потенциала в магнитостатике не равна нулю. И
при этом первое уравнение Максвелла также и в магнитостатике должно
содержать в себе кроме тока проводимости также и ток смещения. Возникает
вопрос будет ли воспроизводиться рассмотренное здесь противоречия, если
использовать выражение 13 для классического векторного потенциала
полученное в калибровке Лоренца? Может ли оказаться так, что калибровку
Лоренца также правомерно применять лишь к полному векторному потенциалу,
но не правомерно к классическому? Но дело в том, что уравнение 13 в
классической электродинамике возникает не только как следствие
применения калибровки Лоренца, но и как решение волнового уравнения
Даламбера. Поэтому прежде чем прежде чем приступить к исследованию
дивергенции классического векторного потенциала в калибровке Лоренца,
необходимо обратить внимание на противоречие, допущенное в классической
электродинамике при выводе волнового уравнения Даламбера О противоречии,
возникающем в классической электродинамике при выводе волнового
уравнения Даламбера для векторного потенциала. Жонглирование
калибровками Предлагаю читателю проследить за выкладками И.Е.Тамма, с
помощью которых он: а) вводит понятие тока смещения в первое равнение
Максвелла §88 б) выводит на основании первого уравнения Максвелла
волновое уравнение Даламбера §94 Итак, токи смещения вводятся в первое
уравнение Максвелла следующим образом
{[}rot~\overrightarrow{H}=\frac{4\pi }{c}\left(
\overrightarrow{j}+\{\{\overrightarrow{j}\}\emph{\{\}\} \right){]}
50150* MERGEFORMAT (.) Затем берётся дивергенция от обоих частей
{[}div~\overrightarrow{j}+div~\{\{\overrightarrow{j}\}}\{\}\}=0{]}
51151* MERGEFORMAT (.) Затем на основании уравнения непрерывности для
токов проводимости и кулоновских зарядов
{[}div~\{\{\overrightarrow{j}\}\emph{\{\}\}=-~div~\overrightarrow{j}=\frac{\partial \rho }{\partial t}{]}
52152* MERGEFORMAT (.) Но дальше начинается интересная манипуляция,
заключающаяся в том, что Тамм для связи плотности кулоновских зарядов с
вектором электрической индукции использует уравнение электростатики
{[}div~\overrightarrow{D}=4\pi \rho {]} 53153* MERGEFORMAT (.) Но в чем
заключается манипуляция? В электростатике не исследуется движение
кулоновских зарядов, следовательно, электростатика ничего не знает ни о
токах проводимости, ни о токах смещения, ни о векторном потенциале, ни о
калибровке Лоренца. И поэтому связь электрической индукции
\(\overrightarrow{D}\) со скалярным потенциалом поля кулоновских зарядов
в электростатике выглядит следующим образом
{[}\overrightarrow{D}=\varepsilon \overrightarrow{E}=-~\varepsilon ~grad~\varphi {]}
54154* MERGEFORMAT (.) Поэтому, когда Тамм в §88 после преобразования
{[}div~\{\{\overrightarrow{j}\}}\{\}\}=\frac{1}{4\pi }\frac{\partial }{\partial t}div~\overrightarrow{D}=div\left(
\frac{1}{4\pi }\frac{\partial \overrightarrow{D}}{\partial t} \right){]}
55155* MERGEFORMAT (.) дает выражение для тока смещения в виде
{[}\{\{\overrightarrow{j}\}\emph{\{\}\}=\frac{1}{4\pi }\frac{\partial \overrightarrow{D}}{\partial t}{]}
56156* MERGEFORMAT (.) то при этом надо понимать, что полученное в итоге
первое уравнение Максвелла
{[}rot~\overrightarrow{H}=\frac{4\pi }{c}\overrightarrow{j}+\frac{\varepsilon }{c}\frac{\partial \overrightarrow{E}}{\partial t}{]}
57157* MERGEFORMAT (.) выведено без применения калибровки Лоренца
{[}div~\overrightarrow{{{A}_{H}}}=-\frac{\varepsilon }{c}\frac{\partial \varphi }{\partial t}{]}
таким образом, что его правильнее было бы записывать в потенциалах
следующим образом
{[}rot~\overrightarrow{H}=\frac{4\pi }{c}\overrightarrow{j}-\frac{\varepsilon }{c}\frac{\partial }{\partial t}grad~\varphi {]}
58158* MERGEFORMAT (.) Однако в § 94 Тамм использует полученное таким
образом (т.е. по сути в калибровке Кулона) первое уравнение Максвелла
для вывода волнового уравнения Даламбера уже в калибровке Лоренца. При
этом вектор электрического поля через потенциалы выражается уже
по-другому
{[}\overrightarrow{E}=-\frac{\mu }{c}\frac{\partial {{\overrightarrow{A}}_{H}}}{\partial t}-grad~\varphi {]}
59159* MERGEFORMAT (.) При этом Тамм поясняет: «скалярный потенциал
зависит лишь от распределения зарядов», а «векторный потенциал -- от
распределения токов проводимости». Однако «напряженность электрического
поля зависит не только от градиента скалярного потенциала, а также и от
производной по времени векторного потенциала. В этом обстоятельстве
проявляется закон электромагнитной индукции». Таким образом, способ
получения волнового уравнения Даламбера для векторного потенциала,
применённый Таммом в § 94 является некорректным результатом, полученным
с помощью жонглирования калибровками. В системе уравнений
электромагнитного поля Тамм приводит первое уравнение Максвелла в
формальной записи 157, в которую входит выражение для электрического
поля \(\overrightarrow{E}\), которое, однако, в § 94 имеет вид 159 и
таким образом включает в себя компоненту, обусловленную явлением
электромагнитной индукции, а в §88 используется другое выражение для
электрического поля 154, которое является результатом решения
дифференциального уравнения электростатики 153 и поэтому компоненту,
обусловленную явлением электромагнитной индукции не включает. Можно ли
при выводе волнового уравнения для векторного потенциала обойтись без
жонглирования калибровками? Вариант 1. Если в качестве выражения для
\(\overrightarrow{E}\) использовать выражение, пришедшее из
электростатики {[}\overrightarrow{E}=-grad~\varphi {]}, то в качестве
волнового уравнения получается уравнение Пуассона
\({{\nabla }^{2}}\overrightarrow{{{A}_{H}}}=-\frac{4\pi }{c}\overrightarrow{j}\),
однако таким уравнением не получится распространение электромагнитных
волн. Поэтому такой вариант непродуктивен Вариант 2. Если при выводе
уравнения для тока смещения использовать
{[}\overrightarrow{E}=-\frac{\mu }{c}\frac{\partial {{\overrightarrow{A}}_{H}}}{\partial t}-grad~\varphi {]},
то
\$di\{\{v\}}\{q\}\}~\overrightarrow{D}=di\{\{v\}\emph{\{q\}\}~\varepsilon \overrightarrow{E}=-\frac{\mu \varepsilon }{c}\frac{\partial }{\partial t}di\{\{v\}}\{q\}\}\overrightarrow{{{A}_{H}}}-di\{\{v\}\emph{\{q\}\}~grad~\varphi \$
Учитывая, что скалярный кулоновский потенциал зависит только лишь от
распределения кулоновских зарядов
\$\{\{\nabla \}\^{}\{2\}\}\varphi =-~4\pi \rho \$ для дивергенции
вектора электрической индукции, дополняя 4-е уравнение Максвелла,
получаем:
\$di\{\{v\}}\{q\}\}~\overrightarrow{D}=-\frac{\mu \varepsilon }{c}\frac{\partial }{\partial t}di\{\{v\}\emph{\{q\}\}\overrightarrow{{{A}_{H}}}+4\pi \rho \$
тогда выражение для дивергенции тока смещения
{[}div~\{\{\overrightarrow{j}\}}\{\}\}=\frac{\partial \rho }{\partial t}=\frac{1}{4\pi }\frac{\partial }{\partial t}\left(
di\{\{v\}\emph{\{q\}\}~\overrightarrow{D}+\frac{\mu \varepsilon }{c}\frac{\partial }{\partial t}di\{\{v\}}\{q\}\}\overrightarrow{{{A}_{H}}}
\right){]} откуда
{[}\{\{\overrightarrow{j}\}\emph{\{\}\}=\frac{1}{4\pi }\frac{\partial }{\partial t}\left(
\overrightarrow{D}+\frac{\mu \varepsilon }{c}\frac{\partial }{\partial t}\overrightarrow{{{A}_{H}}}
\right){]} тогда первое уравнение Максвелла принимает вид
{[}rot~\overrightarrow{H}=\frac{4\pi }{c}\overrightarrow{j}+\frac{\varepsilon }{c}\frac{\partial \overrightarrow{E}}{\partial t}+\frac{\mu \varepsilon }{{{c}^{2}}}\frac{{{\partial }^{2}}}{\partial {{t}^{2}}}\overrightarrow{{{A}_{H}}}{]}
подставляя выражение для электрического поля через потенциалы
{[}rot~\overrightarrow{H}=\frac{4\pi }{c}\overrightarrow{j}-\frac{\varepsilon }{c}\frac{\partial }{\partial t}grad~\varphi {]}
Далее
{[}rot~\overrightarrow{H}=rot~rot~\overrightarrow{{{A}_{H}}}=grad~div\overrightarrow{{{A}_{H}}}-\{\{\nabla \}\^{}\{2\}\}\overrightarrow{{{A}_{H}}}=\frac{4\pi }{c}\overrightarrow{j}-\frac{\varepsilon }{c}\frac{\partial }{\partial t}grad~\varphi {]}
подставляя соотношение между потенциалами соответствующее калибровке
Лоренца
{[}div~\overrightarrow{{{A}_{H}}}=-\frac{\varepsilon }{c}\frac{\partial \varphi }{\partial t}{]}
опять приходим к уравнению Пуассона
\({{\nabla }^{2}}\overrightarrow{{{A}_{H}}}=-\frac{4\pi }{c}\overrightarrow{j}\)
Таким образом показывается, что классическая электродинамика не содержит
корректного вывода волнового уравнения Даламбера. В классическом выводе
волнового уравнения Даламбера есть ещё один спорный
формально-математический момент, на который следует обратить внимание.
Выражение электрического поля через потенциалы в виде
{[}\overrightarrow{E}=-\frac{\mu }{c}\frac{\partial {{\overrightarrow{A}}_{H}}}{\partial t}-grad~\varphi {]}получено
исходя из второго уравнения Максвелла
\(rot\ \overrightarrow{E}=-\frac{1}{c}\frac{\partial \overrightarrow{B}}{\partial t}=rot\left( -\frac{\mu }{c}\frac{\partial \overrightarrow{{{A}_{H}}}}{\partial t} \right)\)
обобщающего опыты электромагнитной индукции Фарадея. Однако в процессе
вывода волнового уравнения Даламбера к этому выражению применяется
формально математическая операция взятия дивергенции, - с целью
приравнять полученную таким формальным способом математическую величину
с дивергенцией электрической индукции из 4 уравнения Максвелла. И здесь
возникает вопрос: если физический смысл уравнения электростатики
{[}div~\overrightarrow{D}=4\pi \rho {]}вполне очевиден, то какой
физический смысл имеет выражение {[}div\overrightarrow{E}=div\left(
-\frac{\mu }{c}\frac{\partial {{\overrightarrow{A}}_{H}}}{\partial t}-grad~\varphi  \right){]}?
Какую физическую ситуацию оно отражает? В каком физическом эксперименте
электромагнитная индукция Фарадея приводит к появлению дивергенции
электрического поля? Соответствуют ли уравнения для запаздывающих
потенциалов калибровке Лоренца? Поле векторного потенциала создаваемое
движущимся равномерно и прямолинейно зарядом, рассчитывается в
классической электродинамике в соответствии с
{[}\overrightarrow{{{A}_{H}}}=\frac{1}{c}\int{\frac{\overrightarrow{j}\left( x',y',z',t-\frac{R}{v} \right)}{R}}~dV'{]}
60160* MERGEFORMAT (.) Математически это выражение является решением
уравнения Даламбера для векторного потенциала. Представляет интерес
задаться вопросом, соответствует ли это выражение калибровке Лоренца?
Для этого проследим за выкладками {[} CITATION Там57 \l 1049 {]} §96, в
котором он ищет выражение для дивергенции от 160
{[}di\{\{v\}}\{a\}\}~\overrightarrow{{{A}_{H}}}=\frac{1}{c}di\{\{v\}\emph{\{a\}\}~\int{\frac{\overrightarrow{j}\left( x',y',z',t-\frac{R}{v} \right)}{R}}~dV'=\frac{1}{c}~\int{di{{v}_{a}}\frac{\overrightarrow{j}\left( x',y',z',t-\frac{R}{v} \right)}{R}}~dV'{]}
61161* MERGEFORMAT (.) Применяя уравнение векторного анализа
{[}div\left( \varphi \overrightarrow{a} \right)=\varphi ~div\left(
\overrightarrow{a} \right)+a~grad\left( \varphi  \right){]}
{[}di\{\{v\}}\{a\}\}\frac{\overrightarrow{j}\left( x',y',z',t-\frac{R}{v} \right)}{R}=\frac{1}{R}\cdot di\{\{v\}\emph{\{a\}\}~\overrightarrow{j}\left(
x',y',z',t-\frac{R}{v} \right)+\overrightarrow{j}\left(
x',y',z',t-\frac{R}{v} \right)\cdot \nabla \frac{1}{R}{]} Так как
аргумент вектора \(\overrightarrow{j}\) зависит от \(x,y,z\) только
через посредством эффективного времени \(t'=t-\frac{R}{v}\), то
{[}di\{\{v\}}\{a\}\}~\overrightarrow{j}\left( x',y',z',t-\frac{R}{v}
\right)=\frac{\partial {{j}_{x}}}{\partial t'}\frac{\partial t'}{\partial x}+\frac{\partial {{j}_{y}}}{\partial t'}\frac{\partial t'}{\partial y}+\frac{\partial {{j}_{z}}}{\partial t'}\frac{\partial t'}{\partial z}=\frac{\partial \overrightarrow{j}}{\partial t'}\cdot \nabla t'=\frac{\partial \overrightarrow{j}}{\partial t'}\cdot \nabla \left(
t-\frac{R}{v}
\right)=-\frac{1}{v}\frac{\partial \overrightarrow{j}}{\partial t'}\cdot \nabla R{]}
Окончательно
{[}di\{\{v\}\emph{\{a\}\}\frac{\overrightarrow{j}}{R}=\overrightarrow{j}\cdot \nabla \frac{1}{R}-\frac{1}{vR}\frac{\partial \overrightarrow{j}}{\partial t'}\cdot \nabla R{]}
С другой стороны, очевидно, что
\(div{{'}_{q}}\ \frac{\overrightarrow{j}}{R}={{\left( div{{'}_{q}}\ \frac{\overrightarrow{j}}{R} \right)}_{t'=const}}+\frac{1}{R}\frac{\partial \overrightarrow{j}}{\partial t'}\cdot \nabla 't'\)
где \(div{{'}_{q}}\) и \(\nabla '\) в отличие от \(di{{v}_{a}}\) и
\$\nabla \$ означает дифференцирование по координатам истока . После
преобразований
{[}div\{\{'\}}\{q\}\}\frac{\overrightarrow{j}}{R}=\overrightarrow{j}\cdot \nabla '\frac{1}{R}+\frac{1}{R}{{\left( div{{'}_{q}}\ \overrightarrow{j} \right)}_{t'=const}}-\frac{1}{vR}\frac{\partial \overrightarrow{j}}{\partial t'}\cdot \nabla 'R{]}
Приняв во внимание, что уравнение
\(gra{{d}_{a}}\ f\left( R \right)=-\ gra{{d}_{q}}\ f\left( R \right)\) в
теперешних обозначениях принимает вид {[}\nabla f\left( R
\right)=-~\nabla 'f\left( R \right){]} и сравнивая выражения для
{[}di\{\{v\}\emph{\{a\}\}\frac{\overrightarrow{j}}{R}{]}и
{[}div\{\{'\}}\{q\}\}\frac{\overrightarrow{j}}{R}{]}
{[}di\{\{v\}\emph{\{a\}\}\frac{\overrightarrow{j}}{R}=-~div\{\{'\}}\{q\}\}\frac{\overrightarrow{j}}{R}+\frac{1}{R}{{\left( div{{'}_{q}}\ \overrightarrow{j} \right)}_{t'=const}}{]}
Стало быть из 161 получено Таммом
{[}di\{\{v\}\emph{\{a\}\}~\overrightarrow{{{A}_{H}}}=-\frac{1}{c}~\int{div{{'}_{q}}\frac{\overrightarrow{j}\left( x',y',z',t-\frac{R}{v} \right)}{R}}~dV'+\frac{1}{c}\int{\frac{dV'}{R}}{{\left( div{{'}_{q}}\ \overrightarrow{j} \right)}_{t'=const}}{]}
Первый из этих интегралов может быть преобразован с помощью теоремы
Гаусса в интеграл по поверхности \(S\), охватывающей объём \(V'\)
{[}\int{div{{'}_{q}}\frac{\overrightarrow{j}\left( x',y',z',t-\frac{R}{v} \right)}{R}}~dV'=\oint{\frac{{{j}_{n}}\left( x',y',z',t-\frac{R}{v} \right)}{R}}~dS{]}
Далее Тамм пишет: «если распространить интегрирование на все бесконечное
пространство, то этот интеграл обратится в нуль, если только все
электрические токи сосредоточены в конечной области пространства».
Однако аналогичный интеграл встречается у Тамма в §46 (в данной работе
это первое слагаемое выражения 149 ), в котором рассматривая
пондемоторное взаимодействие постоянных токов, Тамм не распространял
интегрирование на все бесконечное пространство, а писал, что
«интегрирование должно быть распространено по поверхности всех
обтекаемых током проводников». И это верно в рамках представлений
классической электродинамики, поскольку классическая электродинамика при
вычислении векторного потенциала (как в калибровке Кулона, так и в
калибровке Лоренца) производит интегрирование только лишь по объёму
токов проводимости, занимающих конечный объём, но не производит
интегрирование по объёму токов смещения, занимающих все бесконечное
пространство. Следовательно, попытка Тамма при вычислении интеграла по
объёму токов проводимости «распространить интегрирование на все
бесконечное пространство» является, на мой взгляд, математически
некорректной. Самое интересное, что, перевернув страницу, в продолжении
доказательства Тамма можно встретить ещё одно утверждение, в
математической корректности которого есть определённые сомнения. Тамм
пишет: «обратимся теперь к правой части соотношения (96.3).
{[}\varphi \left( x,y,z,t
\right)=\frac{1}{\varepsilon }\int{\frac{\rho \left( x',y',z',t-\frac{R}{v} \right)}{R}}~dV'{]}
В виду уравнения (95.5) \(t'=t-\frac{R}{v}\) для произвольной функции
\(f\left( t' \right)\) имеем
\(\frac{\partial f\left( t' \right)}{\partial t}=\frac{\partial f\left( t' \right)}{\partial t'}\frac{\partial t'}{\partial t}=\frac{\partial f\left( t' \right)}{\partial t'}\)
Поэтому дифференцируя уравнение (96.3) по времени под знаком интеграла»,
Тамм получает:
{[}-\frac{\varepsilon \mu }{c}\frac{\partial \varphi \left( x,y,z,t \right)}{\partial t}=-\frac{\mu }{c}\int{\frac{dV'}{R}\frac{\partial \rho \left( x',y',z',t-\frac{R}{v} \right)}{\partial t'}}~{]}
Здесь следует обратить внимание, что приравнивание
\(\frac{\partial t'}{\partial t}\)единице не корректно. Чтобы это
показать, найдём производную запаздывающего момента \(t'=t-\frac{R}{c}\)
по времени {[}\frac{\partial t'}{\partial t}{]} Дифференцируя {[}R\left(
\overrightarrow{x},~\overrightarrow{x'}\left( t' \right) \right)=c\left(
t-t' \right){]} по получаем
\(\frac{\partial R}{\partial t}=\frac{\partial R}{\partial t'}\frac{\partial t'}{\partial t}=c\left( 1-\frac{\partial t'}{\partial t} \right)\).
Значение \(\frac{\partial R}{\partial t'}\)находим так
{[}\frac{\partial R\left( t' \right)}{\partial t'}=\frac{\partial }{\partial t'}{{\left( \overrightarrow{R}\cdot \overrightarrow{R} \right)}^{{}^{1}/{}_{2}}}=\frac{1}{2}{{\left( \overrightarrow{R}\cdot \overrightarrow{R} \right)}^{-{}^{1}/{}_{2}}}\frac{\partial }{\partial t'}\left(
\overrightarrow{R}\cdot \overrightarrow{R} \right)=\frac{1}{2R}\left(
\overrightarrow{R}\cdot \frac{\partial \overrightarrow{R}}{\partial t'}+\frac{\partial \overrightarrow{R}}{\partial t'}\cdot \overrightarrow{R}
\right)=\frac{\overrightarrow{R}\left( t' \right)}{R\left( t' \right)}\cdot \frac{\partial \overrightarrow{R}\left( t' \right)}{\partial t'}=-\frac{\overrightarrow{R}\cdot \overrightarrow{v}}{R}{]}
Далее учитывая что
\(\frac{\partial \overrightarrow{R}}{\partial t'}=-\overrightarrow{v}\left( t' \right)\)
имеем
{[}\frac{\partial R}{\partial t'}=-\frac{\overrightarrow{R}}{R}\cdot \overrightarrow{v}{]}
теперь
\(\frac{\partial R}{\partial t'}\frac{\partial t'}{\partial t}=c\left( 1-\frac{\partial t'}{\partial t} \right)\)преобразуется
к
\(-\frac{\overrightarrow{R}\cdot \overrightarrow{v}}{cR}\frac{\partial t'}{\partial t}=1-\frac{\partial t'}{\partial t}\)
решая которое получаем
{[}\frac{\partial t'}{\partial t}=\frac{1}{1-\frac{\overrightarrow{v}\cdot \overrightarrow{R}}{Rc}}{]}
Отсюда правильно производную скалярного потенциала по времени записать
так
{[}-\frac{\varepsilon \mu }{c}\frac{\partial \varphi \left( x,y,z,t \right)}{\partial t}=-\frac{\mu }{c}\int{\frac{dV'}{R-\frac{\overrightarrow{v}\cdot \overrightarrow{R}}{c}}\frac{\partial \rho \left( x',y',z',t-\frac{R}{v} \right)}{\partial t'}}{]}
Применяя соотношение непрерывности в виде {[}\{\{\left(
div\{\{'\}}\{q\}\}~\overrightarrow{j}
\right)\}\emph{\{t'=const\}\}=-\frac{\partial \rho \left( x',y',z',t' \right)}{\partial t'}{]},
убеждаемся что
{[}\frac{1}{c}~\int{div{{'}_{q}}\frac{\overrightarrow{j}\left( x',y',z',t-\frac{R}{v} \right)}{R}}~dV'+di\{\{v\}}\{a\}\}~\overrightarrow{{{A}_{H}}}=-\frac{1}{c}\int{\frac{dV'}{R}}\frac{\partial \rho \left( x',y',z',t' \right)}{\partial t'}{]}
Допуская применимость калибровки Лоренца
{[}div~\overrightarrow{{{A}_{H}}}=-\frac{\varepsilon }{c}\frac{\partial \varphi }{\partial t}{]}.
запишем следующее соотношение
{[}\frac{1}{c}~\int{div{{'}_{q}}\frac{\overrightarrow{j}\left( x',y',z',t-\frac{R}{v} \right)}{R}}~dV'-\frac{1}{c}\int{\frac{dV'}{R-\frac{\overrightarrow{v}\cdot \overrightarrow{R}}{c}}\frac{\partial \rho \left( x',y',z',t-\frac{R}{v} \right)}{\partial t'}}=-\frac{1}{c}\int{\frac{dV'}{R}}\frac{\partial \rho \left( x',y',z',t' \right)}{\partial t'}{]}
Или
{[}\frac{1}{c}~\oint{\frac{{{j}_{n}}\left( x',y',z',t' \right)}{R}}~dS'-\frac{1}{c}\int{\frac{dV'}{R-\frac{\overrightarrow{v}\cdot \overrightarrow{R}}{c}}\frac{\partial \rho \left( x',y',z',t' \right)}{\partial t'}}=-\frac{1}{c}\int{\frac{dV'}{R}}\frac{\partial \rho \left( x',y',z',t' \right)}{\partial t'}{]}
62162* MERGEFORMAT (.) Из вышеизложенного видно, что доказательство
соответствия решения уравнения Даламбера калибровке Лоренца, приведенное
Таммом содержит две математические некорректности. Из вышеизложенного
становится очевидным, что приём распространения интегрирования по токам
проводимости на все бесконечное пространство не верен. Потому что при
обнулении интеграла по {[}dS'{]} в последнем соотношении мы приходим к
очевидному неравенству, следствием которого становится неприменимость
соотношения Лоренца для потенциалов. С другой стороны, справедливость
полученного в данной работе соотношения 162 для произвольных токовых
систем также не является очевидной.

Соответствуют ли потенциалы Лиенара-Вихерта калибровке Лоренца? Чтобы
проверить сомнение о применимости калибровки Лоренца к классическому
векторному потенциалу в электродинамике, удобно воспользоваться таким
математическим инструментом, как потенциалы Лиенара-Вихерта, которые
определены как потенциалы поля, создаваемого движущимся одиночным
(точечным) зарядом. Потенциалы Лиенара-Вихерта, как известно {[}CITATION
ЕМЛ73 \l 1049 {]} получаются в результате интегрирования запаздывающих
потенциалов 12 и 13, которые в свою очередь являются решением волновых
уравнений Даламбера, которые получены из уравнений Максвелла с
применением калибровки Лоренца, хотя и не вполне корректным путём в том
смысле, что там, где авторам было удобно применялась калибровка Лоренца,
а там, где не удобно, - неявно применялась калибровка Кулона из
электростатики. Итак, записываем потенциалы ЛВ
{[}\varphi =\frac{1}{\varepsilon }\frac{e}{\left( R-\frac{\overrightarrow{v}\left( x',y',z',t' \right)\cdot \overrightarrow{R}}{c} \right)}{]}
{[}\overrightarrow{{{A}_{H}}}=\frac{1}{c}\frac{e\ \overrightarrow{v}\left( x',y',z',t' \right)}{\ R-\frac{\overrightarrow{v}\left( x',y',z',t' \right)\cdot \overrightarrow{R}}{c}}{]}
Производная скалярного потенциала по времени равна
{[}\frac{\partial \varphi }{\partial t}=\frac{\partial \varphi }{\partial t'}\frac{\partial t'}{\partial t}{]}
Найдём производную запаздывающего момента \(t'=t-\frac{R}{c}\) по
времени {[}\frac{\partial t'}{\partial t}{]} Дифференцируя {[}R\left(
\overrightarrow{x},~\overrightarrow{x'}\left( t' \right) \right)=c\left(
t-t' \right){]} по \(t\) получаем
\(\frac{\partial R}{\partial t}=\frac{\partial R}{\partial t'}\frac{\partial t'}{\partial t}=c\left( 1-\frac{\partial t'}{\partial t} \right)\).
Значение \(\frac{\partial R}{\partial t'}\)находим так
{[}\frac{\partial R\left( t' \right)}{\partial t'}=\frac{\partial }{\partial t'}{{\left( \overrightarrow{R}\cdot \overrightarrow{R} \right)}^{{}^{1}/{}_{2}}}=\frac{1}{2}{{\left( \overrightarrow{R}\cdot \overrightarrow{R} \right)}^{-{}^{1}/{}_{2}}}\frac{\partial }{\partial t'}\left(
\overrightarrow{R}\cdot \overrightarrow{R} \right)=\frac{1}{2R}\left(
\overrightarrow{R}\cdot \frac{\partial \overrightarrow{R}}{\partial t'}+\frac{\partial \overrightarrow{R}}{\partial t'}\cdot \overrightarrow{R}
\right)=\frac{\overrightarrow{R}\left( t' \right)}{R\left( t' \right)}\cdot \frac{\partial \overrightarrow{R}\left( t' \right)}{\partial t'}=-\frac{\overrightarrow{R}\cdot \overrightarrow{v}}{R}{]}
Далее учитывая что
\(\frac{\partial \overrightarrow{R}}{\partial t'}=-\overrightarrow{v}\left( t' \right)\)
имеем
{[}\frac{\partial R}{\partial t'}=-\frac{\overrightarrow{R}}{R}\cdot \overrightarrow{v}{]}
теперь
\(\frac{\partial R}{\partial t'}\frac{\partial t'}{\partial t}=c\left( 1-\frac{\partial t'}{\partial t} \right)\)преобразуется
к
\(-\frac{\overrightarrow{R}\cdot \overrightarrow{v}}{cR}\frac{\partial t'}{\partial t}=1-\frac{\partial t'}{\partial t}\)
решая которое получаем
{[}\frac{\partial t'}{\partial t}=\frac{1}{1-\frac{\overrightarrow{v}\cdot \overrightarrow{R}}{Rc}}{]}
Для производной скалярного потенциала \$\varphi \$ по \(t'\) вводим
обозначение \(K=R-\frac{\overrightarrow{v}\cdot \overrightarrow{R}}{c}\)
благодаря которому запись скалярного потенциала упрощается
{[}\varphi =\frac{1}{\varepsilon }\frac{e}{K}{]} кроме того
{[}\frac{\partial t'}{\partial t}=\frac{R}{K}{]}
\(\frac{\partial \varphi }{\partial t'}=\frac{\partial \varphi }{\partial K}\frac{\partial K}{\partial t'}=-\frac{e}{\varepsilon {{K}^{2}}}\left\{ \frac{\partial R}{\partial t'}-\frac{1}{c}\frac{\partial }{\partial t'}\left( \overrightarrow{v}\cdot \overrightarrow{R} \right) \right\}\)
\(\frac{\partial }{\partial t'}\left( \overrightarrow{v}\cdot \overrightarrow{R} \right)=\overrightarrow{v}\frac{\partial \overrightarrow{R}}{\partial t'}+\overrightarrow{R}\frac{\partial \overrightarrow{v}}{\partial t'}=-\overrightarrow{v}\cdot \overrightarrow{v}+\overrightarrow{R}\cdot \overrightarrow{a}\)
\(\frac{\partial \varphi }{\partial t'}=-\frac{e}{\varepsilon {{K}^{2}}}\left\{ -\frac{\overrightarrow{R}}{R}\cdot \overrightarrow{v}+\frac{1}{c}\left( \overrightarrow{v}\cdot \overrightarrow{v}-\overrightarrow{R}\cdot \overrightarrow{a} \right) \right\}\)
Итак
{[}\frac{\partial \varphi }{\partial t}=-\frac{e}{\varepsilon }\frac{R\frac{\vec{v}\cdot \vec{v}-\overrightarrow{a}\cdot \overrightarrow{R}}{c}-~\overrightarrow{v}\cdot \overrightarrow{R}}{{{K}^{3}}}{]}
Дивергенция векторного потенциала
{[}\nabla \overrightarrow{{{A}_{H}}}=\frac{e}{c}\nabla \frac{\overrightarrow{v}}{K}=\frac{e}{c}\left\{
\frac{1}{K}\nabla \overrightarrow{v}+\overrightarrow{v}\cdot \nabla \frac{1}{K}
\right\}{]} Дифференцируя скорость точечного заряда по координатам точки
наблюдения {[}\{\{\nabla \}\emph{\{a\}\}~\overrightarrow{v}\left(
x',y',z',t-\frac{R}{v}
\right)=\frac{\partial {{v}_{x}}}{\partial t'}\frac{\partial t'}{\partial x}+\frac{\partial {{v}_{y}}}{\partial t'}\frac{\partial t'}{\partial y}+\frac{\partial {{v}_{z}}}{\partial t'}\frac{\partial t'}{\partial z}=\frac{\partial \overrightarrow{v}}{\partial t'}\cdot \nabla t'{]}
Градиент запаздывающего момента времени {[}\nabla t'{]} {[}R\left(
\overrightarrow{x},~\overrightarrow{x'}\left( t' \right) \right)=c\left(
t-t' \right){]}
{[}t'=t-\frac{R\left( \overrightarrow{x},\overrightarrow{x'}\left( t' \right) \right)}{c}{]}
Дифференцируем по координатам {[}\nabla t'=-\frac{1}{c}\nabla R\left(
\overrightarrow{x},\overrightarrow{x'}\left( t' \right) \right){]}
{[}\nabla R\left( \overrightarrow{x},\overrightarrow{x'}\left( t'
\right) \right)=\nabla R\{\{\left(
\overrightarrow{x},\overrightarrow{x'}\left( t' \right)
\right)\}}\{t'=const\}\}+\nabla R\{\{\left(
\overrightarrow{x},\overrightarrow{x'}\left( t' \right)
\right)\}\_\{\overrightarrow{x}=const\}\}=\frac{\overrightarrow{R}}{R}+\frac{\partial R}{\partial t'}\nabla t'{]}
{[}\nabla t'=-\frac{1}{c}\left(
\frac{\overrightarrow{R}}{R}+\frac{\partial R}{\partial t'}\nabla t'
\right){]} Далее {[}\nabla t'=-\frac{1}{c}\left(
\frac{\overrightarrow{R}}{R}+\frac{\overrightarrow{R}\left( t' \right)}{R\left( t' \right)}\cdot \frac{\partial \overrightarrow{R}\left( t' \right)}{\partial t'}\nabla t'
\right){]} Учитывая что
{[}\frac{\partial \overrightarrow{R}\left( t' \right)}{\partial t'}=-\overrightarrow{v}\left(
t' \right){]} {[}\nabla t'=-\frac{1}{c}\left(
\frac{\overrightarrow{R}}{R}-\frac{\overrightarrow{R}\left( t' \right)}{R\left( t' \right)}\cdot \overrightarrow{v}\left(
t' \right)\nabla t' \right){]}
{[}\nabla t'=-\frac{1}{c}\frac{\overrightarrow{R}}{R}+\frac{1}{c}\frac{\overrightarrow{R}\left( t' \right)}{R\left( t' \right)}\cdot \overrightarrow{v}\left(
t' \right)\nabla t'{]} {[}\nabla t'\left(
1-\frac{1}{c}\frac{\overrightarrow{R}\left( t' \right)\cdot \overrightarrow{v}\left( t' \right)}{R\left( t' \right)}
\right)=-\frac{1}{c}\frac{\overrightarrow{R}}{R}{]}
{[}\nabla t'=-\frac{\ \overrightarrow{R}}{c\ \left( R-\frac{\overrightarrow{v}\cdot \overrightarrow{R}}{c} \right)}=-\frac{\ \overrightarrow{R}}{c\ K}{]}
Таким образом
{[}\nabla \overrightarrow{v}=\frac{\partial \overrightarrow{v}}{\partial t'}\cdot \nabla t'=-\frac{\ \overrightarrow{a}\cdot \overrightarrow{R}}{c\ K}{]}
Далее {[}\nabla \frac{1}{K}=-\frac{1}{{{K}^{2}}}\nabla K{]}
{[}\nabla K=\nabla \left(
R-\frac{\overrightarrow{v}\cdot \overrightarrow{R}}{c}
\right)=\nabla R-\frac{1}{c}\nabla \left(
\overrightarrow{v}\cdot \overrightarrow{R} \right){]} {[}

\begin{align}
  & \nabla R=\frac{\overrightarrow{R}}{R}+\frac{\partial R}{\partial t'}\nabla t'=\frac{\overrightarrow{R}}{R}-\frac{\partial R}{\partial t'}\frac{\ \overrightarrow{R}}{c\ K}=\frac{\overrightarrow{R}}{R}-\left( \frac{\overrightarrow{R}}{R}\cdot \frac{\partial \overrightarrow{R}}{\partial t'} \right)\frac{\ \overrightarrow{R}}{c\ K}=\frac{\overrightarrow{R}}{R}+\left( \frac{\overrightarrow{R}\cdot \overrightarrow{v}}{R} \right)\frac{\ \overrightarrow{R}}{c\ K}=\frac{\overrightarrow{R}}{R}\left( 1+\frac{\overrightarrow{R}\cdot \overrightarrow{v}}{c\ K} \right) \\ 
 & =\frac{\overrightarrow{R}}{R}\left( \frac{c\ K+\overrightarrow{R}\cdot \overrightarrow{v}}{c\ K} \right)=\frac{\overrightarrow{R}}{R}\left( \frac{c\ \left( R-\frac{\overrightarrow{R}\cdot \overrightarrow{v}}{c} \right)+\overrightarrow{R}\cdot \overrightarrow{v}}{c\ K} \right)=\frac{\overrightarrow{R}}{R}\left( \frac{c\ R-\overrightarrow{R}\cdot \overrightarrow{v}+\overrightarrow{R}\cdot \overrightarrow{v}}{c\ K} \right)=\frac{\overrightarrow{R}}{R}\frac{c\,R}{c\,K}=\frac{\overrightarrow{R}\ }{K} \\ 
\end{align}

{]} {[}\nabla \left( \overrightarrow{v}\cdot \overrightarrow{R}
\right)=\overrightarrow{v}\times rot\overrightarrow{R}+\overrightarrow{R}\times rot\overrightarrow{v}+\left(
\overrightarrow{v},\nabla  \right)\overrightarrow{R}+\left(
\overrightarrow{R},\nabla  \right)\overrightarrow{v}{]} Учитывая, что в
декартовой системе координат \({{R}_{x}}=x-x'\left( t' \right)\)
\({{R}_{y}}=y-y'\left( t' \right)\) \({{R}_{x}}=z-z'\left( t' \right)\)
Ищем выражение для ротора радиус вектора
{[}rot\overrightarrow{R}=\left\textbar{}

\begin{matrix}
   \overrightarrow{i} & \overrightarrow{j} & \overrightarrow{k}  \\
   \frac{\partial }{\partial x} & \frac{\partial }{\partial y} & \frac{\partial }{\partial z}  \\
   {{R}_{x}} & {{R}_{y}} & {{R}_{z}}  \\
\end{matrix}

\right\textbar{}=

\begin{matrix}
   \overrightarrow{i}\left( \frac{\partial {{R}_{z}}}{\partial y}-\frac{\partial {{R}_{y}}}{\partial z} \right)  \\
   \overrightarrow{j}\left( \frac{\partial {{R}_{x}}}{\partial z}-\frac{\partial {{R}_{z}}}{\partial x} \right)  \\
   \overrightarrow{k}\left( \frac{\partial {{R}_{y}}}{\partial x}-\frac{\partial {{R}_{x}}}{\partial y} \right)  \\
\end{matrix}

=

\begin{matrix}
   \overrightarrow{i}\left( -\frac{\partial z'}{\partial t'}\frac{\partial t'}{\partial y}+\frac{\partial y'}{\partial t'}\frac{\partial t'}{\partial z} \right)  \\
   \overrightarrow{j}\left( -\frac{\partial x'}{\partial t'}\frac{\partial t'}{\partial z}+\frac{\partial z'}{\partial t'}\frac{\partial t'}{\partial x} \right)  \\
   \overrightarrow{k}\left( -\frac{\partial y'}{\partial t'}\frac{\partial t'}{\partial x}+\frac{\partial x'}{\partial t'}\frac{\partial t'}{\partial y} \right)  \\
\end{matrix}

=

\begin{matrix}
   \frac{\overrightarrow{i}}{c}\left( {{v}_{z}}\frac{\partial R}{\partial y}-{{v}_{y}}\frac{\partial R}{\partial z} \right)  \\
   \frac{\overrightarrow{j}}{c}\left( {{v}_{x}}\frac{\partial R}{\partial z}-{{v}_{z}}\frac{\partial R}{\partial x} \right)  \\
   \frac{\overrightarrow{k}}{c}\left( {{v}_{y}}\frac{\partial R}{\partial x}-{{v}_{x}}\frac{\partial R}{\partial y} \right)  \\
\end{matrix}

{]} Производная \(R\) по координате точки наблюдения
{[}\frac{\partial R}{\partial x}=\frac{\partial }{\partial x}{{\left( \overrightarrow{R}\cdot \overrightarrow{R} \right)}^{{}^{1}/{}_{2}}}=\frac{1}{2}{{\left( \overrightarrow{R}\cdot \overrightarrow{R} \right)}^{-{}^{1}/{}_{2}}}\frac{\partial }{\partial x}\left(
\overrightarrow{R}\cdot \overrightarrow{R} \right)=\frac{1}{2R}\left(
\overrightarrow{R}\cdot \frac{\partial \overrightarrow{R}}{\partial x}+\frac{\partial \overrightarrow{R}}{\partial x}\cdot \overrightarrow{R}
\right)=\frac{\overrightarrow{R}}{R}\cdot \frac{\partial \overrightarrow{R}}{\partial x}{]}
Поскольку в декартовой системе координат
{[}\overrightarrow{R}=\overrightarrow{i}~\{\{R\}\emph{\{x\}\}\left( t'
\right)+\overrightarrow{j}~\{\{R\}}\{y\}\}\left( t'
\right)+\overrightarrow{k}~\{\{R\}\_\{z\}\}\left( t'
\right)=\overrightarrow{i}\left( x-x'\left( t' \right)
\right)+\overrightarrow{j}\left( y-y'\left( t' \right)
\right)+\overrightarrow{k}\left( z-z'\left( t' \right) \right){]}
постольку {[}\frac{\partial \overrightarrow{R}}{\partial x}=

\begin{matrix}
   \overrightarrow{i}\ \frac{\partial {{R}_{x}}\left( t' \right)}{\partial x}  \\
   \overrightarrow{j}\ \frac{\partial {{R}_{y}}\left( t' \right)}{\partial x}  \\
   \overrightarrow{k}\ \frac{\partial {{R}_{z}}\left( t' \right)}{\partial x}  \\
\end{matrix}

=

\begin{matrix}
   \overrightarrow{i}\ \frac{\partial \left( x-x'\left( t' \right) \right)}{\partial x}  \\
   \overrightarrow{j}\ \frac{\partial \left( y-y'\left( t' \right) \right)}{\partial x}  \\
   \overrightarrow{k}\ \frac{\partial \left( z-z'\left( t' \right) \right)}{\partial x}  \\
\end{matrix}

=

\begin{matrix}
   \overrightarrow{i}\ \left( 1-\frac{\partial x'}{\partial t'}\frac{\partial t'}{\partial x} \right)  \\
   -\overrightarrow{j}\ \frac{\partial y'}{\partial t'}\frac{\partial t'}{\partial x}  \\
   -\overrightarrow{k}\ \frac{\partial z'}{\partial t'}\frac{\partial t'}{\partial x}  \\
\end{matrix}

=

\begin{matrix}
   \overrightarrow{i}\ \left( 1+\frac{{{v}_{x}}}{c}\frac{\partial R}{\partial x} \right)  \\
   \overrightarrow{j}\ \frac{{{v}_{y}}}{c}\frac{\partial R}{\partial x}  \\
   \overrightarrow{k}\ \frac{{{v}_{z}}}{c}\frac{\partial R}{\partial x}  \\
\end{matrix}

=\overrightarrow{i}+\frac{\overrightarrow{v}}{c}\frac{\partial R}{\partial x}{]}
Следовательно
{[}\frac{\partial R}{\partial x}=\frac{\overrightarrow{R}}{R}\cdot \frac{\partial \overrightarrow{R}}{\partial x}=\frac{\overrightarrow{R}}{R}\cdot \left(
\overrightarrow{i}+\frac{\overrightarrow{v}}{c}\frac{\partial R}{\partial x}
\right)=\frac{{{R}_{x}}}{R}+\frac{\overrightarrow{R}\cdot \overrightarrow{v}}{Rc}\frac{\partial R}{\partial x}{]}
{[}\frac{\partial R}{\partial x}\left(
1-\frac{\overrightarrow{R}\cdot \overrightarrow{v}}{Rc}
\right)=\frac{{{R}_{x}}}{R}{]} откуда
{[}\frac{\partial R}{\partial x}=\frac{{{R}_{x}}}{R\left( 1-\frac{\overrightarrow{R}\cdot \overrightarrow{v}}{Rc} \right)}=\frac{{{R}_{x}}}{R-\frac{\overrightarrow{R}\cdot \overrightarrow{v}}{c}}=\frac{{{R}_{x}}}{K}{]}
И
{[}\frac{\partial \overrightarrow{R}}{\partial x}=\overrightarrow{i}+\frac{\overrightarrow{v}}{c}\frac{\partial R}{\partial x}=\overrightarrow{i}+\frac{\overrightarrow{v}}{c}\frac{{{R}_{x}}}{K}{]}
Подставляя в выражение для ротора радиус вектора
{[}rot\overrightarrow{R}=

\begin{matrix}
   \frac{\overrightarrow{i}}{c}\left( {{v}_{z}}\frac{\partial R}{\partial y}-{{v}_{y}}\frac{\partial R}{\partial z} \right)  \\
   \frac{\overrightarrow{j}}{c}\left( {{v}_{x}}\frac{\partial R}{\partial z}-{{v}_{z}}\frac{\partial R}{\partial x} \right)  \\
   \frac{\overrightarrow{k}}{c}\left( {{v}_{y}}\frac{\partial R}{\partial x}-{{v}_{x}}\frac{\partial R}{\partial y} \right)  \\
\end{matrix}

=

\begin{matrix}
   \frac{\overrightarrow{i}}{cK}\left( {{v}_{z}}{{R}_{y}}-{{v}_{y}}{{R}_{z}} \right)  \\
   \frac{\overrightarrow{j}}{cK}\left( {{v}_{x}}{{R}_{z}}-{{v}_{z}}{{R}_{x}} \right)  \\
   \frac{\overrightarrow{k}}{cK}\left( {{v}_{y}}{{R}_{x}}-{{v}_{x}}{{R}_{y}} \right)  \\
\end{matrix}

{]} Теперь находим векторное произведение запаздывающей скорости на
ротор радиус вектора {[}cK\left(
\overrightarrow{v}\times rot\overrightarrow{R} \right)=\left\textbar{}

\begin{matrix}
   \overrightarrow{i} & \overrightarrow{j} & \overrightarrow{k}  \\
   {{v}_{x}} & {{v}_{y}} & {{v}_{z}}  \\
   {{v}_{z}}{{R}_{y}}-{{v}_{y}}{{R}_{z}} & {{v}_{x}}{{R}_{z}}-{{v}_{z}}{{R}_{x}} & {{v}_{y}}{{R}_{x}}-{{v}_{x}}{{R}_{y}}  \\
\end{matrix}

\right\textbar{}=

\begin{matrix}
   \overrightarrow{i}\left( {{v}_{y}}\left( {{v}_{y}}{{R}_{x}}-{{v}_{x}}{{R}_{y}} \right)-{{v}_{z}}\left( {{v}_{x}}{{R}_{z}}-{{v}_{z}}{{R}_{x}} \right) \right)  \\
   \overrightarrow{j}\left( {{v}_{z}}\left( {{v}_{z}}{{R}_{y}}-{{v}_{y}}{{R}_{z}} \right)-{{v}_{x}}\left( {{v}_{y}}{{R}_{x}}-{{v}_{x}}{{R}_{y}} \right) \right)  \\
   \overrightarrow{k}\left( {{v}_{x}}\left( {{v}_{x}}{{R}_{z}}-{{v}_{z}}{{R}_{x}} \right)-{{v}_{y}}\left( {{v}_{z}}{{R}_{y}}-{{v}_{y}}{{R}_{z}} \right) \right)  \\
\end{matrix}

{]} {[}=

\begin{matrix}
   \overrightarrow{i}\left( {{R}_{x}}\left( {{v}_{y}}^{2}+{{v}_{z}}^{2} \right)-{{v}_{x}}{{v}_{y}}{{R}_{y}}-{{v}_{x}}{{v}_{z}}{{R}_{z}} \right)  \\
   \overrightarrow{j}\left( {{R}_{y}}\left( {{v}_{x}}^{2}+{{v}_{z}}^{2} \right)-{{v}_{y}}{{v}_{x}}{{R}_{x}}-{{v}_{y}}{{v}_{z}}{{R}_{z}} \right)  \\
   \overrightarrow{k}\left( {{R}_{z}}\left( {{v}_{x}}^{2}+{{v}_{y}}^{2} \right)-{{v}_{z}}{{v}_{x}}{{R}_{x}}-{{v}_{z}}{{v}_{y}}{{R}_{y}} \right)  \\
\end{matrix}

=

\begin{matrix}
   \overrightarrow{i}\left( {{R}_{x}}\left( {{v}^{2}}-{{v}_{x}}^{2} \right)-{{v}_{x}}\left( {{v}_{y}}{{R}_{y}}-{{v}_{z}}{{R}_{z}} \right) \right)  \\
   \overrightarrow{j}\left( {{R}_{y}}\left( {{v}^{2}}-{{v}_{y}}^{2} \right)-{{v}_{y}}\left( {{v}_{x}}{{R}_{x}}-{{v}_{z}}{{R}_{z}} \right) \right)  \\
   \overrightarrow{k}\left( {{R}_{z}}\left( {{v}^{2}}-{{v}_{z}}^{2} \right)-{{v}_{z}}\left( {{v}_{x}}{{R}_{x}}-{{v}_{y}}{{R}_{y}} \right) \right)  \\
\end{matrix}

=

\begin{matrix}
   \overrightarrow{i}\left( {{R}_{x}}{{v}^{2}}-{{v}_{x}}\left( \overrightarrow{v}\cdot \overrightarrow{R} \right) \right)  \\
   \overrightarrow{j}\left( {{R}_{y}}{{v}^{2}}-{{v}_{y}}\left( \overrightarrow{v}\cdot \overrightarrow{R} \right) \right)  \\
   \overrightarrow{k}\left( {{R}_{z}}{{v}^{2}}-{{v}_{z}}\left( \overrightarrow{v}\cdot \overrightarrow{R} \right) \right)  \\
\end{matrix}

{]}
{[}\overrightarrow{v}\times rot\overrightarrow{R}=\frac{{{v}^{2}}\ \overrightarrow{R}-\left( \overrightarrow{v}\cdot \overrightarrow{R} \right)\ \overrightarrow{v}}{cK}{]}
Теперь ищем выражение для ротора запаздывающей скорости
{[}rot~\overrightarrow{v}=

\begin{matrix}
   \overrightarrow{i}\left( \frac{\partial {{v}_{z}}}{\partial y}-\frac{\partial {{v}_{y}}}{\partial z} \right)  \\
   \overrightarrow{j}\left( \frac{\partial {{v}_{x}}}{\partial z}-\frac{\partial {{v}_{z}}}{\partial x} \right)  \\
   \overrightarrow{k}\left( \frac{\partial {{v}_{y}}}{\partial x}-\frac{\partial {{v}_{x}}}{\partial y} \right)  \\
\end{matrix}

=

\begin{matrix}
   \overrightarrow{i}\left( \frac{\partial {{v}_{z}}}{\partial t'}\frac{\partial t'}{\partial y}-\frac{\partial {{v}_{y}}}{\partial t'}\frac{\partial t'}{\partial z} \right)  \\
   \overrightarrow{j}\left( \frac{\partial {{v}_{x}}}{\partial t'}\frac{\partial t'}{\partial z}-\frac{\partial {{v}_{z}}}{\partial t'}\frac{\partial t'}{\partial x} \right)  \\
   \overrightarrow{k}\left( \frac{\partial {{v}_{y}}}{\partial t'}\frac{\partial t'}{\partial x}-\frac{\partial {{v}_{x}}}{\partial t'}\frac{\partial t'}{\partial y} \right)  \\
\end{matrix}

=

\begin{matrix}
   \frac{\overrightarrow{i}}{c}\left( -{{a}_{z}}\frac{\partial R}{\partial y}+{{a}_{y}}\frac{\partial R}{\partial z} \right)  \\
   \frac{\overrightarrow{j}}{c}\left( -{{a}_{x}}\frac{\partial R}{\partial z}+{{a}_{z}}\frac{\partial R}{\partial x} \right)  \\
   \frac{\overrightarrow{k}}{c}\left( -{{a}_{y}}\frac{\partial R}{\partial x}+{{a}_{x}}\frac{\partial R}{\partial y} \right)  \\
\end{matrix}

=

\begin{matrix}
   \frac{\overrightarrow{i}}{cK}\left( -{{a}_{z}}{{R}_{y}}+{{a}_{y}}{{R}_{z}} \right)  \\
   \frac{\overrightarrow{j}}{cK}\left( -{{a}_{x}}{{R}_{z}}+{{a}_{z}}{{R}_{x}} \right)  \\
   \frac{\overrightarrow{k}}{cK}\left( -{{a}_{y}}{{R}_{x}}+{{a}_{x}}{{R}_{y}} \right)  \\
\end{matrix}

{]} Далее векторное произведение радиус-вектора на ротор запаздывающей
скорости {[}cK\left( \overrightarrow{R}\times rot~\overrightarrow{v}
\right)=\left\textbar{}

\begin{matrix}
   \overrightarrow{i} & \overrightarrow{j} & \overrightarrow{k}  \\
   {{R}_{x}} & {{R}_{y}} & {{R}_{z}}  \\
   {{a}_{y}}{{R}_{z}}-{{a}_{z}}{{R}_{y}} & {{a}_{z}}{{R}_{x}}-{{a}_{x}}{{R}_{z}} & {{a}_{x}}{{R}_{y}}-{{a}_{y}}{{R}_{x}}  \\
\end{matrix}

\right\textbar{}=

\begin{matrix}
   \overrightarrow{i}\left( {{R}_{y}}\left( {{a}_{x}}{{R}_{y}}-{{a}_{y}}{{R}_{x}} \right)-{{R}_{z}}\left( {{a}_{z}}{{R}_{x}}-{{a}_{x}}{{R}_{z}} \right) \right)  \\
   \overrightarrow{j}\left( {{R}_{z}}\left( {{a}_{y}}{{R}_{z}}-{{a}_{z}}{{R}_{y}} \right)-{{R}_{x}}\left( {{a}_{x}}{{R}_{y}}-{{a}_{y}}{{R}_{x}} \right) \right)  \\
   \overrightarrow{k}\left( {{R}_{x}}\left( {{a}_{z}}{{R}_{x}}-{{a}_{x}}{{R}_{z}} \right)-{{R}_{y}}\left( {{a}_{y}}{{R}_{z}}-{{a}_{z}}{{R}_{y}} \right) \right)  \\
\end{matrix}

{]} {[}=

\begin{matrix}
   \overrightarrow{i}\left( {{a}_{x}}\left( {{R}_{y}}^{2}+{{R}_{z}}^{2} \right)-{{R}_{x}}{{R}_{y}}{{a}_{y}}-{{R}_{x}}{{R}_{z}}{{a}_{z}} \right)  \\
   \overrightarrow{j}\left( {{a}_{y}}\left( {{R}_{x}}^{2}+{{R}_{z}}^{2} \right)-{{R}_{y}}{{R}_{x}}{{a}_{x}}-{{R}_{y}}{{R}_{z}}{{a}_{z}} \right)  \\
   \overrightarrow{k}\left( {{a}_{z}}\left( {{R}_{x}}^{2}+{{R}_{y}}^{2} \right)-{{R}_{z}}{{R}_{x}}{{a}_{x}}-{{R}_{z}}{{R}_{y}}{{a}_{y}} \right)  \\
\end{matrix}

=

\begin{matrix}
   \overrightarrow{i}\left( {{a}_{x}}\left( {{R}^{2}}-{{R}_{x}}^{2} \right)-{{R}_{x}}\left( {{a}_{y}}{{R}_{y}}+{{a}_{z}}{{R}_{z}} \right) \right)  \\
   \overrightarrow{j}\left( {{a}_{y}}\left( {{R}^{2}}-{{R}_{y}}^{2} \right)-{{R}_{y}}\left( {{a}_{x}}{{R}_{x}}+{{a}_{z}}{{R}_{z}} \right) \right)  \\
   \overrightarrow{k}\left( {{a}_{z}}\left( {{R}^{2}}-{{R}_{z}}^{2} \right)-{{R}_{z}}\left( {{a}_{x}}{{R}_{x}}+{{a}_{y}}{{R}_{y}} \right) \right)  \\
\end{matrix}

=

\begin{matrix}
   \overrightarrow{i}\left( {{a}_{x}}{{R}^{2}}-{{R}_{x}}\left( \overrightarrow{a}\cdot \overrightarrow{R} \right) \right)  \\
   \overrightarrow{j}\left( {{a}_{y}}{{R}^{2}}-{{R}_{y}}\left( \overrightarrow{a}\cdot \overrightarrow{R} \right) \right)  \\
   \overrightarrow{k}\left( {{a}_{z}}{{R}^{2}}-{{R}_{z}}\left( \overrightarrow{a}\cdot \overrightarrow{R} \right) \right)  \\
\end{matrix}

{]}
{[}\overrightarrow{R}\times rot~\overrightarrow{v}=\frac{{{R}^{2}}\overrightarrow{a}-\left( \overrightarrow{a}\cdot \overrightarrow{R} \right)\overrightarrow{R}}{cK}{]}
Субстанционная производная радиус вектора {[}\left(
\overrightarrow{v},\nabla  \right)\overrightarrow{R}=\{\{v\}\emph{\{x\}\}\frac{\partial }{\partial x}\overrightarrow{R}+\{\{v\}}\{y\}\}\frac{\partial }{\partial y}\overrightarrow{R}+\{\{v\}\_\{z\}\}\frac{\partial }{\partial z}\overrightarrow{R}{]}
Учитывая ранее найденные выражения производной радиус вектора по
координате
{[}\frac{\partial \overrightarrow{R}}{\partial x}=\overrightarrow{i}+\frac{\overrightarrow{v}}{c}\frac{\partial R}{\partial x}{]}
и если учесть что
{[}\frac{\partial R}{\partial x}=\frac{{{R}_{x}}}{K}{]} то можно
записать
{[}\frac{\partial \overrightarrow{R}}{\partial x}=\overrightarrow{i}+\frac{\overrightarrow{v}}{c}\frac{{{R}_{x}}}{K}{]}
{[}

\begin{align}
  & \left( \overrightarrow{v},\nabla  \right)\overrightarrow{R}={{v}_{x}}\left( \overrightarrow{i}+\frac{\overrightarrow{v}}{c}\frac{{{R}_{x}}}{K} \right)+{{v}_{y}}\left( \overrightarrow{j}+\frac{\overrightarrow{v}}{c}\frac{{{R}_{y}}}{K} \right)+{{v}_{z}}\left( \overrightarrow{k}+\frac{\overrightarrow{v}}{c}\frac{{{R}_{z}}}{K} \right)=\overrightarrow{v}+\frac{\overrightarrow{v}}{c}\frac{\left( \overrightarrow{v}\cdot \overrightarrow{R} \right)}{K}=\overrightarrow{v}\left( 1+\frac{\left( \overrightarrow{v}\cdot \overrightarrow{R} \right)}{cK} \right) \\ 
 & =\overrightarrow{v}\left( \frac{c\left( R-\frac{\overrightarrow{v}\cdot \overrightarrow{R}}{c} \right)}{cK}+\frac{\left( \overrightarrow{v}\cdot \overrightarrow{R} \right)}{cK} \right)=\overrightarrow{v}\left( \frac{cR-\overrightarrow{v}\cdot \overrightarrow{R}}{cK}+\frac{\left( \overrightarrow{v}\cdot \overrightarrow{R} \right)}{cK} \right)=\overrightarrow{v}\left( \frac{cR}{cK} \right)=\overrightarrow{v}\frac{R}{K} \\ 
\end{align}

{]} Теперь {[}\left(
\overrightarrow{R},\nabla  \right)\overrightarrow{v}=\{\{R\}\emph{\{x\}\}\frac{\partial }{\partial x}\overrightarrow{v}+\{\{R\}}\{y\}\}\frac{\partial }{\partial y}\overrightarrow{v}+\{\{R\}\emph{\{z\}\}\frac{\partial }{\partial z}\overrightarrow{v}{]}
{[}\frac{\partial }{\partial x}\overrightarrow{v}=\frac{\partial \overrightarrow{v}}{\partial t'}\frac{\partial t'}{\partial x}=\overrightarrow{a}\frac{\partial t'}{\partial x}=-\frac{\overrightarrow{a}}{c}\frac{\partial R}{\partial x}=-\frac{\overrightarrow{a}}{c}\frac{{{R}_{x}}}{K}{]}
{[}\left(
\overrightarrow{R},\nabla  \right)\overrightarrow{v}=-\{\{R\}}\{x\}\}\frac{\overrightarrow{a}}{c}\frac{{{R}_{x}}}{K}-\{\{R\}\emph{\{y\}\}\frac{\overrightarrow{a}}{c}\frac{{{R}_{y}}}{K}-\{\{R\}}\{z\}\}\frac{\overrightarrow{a}}{c}\frac{{{R}_{z}}}{K}=-\frac{\overrightarrow{a}}{c}\frac{{{R}^{2}}}{K}{]}
Таким образом, градиент скалярного произведения скорости на радиус
вектор {[}

\begin{align}
  & \nabla \left( \overrightarrow{v}\cdot \overrightarrow{R} \right)=\overrightarrow{v}\times rot\overrightarrow{R}+\overrightarrow{R}\times rot\overrightarrow{v}+\left( \overrightarrow{v},\nabla  \right)\overrightarrow{R}+\left( \overrightarrow{R},\nabla  \right)\overrightarrow{v}= \\ 
 & =\left( \frac{{{v}^{2}}\ \overrightarrow{R}-\left( \overrightarrow{v}\cdot \overrightarrow{R} \right)\ \overrightarrow{v}}{cK} \right)+\left( \frac{{{R}^{2}}\overrightarrow{a}-\left( \overrightarrow{a}\cdot \overrightarrow{R} \right)\overrightarrow{R}}{cK} \right)+\left( \overrightarrow{v}+\frac{\overrightarrow{v}}{c}\frac{\left( \overrightarrow{v}\cdot \overrightarrow{R} \right)}{K} \right)+\left( -\frac{\overrightarrow{a}}{c}\frac{{{R}^{2}}}{K} \right) \\ 
 & =\left( \frac{{{v}^{2}}\ \overrightarrow{R}}{cK} \right)-\left( \frac{\left( \overrightarrow{a}\cdot \overrightarrow{R} \right)\overrightarrow{R}}{cK} \right)+\overrightarrow{v} \\ 
\end{align}

{]} Интересно, что этот градиент можно также найти значительно более
коротким путём - через сумму производных по координате {[}\nabla \left(
\overrightarrow{v}\cdot \overrightarrow{R}
\right)=\overrightarrow{i}\frac{\partial }{\partial x}\left(
\overrightarrow{v}\cdot \overrightarrow{R}
\right)+\overrightarrow{j}\frac{\partial }{\partial y}\left(
\overrightarrow{v}\cdot \overrightarrow{R}
\right)+\overrightarrow{k}\frac{\partial }{\partial z}\left(
\overrightarrow{v}\cdot \overrightarrow{R} \right){]} производная по
координате {[}\frac{\partial }{\partial x}\left(
\overrightarrow{v}\cdot \overrightarrow{R}
\right)=\overrightarrow{v}\frac{\partial \overrightarrow{R}}{\partial x}+\overrightarrow{R}\frac{\partial \overrightarrow{v}}{\partial x}{]}
{[}\frac{\partial \overrightarrow{v}}{\partial x}=\frac{\partial \overrightarrow{v}}{\partial t'}\frac{\partial t'}{\partial x}=\overrightarrow{a}\left(
\frac{\partial }{\partial x}\left\{ t-\frac{R}{c} \right\}
\right)=\overrightarrow{a}\left(
-\frac{1}{c}\frac{\partial R}{\partial x}
\right)=-\frac{\overrightarrow{a}}{c}\frac{{{R}_{x}}}{K}{]}
{[}\frac{\partial }{\partial x}\left(
\overrightarrow{v}\cdot \overrightarrow{R}
\right)=\overrightarrow{v}\left(
\overrightarrow{i}+\frac{\overrightarrow{v}}{c}\frac{{{R}_{x}}}{K}
\right)+\overrightarrow{R}\left(
-\frac{\overrightarrow{a}}{c}\frac{{{R}_{x}}}{K}
\right)=\{\{v\}\_\{x\}\}+\left(
\frac{{{v}^{2}}-\overrightarrow{R}\cdot \overrightarrow{a}}{c}
\right)\frac{{{R}_{x}}}{K}{]} Складывая три компоненты приходим к тому
же выражению {[}\nabla \left( \overrightarrow{v}\cdot \overrightarrow{R}
\right)=\overrightarrow{v}+\left(
\frac{{{v}^{2}}-\overrightarrow{R}\cdot \overrightarrow{a}}{c}
\right)\frac{\overrightarrow{R}}{K}{]} Откуда
{[}\nabla K=\nabla R-\frac{1}{c}\nabla \left(
\overrightarrow{v}\cdot \overrightarrow{R} \right)=\left(
\frac{\overrightarrow{R}\ }{K} \right)-\frac{1}{c}\left\{
\overrightarrow{v}+\left(
\frac{{{v}^{2}}-\overrightarrow{R}\cdot \overrightarrow{a}}{c}
\right)\frac{\overrightarrow{R}}{K}
\right\}=\frac{\overrightarrow{R}\ }{K}\left(
1+\frac{\overrightarrow{a}\cdot \overrightarrow{R}-{{v}^{2}}}{{{c}^{2}}}
\right)-\frac{\overrightarrow{v}}{c}{]} теперь
{[}\nabla \frac{1}{\left( R-\frac{\overrightarrow{v}\cdot \overrightarrow{R}}{c} \right)}=-\frac{1}{{{\left( R-\frac{\overrightarrow{v}\cdot \overrightarrow{R}}{c} \right)}^{2}}}\left\{
\frac{\overrightarrow{R}}{\left( R-\frac{\overrightarrow{v}\cdot \overrightarrow{R}}{c} \right)}\left(
1+\frac{\overrightarrow{a}\cdot \overrightarrow{R}}{{{c}^{2}}}-\frac{\overrightarrow{v}\cdot \overrightarrow{v}}{{{c}^{2}}}
\right)-\frac{\overrightarrow{v}}{c} \right\}{]} Откуда
{[}\nabla \overrightarrow{{{A}_{H}}}=\frac{e}{c}\left\{
-\frac{\overrightarrow{a}\cdot \overrightarrow{R}}{c{{\left( R-\frac{\overrightarrow{v}\cdot \overrightarrow{R}}{c} \right)}^{2}}}-\frac{\overrightarrow{v}}{{{\left( R-\frac{\overrightarrow{v}\cdot \overrightarrow{R}}{c} \right)}^{2}}}\cdot \left\{
\frac{\overrightarrow{R}}{\left( R-\frac{\overrightarrow{v}\cdot \overrightarrow{R}}{c} \right)}\left(
1+\frac{\overrightarrow{a}\cdot \overrightarrow{R}}{{{c}^{2}}}-\frac{\overrightarrow{v}\cdot \overrightarrow{v}}{{{c}^{2}}}
\right)-\frac{\overrightarrow{v}}{c} \right\} \right\}{]} На слагаемые
{[}\nabla \overrightarrow{{{A}_{H}}}=\frac{e}{c}\left\{
-\frac{\overrightarrow{a}\cdot \overrightarrow{R}}{c{{\left( R-\frac{\overrightarrow{v}\cdot \overrightarrow{R}}{c} \right)}^{2}}}-\frac{\overrightarrow{v}\cdot \overrightarrow{R}}{{{\left( R-\frac{\overrightarrow{v}\cdot \overrightarrow{R}}{c} \right)}^{3}}}\left(
1+\frac{\overrightarrow{a}\cdot \overrightarrow{R}}{{{c}^{2}}}-\frac{\overrightarrow{v}\cdot \overrightarrow{v}}{{{c}^{2}}}
\right)+\frac{\overrightarrow{v}\cdot \overrightarrow{v}}{{{\left( R-\frac{\overrightarrow{v}\cdot \overrightarrow{R}}{c} \right)}^{2}}}
\right\}{]} {[}\nabla \overrightarrow{{{A}_{H}}}=\frac{e}{c}\left\{
\frac{\overrightarrow{v}\cdot \overrightarrow{v}-\overrightarrow{a}\cdot \overrightarrow{R}}{c{{\left( R-\frac{\overrightarrow{v}\cdot \overrightarrow{R}}{c} \right)}^{2}}}-\frac{\overrightarrow{v}\cdot \overrightarrow{R}}{{{\left( R-\frac{\overrightarrow{v}\cdot \overrightarrow{R}}{c} \right)}^{3}}}\left(
1+\frac{\overrightarrow{a}\cdot \overrightarrow{R}}{{{c}^{2}}}-\frac{\overrightarrow{v}\cdot \overrightarrow{v}}{{{c}^{2}}}
\right) \right\}{]} Теперь проверяем полученные выражения на
соответствие калибровке Лоренца
{[}div~\overrightarrow{{{A}_{H}}}=-\frac{\varepsilon }{c}\frac{\partial \varphi }{\partial t}{]}
Правая часть этого выражения
{[}-\frac{\varepsilon }{c}\frac{\partial \varphi }{\partial t}=\frac{e}{}\frac{R\frac{\vec{v}\cdot \vec{v}-\overrightarrow{a}\cdot \overrightarrow{R}}{c}-~\overrightarrow{v}\cdot \overrightarrow{R}}{{{\left( R-\frac{\overrightarrow{v}\cdot \overrightarrow{R}}{c} \right)}^{3}}}{]}
Итак, проверим тождество
{[}\frac{\overrightarrow{v}\cdot \overrightarrow{v}-\overrightarrow{a}\cdot \overrightarrow{R}}{c{{\left( R-\frac{\overrightarrow{v}\cdot \overrightarrow{R}}{c} \right)}^{2}}}-\frac{\overrightarrow{v}\cdot \overrightarrow{R}}{{{\left( R-\frac{\overrightarrow{v}\cdot \overrightarrow{R}}{c} \right)}^{3}}}\left(
1+\frac{\overrightarrow{a}\cdot \overrightarrow{R}}{{{c}^{2}}}-\frac{\overrightarrow{v}\cdot \overrightarrow{v}}{{{c}^{2}}}
\right)=\frac{R\frac{\vec{v}\cdot \vec{v}-\overrightarrow{a}\cdot \overrightarrow{R}}{c}-~\overrightarrow{v}\cdot \overrightarrow{R}}{{{\left( R-\frac{\overrightarrow{v}\cdot \overrightarrow{R}}{c} \right)}^{3}}}{]}

{[}\frac{\overrightarrow{v}\cdot \overrightarrow{v}-\overrightarrow{a}\cdot \overrightarrow{R}}{c{{\left( R-\frac{\overrightarrow{v}\cdot \overrightarrow{R}}{c} \right)}^{2}}}-\frac{\overrightarrow{v}\cdot \overrightarrow{R}}{{{\left( R-\frac{\overrightarrow{v}\cdot \overrightarrow{R}}{c} \right)}^{3}}}\left(
1+\frac{\overrightarrow{a}\cdot \overrightarrow{R}}{{{c}^{2}}}-\frac{\overrightarrow{v}\cdot \overrightarrow{v}}{{{c}^{2}}}
\right)=\frac{R\left( \vec{v}\cdot \vec{v}-\overrightarrow{a}\cdot \overrightarrow{R} \right)}{{{\left( R-\frac{\overrightarrow{v}\cdot \overrightarrow{R}}{c} \right)}^{3}}}-\frac{~\overrightarrow{v}\cdot \overrightarrow{R}}{{{\left( R-\frac{\overrightarrow{v}\cdot \overrightarrow{R}}{c} \right)}^{3}}}{]}

{[}\frac{\overrightarrow{v}\cdot \overrightarrow{v}-\overrightarrow{a}\cdot \overrightarrow{R}}{c{{\left( R-\frac{\overrightarrow{v}\cdot \overrightarrow{R}}{c} \right)}^{2}}}-\frac{\overrightarrow{v}\cdot \overrightarrow{R}}{{{\left( R-\frac{\overrightarrow{v}\cdot \overrightarrow{R}}{c} \right)}^{3}}}\left(
\frac{\overrightarrow{a}\cdot \overrightarrow{R}}{{{c}^{2}}}-\frac{\overrightarrow{v}\cdot \overrightarrow{v}}{{{c}^{2}}}
\right)=\frac{R\left( \vec{v}\cdot \vec{v}-\overrightarrow{a}\cdot \overrightarrow{R} \right)}{{{\left( R-\frac{\overrightarrow{v}\cdot \overrightarrow{R}}{c} \right)}^{3}}}{]}

{[}\frac{\overrightarrow{v}\cdot \overrightarrow{v}-\overrightarrow{a}\cdot \overrightarrow{R}}{c{{\left( R-\frac{\overrightarrow{v}\cdot \overrightarrow{R}}{c} \right)}^{2}}}+\frac{\overrightarrow{v}\cdot \overrightarrow{R}}{{{\left( R-\frac{\overrightarrow{v}\cdot \overrightarrow{R}}{c} \right)}^{3}}}\left(
\frac{\overrightarrow{v}\cdot \overrightarrow{v}-\overrightarrow{a}\cdot \overrightarrow{R}}{{{c}^{2}}}
\right)=\frac{R\left( \vec{v}\cdot \vec{v}-\overrightarrow{a}\cdot \overrightarrow{R} \right)}{{{\left( R-\frac{\overrightarrow{v}\cdot \overrightarrow{R}}{c} \right)}^{3}}}{]}
Сокращаем на
{[}\frac{\overrightarrow{v}\cdot \overrightarrow{v}-\overrightarrow{a}\cdot \overrightarrow{R}}{{{c}^{2}}}{]}
{[}\frac{c}{{{\left( R-\frac{\overrightarrow{v}\cdot \overrightarrow{R}}{c} \right)}^{2}}}+\frac{\overrightarrow{v}\cdot \overrightarrow{R}}{{{\left( R-\frac{\overrightarrow{v}\cdot \overrightarrow{R}}{c} \right)}^{3}}}=\frac{cR}{{{\left( R-\frac{\overrightarrow{v}\cdot \overrightarrow{R}}{c} \right)}^{3}}}{]}
Приводим к общему знаменателю
{[}\frac{c\left( R-\frac{\overrightarrow{v}\cdot \overrightarrow{R}}{c} \right)}{{{\left( R-\frac{\overrightarrow{v}\cdot \overrightarrow{R}}{c} \right)}^{3}}}+\frac{\overrightarrow{v}\cdot \overrightarrow{R}}{{{\left( R-\frac{\overrightarrow{v}\cdot \overrightarrow{R}}{c} \right)}^{3}}}=\frac{cR}{{{\left( R-\frac{\overrightarrow{v}\cdot \overrightarrow{R}}{c} \right)}^{3}}}{]}

{[}\frac{cR-\overrightarrow{v}\cdot \overrightarrow{R}+\overrightarrow{v}\cdot \overrightarrow{R}}{{{\left( R-\frac{\overrightarrow{v}\cdot \overrightarrow{R}}{c} \right)}^{3}}}=\frac{cR}{{{\left( R-\frac{\overrightarrow{v}\cdot \overrightarrow{R}}{c} \right)}^{3}}}{]}
Тождество доказано

О конвективной производной вектора электрического поля Николаев в своей
работе {[} CITATION ГНи \l 1049 {]} приводит формулу 1.13 для выражающую
ток смещения через конвективную производную вектора электрического поля
\(\frac{\partial \overrightarrow{E}}{\partial t}=-\left( \overrightarrow{v},\nabla \right)\overrightarrow{E}\)
При попытке дать этой формуле осмысленную физическую интерпретацию
возникает вопрос, что следует понимать под скоростью
\(\overrightarrow{v}\)в данном случае? Скорость заряда? Или может быть
скорость поля? Для прояснения этого вопроса представляет интерес
проверка данной формулы непосредственным дифференцированием: а) формулы
закона Кулона б) градиента скалярного потенциала Лиенара-Вихерта Минус
градиент скалярного потенциала ЛВ
{[}\overrightarrow{{{E}_{1}}}=-\nabla \varphi =-\nabla \frac{e}{K}=\frac{e}{{{K}^{2}}}\left\{
\frac{\overrightarrow{R}}{K}\left(
1+\frac{\overrightarrow{a}\cdot \overrightarrow{R}}{{{c}^{2}}}-\frac{\overrightarrow{v}\cdot \overrightarrow{v}}{{{c}^{2}}}
\right)-\frac{\overrightarrow{v}}{c} \right\}{]}
{[}\overrightarrow{{{E}_{1}}}=e\frac{\overrightarrow{R}}{{{K}^{3}}}\left(
1+\frac{\overrightarrow{a}\cdot \overrightarrow{R}}{{{c}^{2}}}-\frac{\overrightarrow{v}\cdot \overrightarrow{v}}{{{c}^{2}}}
\right)-\frac{e}{{{K}^{2}}}\frac{\overrightarrow{v}}{c}{]} {[}

\begin{align}
  & \frac{1}{e}\frac{\partial \overrightarrow{{{E}_{1}}}}{\partial t}=\frac{\partial }{\partial t}\left[ \frac{\overrightarrow{R}}{{{K}^{3}}}\left( 1+\frac{\overrightarrow{a}\cdot \overrightarrow{R}}{{{c}^{2}}}-\frac{\overrightarrow{v}\cdot \overrightarrow{v}}{{{c}^{2}}} \right)-\frac{1}{{{K}^{2}}}\frac{\overrightarrow{v}}{c} \right]= \\ 
 & =\frac{\overrightarrow{R}}{{{K}^{3}}}\frac{\partial }{\partial t}\left( 1+\frac{\overrightarrow{a}\cdot \overrightarrow{R}}{{{c}^{2}}}-\frac{\overrightarrow{v}\cdot \overrightarrow{v}}{{{c}^{2}}} \right)+\left( 1+\frac{\overrightarrow{a}\cdot \overrightarrow{R}}{{{c}^{2}}}-\frac{\overrightarrow{v}\cdot \overrightarrow{v}}{{{c}^{2}}} \right)\cdot \frac{\partial }{\partial t}\left( \frac{\overrightarrow{R}}{{{K}^{3}}} \right)-\frac{\partial }{\partial t}\left( \frac{1}{{{K}^{2}}}\frac{\overrightarrow{v}}{c} \right) \\ 
\end{align}

{]}

{[}\frac{\partial }{\partial t}\left(
\frac{\overrightarrow{R}}{{{K}^{3}}}
\right)=\frac{1}{{{K}^{3}}}\frac{\partial \overrightarrow{R}}{\partial t}+\overrightarrow{R}\frac{\partial }{\partial t}\left(
\frac{1}{{{K}^{3}}}
\right)=\frac{1}{{{K}^{3}}}\frac{\partial \overrightarrow{R}}{\partial t'}\frac{\partial t'}{\partial t}-\frac{3\overrightarrow{R}}{{{K}^{4}}}\frac{\partial K}{\partial t}=-\frac{\overrightarrow{v}}{{{K}^{3}}}\frac{R}{K}-\frac{3\overrightarrow{R}}{{{K}^{4}}}\frac{\partial K}{\partial t'}\frac{\partial t'}{\partial t}=-\frac{\overrightarrow{v}R}{{{K}^{4}}}-\frac{3\overrightarrow{R}R}{{{K}^{5}}}\frac{\partial K}{\partial t'}{]}
{[}\frac{\partial K}{\partial t'}=\frac{\partial }{\partial t'}\left(
R-\frac{\overrightarrow{v}\cdot \overrightarrow{R}}{c}
\right)=\frac{\overrightarrow{v}\cdot \overrightarrow{v}-\overrightarrow{R}\cdot \overrightarrow{a}}{c}-\frac{\overrightarrow{R}\cdot \overrightarrow{v}}{R}{]}
{[}\frac{\partial }{\partial t}\left(
\frac{\overrightarrow{R}}{{{K}^{3}}}
\right)=-\frac{\overrightarrow{v}R}{{{K}^{4}}}-\frac{3\overrightarrow{R}R}{{{K}^{5}}}\left(
\frac{\overrightarrow{v}\cdot \overrightarrow{v}-\overrightarrow{R}\cdot \overrightarrow{a}}{c}-\frac{\overrightarrow{R}\cdot \overrightarrow{v}}{R}
\right){]} Тут исправлена ошибка!!!!! {[}

\begin{align}
  & \frac{\partial }{\partial t}\left( \frac{1}{{{K}^{2}}}\frac{\overrightarrow{v}}{c} \right)=\frac{\overrightarrow{v}}{c}\frac{\partial }{\partial t}\left( \frac{1}{{{K}^{2}}} \right)+\frac{1}{{{K}^{2}}}\frac{\partial }{\partial t}\left( \frac{\overrightarrow{v}}{c} \right)=-\frac{2}{{{K}^{3}}}\frac{\overrightarrow{v}}{c}\frac{\partial K}{\partial t'}\frac{\partial t'}{\partial t}+\frac{\overrightarrow{a}}{c{{K}^{2}}}\frac{\partial t'}{\partial t}=-\frac{2R}{{{K}^{4}}}\frac{\overrightarrow{v}}{c}\left( \frac{\overrightarrow{v}\cdot \overrightarrow{v}-\overrightarrow{R}\cdot \overrightarrow{a}}{c}-\frac{\overrightarrow{R}\cdot \overrightarrow{v}}{R} \right)+\frac{R\overrightarrow{a}}{c{{K}^{3}}} \\ 
 &  \\ 
\end{align}

{]} Третья промежуточная производная
{[}\frac{\partial }{\partial t}\left(
\frac{1}{{{R}^{*}}^{2}}\frac{\overrightarrow{v}}{c}
\right)=\frac{\overrightarrow{v}}{c}\frac{\partial }{\partial t}\left(
\frac{1}{{{R}^{*}}^{2}}
\right)+\frac{1}{{{R}^{*}}^{2}}\frac{\partial }{\partial t}\left(
\frac{\overrightarrow{v}}{c} \right)=\left(
-\frac{2}{{{R}^{*}}^{3}}\frac{\overrightarrow{v}}{c}\frac{\partial {{R}^{*}}}{\partial t'}+\frac{1}{c{{R}^{*}}^{2}}\frac{\partial \overrightarrow{v}}{\partial t'}
\right)\frac{\partial t'}{\partial t}=\left\{
-\frac{2}{{{R}^{*}}^{3}}\frac{\overrightarrow{v}}{c}\left(
\frac{\overrightarrow{v}\cdot \overrightarrow{v}-\overrightarrow{R}\cdot \overrightarrow{a}}{c}-\frac{\overrightarrow{R}\cdot \overrightarrow{v}}{R}
\right)+\frac{\overrightarrow{a}}{c{{R}^{*}}^{2}}
\right\}\frac{R}{{{R}^{*}}}{]} {[}

\begin{align}
  & \frac{\partial }{\partial t}\left( 1+\frac{\overrightarrow{a}\cdot \overrightarrow{R}}{{{c}^{2}}}-\frac{\overrightarrow{v}\cdot \overrightarrow{v}}{{{c}^{2}}} \right)=\frac{1}{{{c}^{2}}}\frac{\partial }{\partial t}\left( \overrightarrow{a}\cdot \overrightarrow{R}-\overrightarrow{v}\cdot \overrightarrow{v} \right)=\frac{1}{{{c}^{2}}}\left( \overrightarrow{a}\frac{\partial \overrightarrow{R}}{\partial t}+\overrightarrow{R}\frac{\partial \overrightarrow{a}}{\partial t}-2\overrightarrow{v}\frac{\partial \overrightarrow{v}}{\partial t} \right)=\frac{1}{{{c}^{2}}}\left( \overrightarrow{a}\frac{\partial \overrightarrow{R}}{\partial t'}+\overrightarrow{R}\frac{\partial \overrightarrow{a}}{\partial t'}-2\overrightarrow{v}\frac{\partial \overrightarrow{v}}{\partial t'} \right)\frac{\partial t'}{\partial t} \\ 
 & =\frac{1}{{{c}^{2}}}\frac{R}{K}\left( -\,\overrightarrow{a}\cdot \overrightarrow{v}+\overrightarrow{R}\frac{\partial \overrightarrow{a}}{\partial t'}-2\overrightarrow{v}\cdot \overrightarrow{a} \right)=\frac{1}{{{c}^{2}}}\frac{R}{K}\left( \overrightarrow{R}\frac{\partial \overrightarrow{a}}{\partial t'}-3\ \overrightarrow{v}\cdot \overrightarrow{a} \right) \\ 
\end{align}

{]} И, следовательно, тут тоже исправлена ошибка!!!!! {[}

\begin{align}
  & \frac{1}{e}\frac{\partial \overrightarrow{{{E}_{1}}}}{\partial t}=\frac{\overrightarrow{R}}{{{K}^{4}}}\frac{R}{{{c}^{2}}}\left( \overrightarrow{R}\frac{\partial \overrightarrow{a}}{\partial t'}-3\ \overrightarrow{v}\cdot \overrightarrow{a} \right)+ \\ 
 & +\left( 1+\frac{\overrightarrow{a}\cdot \overrightarrow{R}}{{{c}^{2}}}-\frac{\overrightarrow{v}\cdot \overrightarrow{v}}{{{c}^{2}}} \right)\cdot \left\{ -\frac{\overrightarrow{v}R}{{{K}^{4}}}-\frac{3\overrightarrow{R}R}{{{K}^{5}}}\left( \frac{\overrightarrow{v}\cdot \overrightarrow{v}-\overrightarrow{R}\cdot \overrightarrow{a}}{c}-\frac{\overrightarrow{R}\cdot \overrightarrow{v}}{R} \right) \right\}+\frac{2R}{{{K}^{4}}}\frac{\overrightarrow{v}}{c}\left( \frac{\overrightarrow{v}\cdot \overrightarrow{v}-\overrightarrow{R}\cdot \overrightarrow{a}}{c}-\frac{\overrightarrow{R}\cdot \overrightarrow{v}}{R} \right)-\frac{R\overrightarrow{a}}{c{{K}^{3}}} \\ 
\end{align}

{]} {[}

\begin{align}
  & \frac{1}{e}\frac{\partial \overrightarrow{{{E}_{1}}}}{\partial t}=\frac{\overrightarrow{R}}{{{R}^{*}}^{3}}\frac{1}{{{c}^{2}}}\frac{R}{{{R}^{*}}}\left( \overrightarrow{R}\frac{\partial \overrightarrow{a}}{\partial t'}-3\ \overrightarrow{v}\cdot \overrightarrow{a} \right)-\left( 1+\frac{\overrightarrow{a}\cdot \overrightarrow{R}}{{{c}^{2}}}-\frac{\overrightarrow{v}\cdot \overrightarrow{v}}{{{c}^{2}}} \right)\cdot \left\{ \frac{\overrightarrow{v}R}{{{R}^{*}}^{4}}+\frac{3\overrightarrow{R}R}{{{R}^{*}}^{5}}\left( \frac{\overrightarrow{v}\cdot \overrightarrow{v}-\overrightarrow{R}\cdot \overrightarrow{a}}{c}-\frac{\overrightarrow{R}\cdot \overrightarrow{v}}{R} \right) \right\}+ \\ 
 & +\left\{ \frac{2}{{{R}^{*}}^{3}}\frac{\overrightarrow{v}}{c}\left( \frac{\overrightarrow{v}\cdot \overrightarrow{v}-\overrightarrow{R}\cdot \overrightarrow{a}}{c}-\frac{\overrightarrow{R}\cdot \overrightarrow{v}}{R} \right)-\frac{\overrightarrow{a}}{c{{R}^{*}}^{2}} \right\}\frac{R}{{{R}^{*}}} \\ 
\end{align}

{]} Далее выразим конвективную производную. При этом искомую скорость
(заряда или поля?) обозначим как \(\overrightarrow{W}\)
\(-\left( \overrightarrow{W},\nabla \right)\overrightarrow{E}=-{{W}_{x}}\frac{\partial \overrightarrow{E}}{\partial x}-{{W}_{y}}\frac{\partial \overrightarrow{E}}{\partial y}-{{W}_{z}}\frac{\partial \overrightarrow{E}}{\partial z}\)
Здесь найдем производную первой компоненты электрического поля (минус
градиент скалярного потенциала ЛВ) по координате \(x\) точки наблюдения
{[}

\begin{align}
  & \frac{1}{e}\frac{\partial \overrightarrow{{{E}_{1}}}}{\partial x}=\frac{\partial }{\partial x}\left[ \frac{\overrightarrow{R}}{{{K}^{3}}}\left( 1+\frac{\overrightarrow{a}\cdot \overrightarrow{R}}{{{c}^{2}}}-\frac{\overrightarrow{v}\cdot \overrightarrow{v}}{{{c}^{2}}} \right)-\frac{1}{{{K}^{2}}}\frac{\overrightarrow{v}}{c} \right]= \\ 
 & =\frac{\overrightarrow{R}}{{{K}^{3}}}\frac{\partial }{\partial x}\left( 1+\frac{\overrightarrow{a}\cdot \overrightarrow{R}}{{{c}^{2}}}-\frac{\overrightarrow{v}\cdot \overrightarrow{v}}{{{c}^{2}}} \right)+\left( 1+\frac{\overrightarrow{a}\cdot \overrightarrow{R}}{{{c}^{2}}}-\frac{\overrightarrow{v}\cdot \overrightarrow{v}}{{{c}^{2}}} \right)\cdot \frac{\partial }{\partial x}\left( \frac{\overrightarrow{R}}{{{K}^{3}}} \right)-\frac{\partial }{\partial x}\left( \frac{1}{{{K}^{2}}}\frac{\overrightarrow{v}}{c} \right) \\ 
\end{align}

{]}.

{[}\frac{\partial }{\partial x}\left(
\frac{\overrightarrow{R}}{{{K}^{3}}}
\right)=\frac{1}{{{K}^{3}}}\frac{\partial \overrightarrow{R}}{\partial x}+\overrightarrow{R}\frac{\partial }{\partial x}\left(
\frac{1}{{{K}^{3}}} \right)=\frac{1}{{{K}^{3}}}\left(
\overrightarrow{i}+\frac{\overrightarrow{v}}{c}\frac{{{R}_{x}}}{K}
\right)-\overrightarrow{R}\frac{3}{{{K}^{4}}}\frac{\partial K}{\partial x}{]}
{[}\frac{\partial K}{\partial x}=\frac{\partial }{\partial x}\left(
R-\frac{\overrightarrow{v}\cdot \overrightarrow{R}}{c}
\right)=\frac{{{R}_{x}}}{K}-\frac{1}{c}\frac{\partial }{\partial x}\left(
\overrightarrow{v}\cdot \overrightarrow{R} \right){]}
{[}\frac{\partial \overrightarrow{v}}{\partial x}=\frac{\partial \overrightarrow{v}}{\partial t'}\frac{\partial t'}{\partial x}=\overrightarrow{a}\left(
\frac{\partial }{\partial x}\left\{ t-\frac{R}{c} \right\}
\right)=\overrightarrow{a}\left(
-\frac{1}{c}\frac{\partial R}{\partial x}
\right)=-\frac{\overrightarrow{a}}{c}\frac{{{R}_{x}}}{K}{]}
{[}\frac{\partial }{\partial x}\left(
\overrightarrow{v}\cdot \overrightarrow{R}
\right)=\overrightarrow{v}\frac{\partial \overrightarrow{R}}{\partial x}+\overrightarrow{R}\frac{\partial \overrightarrow{v}}{\partial x}=\overrightarrow{v}\left(
\overrightarrow{i}+\frac{\overrightarrow{v}}{c}\frac{{{R}_{x}}}{K}
\right)+\overrightarrow{R}\left(
-\frac{\overrightarrow{a}}{c}\frac{{{R}_{x}}}{K}
\right)=\{\{v\}\emph{\{x\}\}+\left(
\frac{{{v}^{2}}-\overrightarrow{R}\cdot \overrightarrow{a}}{c}
\right)\frac{{{R}_{x}}}{K}{]}
{[}\frac{\partial K}{\partial x}=\frac{{{R}_{x}}}{K}-\frac{1}{c}\frac{\partial }{\partial x}\left(
\overrightarrow{v}\cdot \overrightarrow{R}
\right)=\frac{{{R}_{x}}}{K}-\frac{1}{c}\left( \{\{v\}}\{x\}\}+\left(
\frac{{{v}^{2}}-\overrightarrow{R}\cdot \overrightarrow{a}}{c}
\right)\frac{{{R}_{x}}}{K}
\right)=\frac{{{R}_{x}}}{K}-\frac{{{v}_{x}}}{c}+\frac{\overrightarrow{R}\cdot \overrightarrow{a}-{{v}^{2}}}{{{c}^{2}}}\frac{{{R}_{x}}}{K}=\frac{{{R}_{x}}}{K}\left(
1+\frac{\overrightarrow{R}\cdot \overrightarrow{a}-{{v}^{2}}}{{{c}^{2}}}
\right)-\frac{{{v}_{x}}}{c}{]} {[}\frac{\partial }{\partial x}\left(
\frac{\overrightarrow{R}}{{{K}^{3}}} \right)=\frac{1}{{{K}^{3}}}\left(
\overrightarrow{i}+\frac{\overrightarrow{v}}{c}\frac{{{R}_{x}}}{K}
\right)-\overrightarrow{R}\frac{3}{{{K}^{4}}}\left(
\frac{{{R}_{x}}}{K}+\left(
\frac{\overrightarrow{R}\cdot \overrightarrow{a}-{{v}^{2}}}{{{c}^{2}}}
\right)\frac{{{R}_{x}}}{K}-\frac{{{v}_{x}}}{c} \right){]}

{[}

\begin{align}
  & \frac{\partial }{\partial x}\left( \frac{1}{{{K}^{2}}}\frac{\overrightarrow{v}}{c} \right)=\frac{1}{{{K}^{2}}}\frac{\partial }{\partial x}\left( \frac{\overrightarrow{v}}{c} \right)+\frac{\overrightarrow{v}}{c}\frac{\partial }{\partial x}\left( \frac{1}{{{K}^{2}}} \right)=\frac{1}{c{{K}^{2}}}\left( -\frac{\overrightarrow{a}}{c}\frac{{{R}_{x}}}{K} \right)+\frac{\overrightarrow{v}}{c}\left( \frac{-2}{{{K}^{3}}} \right)\left\{ \frac{{{R}_{x}}}{K}\left( 1+\frac{\overrightarrow{R}\cdot \overrightarrow{a}-{{v}^{2}}}{{{c}^{2}}} \right)-\frac{{{v}_{x}}}{c} \right\}= \\ 
 & =-\frac{{{R}_{x}}}{{{K}^{3}}}\frac{\overrightarrow{a}}{{{c}^{2}}}-\frac{2}{{{K}^{3}}}\frac{\overrightarrow{v}}{c}\left\{ \frac{{{R}_{x}}}{K}\left( 1+\frac{\overrightarrow{R}\cdot \overrightarrow{a}-{{v}^{2}}}{{{c}^{2}}} \right)-\frac{{{v}_{x}}}{c} \right\} \\ 
\end{align}

{]}

{[}

\begin{align}
  & \frac{\partial }{\partial x}\left( 1+\frac{\overrightarrow{a}\cdot \overrightarrow{R}}{{{c}^{2}}}-\frac{\overrightarrow{v}\cdot \overrightarrow{v}}{{{c}^{2}}} \right)=\frac{1}{{{c}^{2}}}\frac{\partial }{\partial x}\left( \overrightarrow{a}\cdot \overrightarrow{R}-\overrightarrow{v}\cdot \overrightarrow{v} \right)=\frac{1}{{{c}^{2}}}\left\{ \overrightarrow{R}\frac{\partial \overrightarrow{a}}{\partial x}+\overrightarrow{a}\frac{\partial \overrightarrow{R}}{\partial x}-2\overrightarrow{v}\frac{\partial \overrightarrow{v}}{\partial x} \right\}= \\ 
 & =\frac{1}{{{c}^{2}}}\left\{ \overrightarrow{R}\frac{\partial \overrightarrow{a}}{\partial t'}\left( -\frac{1}{c}\frac{\partial R}{\partial x} \right)+\overrightarrow{a}\left( \overrightarrow{i}+\frac{\overrightarrow{v}}{c}\frac{{{R}_{x}}}{K} \right)+2\overrightarrow{v}\frac{\overrightarrow{a}}{c}\frac{{{R}_{x}}}{K} \right\}=\frac{1}{{{c}^{2}}}\left\{ -\frac{\overrightarrow{R}}{c}\frac{\partial \overrightarrow{a}}{\partial t'}\frac{{{R}_{x}}}{K}+{{a}_{x}}+3\overrightarrow{v}\frac{\overrightarrow{a}}{c}\frac{{{R}_{x}}}{K} \right\} \\ 
\end{align}

{]}

{[}

\begin{align}
  & \frac{1}{e}\frac{\partial \overrightarrow{{{E}_{1}}}}{\partial x}=\frac{\overrightarrow{R}}{{{K}^{3}}}\frac{\partial }{\partial x}\left( 1+\frac{\overrightarrow{a}\cdot \overrightarrow{R}}{{{c}^{2}}}-\frac{\overrightarrow{v}\cdot \overrightarrow{v}}{{{c}^{2}}} \right)+\left( 1+\frac{\overrightarrow{a}\cdot \overrightarrow{R}}{{{c}^{2}}}-\frac{\overrightarrow{v}\cdot \overrightarrow{v}}{{{c}^{2}}} \right)\cdot \frac{\partial }{\partial x}\left( \frac{\overrightarrow{R}}{{{K}^{3}}} \right)-\frac{\partial }{\partial x}\left( \frac{1}{{{K}^{2}}}\frac{\overrightarrow{v}}{c} \right)= \\ 
 &  \\ 
 & =\frac{\overrightarrow{R}}{{{K}^{3}}}\frac{1}{{{c}^{2}}}\left\{ -\frac{\overrightarrow{R}}{c}\frac{\partial \overrightarrow{a}}{\partial t'}\frac{{{R}_{x}}}{K}+{{a}_{x}}+3\overrightarrow{v}\frac{\overrightarrow{a}}{c}\frac{{{R}_{x}}}{K} \right\}+ \\ 
 & +\left( 1+\frac{\overrightarrow{a}\cdot \overrightarrow{R}}{{{c}^{2}}}-\frac{\overrightarrow{v}\cdot \overrightarrow{v}}{{{c}^{2}}} \right)\cdot \left\{ \frac{1}{{{K}^{3}}}\left( \overrightarrow{i}+\frac{\overrightarrow{v}}{c}\frac{{{R}_{x}}}{K} \right)-\frac{3\overrightarrow{R}}{{{K}^{4}}}\left( \frac{{{R}_{x}}}{K}\left( 1+\frac{\overrightarrow{R}\cdot \overrightarrow{a}-{{v}^{2}}}{{{c}^{2}}} \right)-\frac{{{v}_{x}}}{c} \right) \right\}+ \\ 
 & +\frac{{{R}_{x}}}{{{K}^{3}}}\frac{\overrightarrow{a}}{{{c}^{2}}}+\frac{2}{{{K}^{3}}}\frac{\overrightarrow{v}}{c}\left\{ \frac{{{R}_{x}}}{K}\left( 1+\frac{\overrightarrow{R}\cdot \overrightarrow{a}-{{v}^{2}}}{{{c}^{2}}} \right)-\frac{{{v}_{x}}}{c} \right\}= \\ 
 &  \\ 
 & =-\frac{\overrightarrow{R}}{{{K}^{4}}}\frac{{{R}_{x}}}{{{c}^{2}}}\left( \frac{\overrightarrow{R}}{c}\frac{\partial \overrightarrow{a}}{\partial t'} \right)+\frac{\overrightarrow{R}}{{{K}^{4}}}\frac{{{R}_{x}}}{{{c}^{2}}}\left( \frac{3\ \overrightarrow{v}\cdot \overrightarrow{a}}{c} \right)+\frac{\overrightarrow{R}}{{{K}^{3}}}\frac{{{a}_{x}}}{{{c}^{2}}}+ \\ 
 & +\left( 1+\frac{\overrightarrow{a}\cdot \overrightarrow{R}}{{{c}^{2}}}-\frac{\overrightarrow{v}\cdot \overrightarrow{v}}{{{c}^{2}}} \right)\cdot \left\{ \frac{\overrightarrow{i}}{{{K}^{3}}}+\frac{\overrightarrow{v}}{c}\frac{{{R}_{x}}}{{{K}^{4}}}-\frac{3\overrightarrow{R}}{K}\frac{{{R}_{x}}}{{{K}^{4}}}\left( 1+\frac{\overrightarrow{R}\cdot \overrightarrow{a}-{{v}^{2}}}{{{c}^{2}}} \right)+\frac{3\overrightarrow{R}}{{{K}^{4}}}\frac{{{v}_{x}}}{c} \right\}+ \\ 
 & +\frac{{{R}_{x}}}{{{K}^{3}}}\frac{\overrightarrow{a}}{{{c}^{2}}}+\frac{2{{R}_{x}}}{{{K}^{4}}}\frac{\overrightarrow{v}}{c}\left( 1+\frac{\overrightarrow{R}\cdot \overrightarrow{a}-{{v}^{2}}}{{{c}^{2}}} \right)-\frac{2}{{{K}^{3}}}\frac{\overrightarrow{v}}{c}\frac{{{v}_{x}}}{c}= \\ 
 &  \\ 
 & =-\frac{\overrightarrow{R}}{{{K}^{4}}}\frac{{{R}_{x}}}{{{c}^{2}}}\left( \frac{\overrightarrow{R}}{c}\frac{\partial \overrightarrow{a}}{\partial t'} \right)+\frac{\overrightarrow{R}}{{{K}^{4}}}\frac{{{R}_{x}}}{{{c}^{2}}}\left( \frac{3\ \overrightarrow{v}\cdot \overrightarrow{a}}{c} \right)+\frac{\overrightarrow{R}}{{{K}^{3}}}\frac{{{a}_{x}}}{{{c}^{2}}}+ \\ 
 & +\left( 1+\frac{\overrightarrow{a}\cdot \overrightarrow{R}}{{{c}^{2}}}-\frac{\overrightarrow{v}\cdot \overrightarrow{v}}{{{c}^{2}}} \right)\cdot \left\{ \frac{\overrightarrow{i}}{{{K}^{3}}}+\frac{3{{R}_{x}}}{{{K}^{4}}}\frac{\overrightarrow{v}}{c}-\frac{3\overrightarrow{R}}{K}\frac{{{R}_{x}}}{{{K}^{4}}}\left( 1+\frac{\overrightarrow{R}\cdot \overrightarrow{a}-{{v}^{2}}}{{{c}^{2}}} \right)+\frac{3\overrightarrow{R}}{{{K}^{4}}}\frac{{{v}_{x}}}{c} \right\}+ \\ 
 & +\frac{{{R}_{x}}}{{{K}^{3}}}\frac{\overrightarrow{a}}{{{c}^{2}}}-\frac{2}{{{K}^{3}}}\frac{\overrightarrow{v}}{c}\frac{{{v}_{x}}}{c}= \\ 
 &  \\ 
 & =-\frac{\overrightarrow{R}}{{{K}^{4}}}\frac{{{R}_{x}}}{{{c}^{2}}}\left( \frac{\overrightarrow{R}}{c}\frac{\partial \overrightarrow{a}}{\partial t'} \right)+\frac{\overrightarrow{R}}{{{K}^{4}}}\frac{{{R}_{x}}}{{{c}^{2}}}\left( \frac{3\ \overrightarrow{v}\cdot \overrightarrow{a}}{c} \right)+\frac{\overrightarrow{R}}{{{K}^{3}}}\frac{{{a}_{x}}}{{{c}^{2}}}+\frac{{{R}_{x}}}{{{K}^{3}}}\frac{\overrightarrow{a}}{{{c}^{2}}}-\frac{2}{{{K}^{3}}}\frac{\overrightarrow{v}}{c}\frac{{{v}_{x}}}{c}+ \\ 
 & +\left( 1+\frac{\overrightarrow{a}\cdot \overrightarrow{R}}{{{c}^{2}}}-\frac{\overrightarrow{v}\cdot \overrightarrow{v}}{{{c}^{2}}} \right)\cdot \left\{ \frac{\overrightarrow{i}}{{{K}^{3}}}+\frac{3{{R}_{x}}}{{{K}^{4}}}\frac{\overrightarrow{v}}{c}-\frac{3\overrightarrow{R}}{K}\frac{{{R}_{x}}}{{{K}^{4}}}\left( 1+\frac{\overrightarrow{R}\cdot \overrightarrow{a}-{{v}^{2}}}{{{c}^{2}}} \right)+\frac{3\overrightarrow{R}}{{{K}^{4}}}\frac{{{v}_{x}}}{c} \right\} \\ 
 &  \\ 
\end{align}

{]} {[}

\begin{align}
  & \frac{1}{e}\frac{\partial \overrightarrow{{{E}_{1}}}}{\partial t}=\frac{\overrightarrow{R}}{{{K}^{4}}}\frac{R}{{{c}^{2}}}\left( \overrightarrow{R}\frac{\partial \overrightarrow{a}}{\partial t'}-3\ \overrightarrow{v}\cdot \overrightarrow{a} \right)+ \\ 
 & +\left( 1+\frac{\overrightarrow{a}\cdot \overrightarrow{R}}{{{c}^{2}}}-\frac{\overrightarrow{v}\cdot \overrightarrow{v}}{{{c}^{2}}} \right)\cdot \left\{ -\frac{\overrightarrow{v}R}{{{K}^{4}}}-\frac{3\overrightarrow{R}R}{{{K}^{5}}}\left( \frac{\overrightarrow{v}\cdot \overrightarrow{v}-\overrightarrow{R}\cdot \overrightarrow{a}}{c}-\frac{\overrightarrow{R}\cdot \overrightarrow{v}}{R} \right) \right\}+\frac{2R}{{{K}^{4}}}\frac{\overrightarrow{v}}{c}\left( \frac{\overrightarrow{v}\cdot \overrightarrow{v}-\overrightarrow{R}\cdot \overrightarrow{a}}{c}-\frac{\overrightarrow{R}\cdot \overrightarrow{v}}{R} \right)-\frac{\overrightarrow{a}}{c{{K}^{2}}} \\ 
\end{align}

{]}

Четвертое уравнение Максвелла На форуме sciteclibrary.ru 04.06.18
появилась тема «Как запаздывающий Лиенар-Вихерт становится
"незапаздывающим". Визуализация»
http://www.sciteclibrary.ru/cgi-bin/yabb2/YaBB.pl?num=1528093569/0 Автор
темы, Дробышев, в частности, пишет «Забавно, но в процессе перехода
возникают локальные минимумы скалярного потенциала, например, здесь:

Изображение взято для момента \(t=7{,}5\), когда заряд уже движется
равномерно. Естественно, возникновение потенциальных ям не связано с
какими-то дополнительными зарядами, материализующимися "из воздуха". Это
ведь не статика, где \(\vec{E}=-\nabla \varphi\), а динамика с
\(\vec{E}=-\nabla \varphi -(1/c)\partial\vec{A}/ \partial t\), поэтому с
уравнением \(\nabla\cdot\vec{E}=0\) по-прежнему все в порядке.». Конец
цитаты. Приведенное наблюдение наталкивает на мысль проверить
непосредственным дифференцированием, равна ли нулю дивергенция
электрического поля в вакууме при ускоренном движении одиночного заряда.
Градиент скалярного потенциала ЛВ (первое слагаемое выражения
электрического поля через потенциалы)
{[}\overrightarrow{{{E}_{1}}}=-\nabla \varphi =-\nabla \frac{e}{K}=\frac{e}{{{K}^{2}}}\left\{
\frac{\overrightarrow{R}}{K}\left(
1+\frac{\overrightarrow{a}\cdot \overrightarrow{R}}{{{c}^{2}}}-\frac{\overrightarrow{v}\cdot \overrightarrow{v}}{{{c}^{2}}}
\right)-\frac{\overrightarrow{v}}{c} \right\}{]}
{[}\overrightarrow{{{E}_{1}}}=e\frac{\overrightarrow{R}}{{{K}^{3}}}\left(
1+\frac{\overrightarrow{a}\cdot \overrightarrow{R}}{{{c}^{2}}}-\frac{\overrightarrow{v}\cdot \overrightarrow{v}}{{{c}^{2}}}
\right)-\frac{e}{{{K}^{2}}}\frac{\overrightarrow{v}}{c}{]} Для расчёта
дивергенции выразим производную первого слагаемого электрического поля
по координате точки наблюдения {[}

\begin{align}
  & \frac{1}{e}\frac{\partial \overrightarrow{{{E}_{1}}}}{\partial x}=-\frac{\overrightarrow{R}}{{{K}^{4}}}\frac{{{R}_{x}}}{{{c}^{2}}}\left( \frac{\overrightarrow{R}}{c}\frac{\partial \overrightarrow{a}}{\partial t'} \right)+\frac{\overrightarrow{R}}{{{K}^{4}}}\frac{{{R}_{x}}}{{{c}^{2}}}\left( \frac{3\ \overrightarrow{v}\cdot \overrightarrow{a}}{c} \right)+\frac{\overrightarrow{R}}{{{K}^{3}}}\frac{{{a}_{x}}}{{{c}^{2}}}+\frac{{{R}_{x}}}{{{K}^{3}}}\frac{\overrightarrow{a}}{{{c}^{2}}}-\frac{2}{{{K}^{3}}}\frac{\overrightarrow{v}}{c}\frac{{{v}_{x}}}{c}+ \\ 
 & +\left( 1+\frac{\overrightarrow{a}\cdot \overrightarrow{R}}{{{c}^{2}}}-\frac{\overrightarrow{v}\cdot \overrightarrow{v}}{{{c}^{2}}} \right)\cdot \left\{ \frac{\overrightarrow{i}}{{{K}^{3}}}-\frac{3\overrightarrow{R}}{K}\frac{{{R}_{x}}}{{{K}^{4}}}\left( 1+\frac{\overrightarrow{R}\cdot \overrightarrow{a}-{{v}^{2}}}{{{c}^{2}}} \right)+\frac{3\overrightarrow{R}}{{{K}^{4}}}\frac{{{v}_{x}}}{c}+\frac{3{{R}_{x}}}{{{K}^{4}}}\frac{\overrightarrow{v}}{c} \right\} \\ 
\end{align}

{]} Выражения производной первого слагаемого электрического поля по
остальным координатам \(y\) и \(z\) точки наблюдения легко получить
меняя в данном выражении \({{R}_{x}}\) \({{a}_{x}}\) \({{v}_{x}}\)
\(\overrightarrow{i}\) на \({{R}_{y}}\) \({{a}_{y}}\) \({{v}_{y}}\)
\(\overrightarrow{j}\)и на \({{R}_{z}}\) \({{a}_{z}}\) \({{v}_{z}}\)
\(\overrightarrow{k}\)соответственно Возьмём теперь производную \(x\)
компоненты первого слагаемого электрического поля по \(x\) координате
точки наблюдения {[}

\begin{align}
  & \frac{1}{e}\frac{\partial {{E}_{{{1}_{x}}}}}{\partial x}=-\frac{{{R}_{x}}}{{{K}^{4}}}\frac{{{R}_{x}}}{{{c}^{2}}}\left( \frac{\overrightarrow{R}}{c}\frac{\partial \overrightarrow{a}}{\partial t'} \right)+\frac{{{R}_{x}}}{{{K}^{4}}}\frac{{{R}_{x}}}{{{c}^{2}}}\left( \frac{3\ \overrightarrow{v}\cdot \overrightarrow{a}}{c} \right)+\frac{{{R}_{x}}}{{{K}^{3}}}\frac{{{a}_{x}}}{{{c}^{2}}}+\frac{{{R}_{x}}}{{{K}^{3}}}\frac{{{a}_{x}}}{{{c}^{2}}}-\frac{2}{{{K}^{3}}}\frac{{{v}_{x}}}{c}\frac{{{v}_{x}}}{c}+ \\ 
 & +\left( 1+\frac{\overrightarrow{a}\cdot \overrightarrow{R}}{{{c}^{2}}}-\frac{\overrightarrow{v}\cdot \overrightarrow{v}}{{{c}^{2}}} \right)\cdot \left\{ \frac{1}{{{K}^{3}}}-\frac{3{{R}_{x}}}{K}\frac{{{R}_{x}}}{{{K}^{4}}}\left( 1+\frac{\overrightarrow{R}\cdot \overrightarrow{a}-{{v}^{2}}}{{{c}^{2}}} \right)+\frac{3{{R}_{x}}}{{{K}^{4}}}\frac{{{v}_{x}}}{c}+\frac{3{{R}_{x}}}{{{K}^{4}}}\frac{{{v}_{x}}}{c} \right\} \\ 
\end{align}

{]} Выражения производных \(y\) и \(z\) компонент первого слагаемого
электрического поля по \(y\)и \(z\) координатам точки наблюдения легко
получить аналогичной заменой индексов в данном выражении Скалярно
складывая производные\(x\) , \(y\) и \(z\)компонент первого слагаемого
электрического поля по соответственно \(x\) , \(y\) и \(z\)координатам
точки наблюдения получаем выражение для дивергенции первого слагаемого
электрического поля {[}

\begin{align}
  & \frac{1}{e}{{\left( \frac{\partial {{E}_{x}}}{\partial x}+\frac{\partial {{E}_{y}}}{\partial y}+\frac{\partial {{E}_{z}}}{\partial z} \right)}_{1}}=-\frac{{{R}^{2}}}{{{K}^{4}}{{c}^{2}}}\left( \frac{\overrightarrow{R}}{c}\frac{\partial \overrightarrow{a}}{\partial t'} \right)+\frac{{{R}^{2}}}{{{K}^{4}}{{c}^{2}}}\left( \frac{3\ \overrightarrow{v}\cdot \overrightarrow{a}}{c} \right)+\frac{2}{{{K}^{3}}}\frac{\overrightarrow{a}\cdot \overrightarrow{R}}{{{c}^{2}}}-\frac{2}{{{K}^{3}}}\frac{{{v}^{2}}}{{{c}^{2}}}+ \\ 
 & +\left( 1+\frac{\overrightarrow{a}\cdot \overrightarrow{R}}{{{c}^{2}}}-\frac{\overrightarrow{v}\cdot \overrightarrow{v}}{{{c}^{2}}} \right)\cdot \left\{ \frac{3}{{{K}^{3}}}-\frac{3{{R}^{2}}}{{{K}^{5}}}\left( 1+\frac{\overrightarrow{R}\cdot \overrightarrow{a}-{{v}^{2}}}{{{c}^{2}}} \right)+6\frac{\overrightarrow{v}\cdot \overrightarrow{R}}{c{{K}^{4}}} \right\} \\ 
\end{align}

{]}

{[}

\begin{align}
  & \frac{1}{e}{{\left( \frac{\partial {{E}_{x}}}{\partial x}+\frac{\partial {{E}_{y}}}{\partial y}+\frac{\partial {{E}_{z}}}{\partial z} \right)}_{1}}=-\frac{{{R}^{2}}}{{{K}^{4}}{{c}^{2}}}\left( \frac{\overrightarrow{R}}{c}\frac{\partial \overrightarrow{a}}{\partial t'} \right)+\frac{{{R}^{2}}}{{{K}^{4}}{{c}^{2}}}\left( \frac{3\ \overrightarrow{v}\cdot \overrightarrow{a}}{c} \right)+\frac{2}{{{K}^{3}}}\left( \frac{\overrightarrow{a}\cdot \overrightarrow{R}}{{{c}^{2}}}-\frac{{{v}^{2}}}{{{c}^{2}}} \right)+ \\ 
 & +\left( 1+\frac{\overrightarrow{a}\cdot \overrightarrow{R}}{{{c}^{2}}}-\frac{\overrightarrow{v}\cdot \overrightarrow{v}}{{{c}^{2}}} \right)\cdot \left\{ \frac{3}{{{K}^{3}}}-\frac{3{{R}^{2}}}{{{K}^{5}}}\left( 1+\frac{\overrightarrow{R}\cdot \overrightarrow{a}-{{v}^{2}}}{{{c}^{2}}} \right)+6\frac{\overrightarrow{v}\cdot \overrightarrow{R}}{c{{K}^{4}}} \right\} \\ 
\end{align}

{]}

{[}

\begin{align}
  & \frac{1}{e}{{\left( \frac{\partial {{E}_{x}}}{\partial x}+\frac{\partial {{E}_{y}}}{\partial y}+\frac{\partial {{E}_{z}}}{\partial z} \right)}_{1}}=-\frac{{{R}^{2}}}{{{K}^{4}}{{c}^{2}}}\left( \frac{\overrightarrow{R}}{c}\frac{\partial \overrightarrow{a}}{\partial t'} \right)+\frac{{{R}^{2}}}{{{K}^{4}}{{c}^{2}}}\left( \frac{3\ \overrightarrow{v}\cdot \overrightarrow{a}}{c} \right)+\frac{2}{{{K}^{3}}}\left( 1+\frac{\overrightarrow{a}\cdot \overrightarrow{R}}{{{c}^{2}}}-\frac{{{v}^{2}}}{{{c}^{2}}} \right)-\frac{2}{{{K}^{3}}}+ \\ 
 & +\left( 1+\frac{\overrightarrow{a}\cdot \overrightarrow{R}}{{{c}^{2}}}-\frac{\overrightarrow{v}\cdot \overrightarrow{v}}{{{c}^{2}}} \right)\cdot \left\{ \frac{3}{{{K}^{3}}}-\frac{3{{R}^{2}}}{{{K}^{5}}}\left( 1+\frac{\overrightarrow{R}\cdot \overrightarrow{a}-{{v}^{2}}}{{{c}^{2}}} \right)+6\frac{\overrightarrow{v}\cdot \overrightarrow{R}}{c{{K}^{4}}} \right\} \\ 
\end{align}

{]}

{[}

\begin{align}
  & \frac{1}{e}{{\left( \frac{\partial {{E}_{x}}}{\partial x}+\frac{\partial {{E}_{y}}}{\partial y}+\frac{\partial {{E}_{z}}}{\partial z} \right)}_{1}}=-\frac{{{R}^{2}}}{{{K}^{4}}{{c}^{2}}}\left( \frac{\overrightarrow{R}}{c}\frac{\partial \overrightarrow{a}}{\partial t'} \right)+\frac{{{R}^{2}}}{{{K}^{4}}{{c}^{2}}}\left( \frac{3\ \overrightarrow{v}\cdot \overrightarrow{a}}{c} \right)-\frac{2}{{{K}^{3}}}+ \\ 
 & +\left( 1+\frac{\overrightarrow{a}\cdot \overrightarrow{R}}{{{c}^{2}}}-\frac{\overrightarrow{v}\cdot \overrightarrow{v}}{{{c}^{2}}} \right)\cdot \left\{ \frac{5}{{{K}^{3}}}-\frac{3{{R}^{2}}}{{{K}^{5}}}\left( 1+\frac{\overrightarrow{R}\cdot \overrightarrow{a}-{{v}^{2}}}{{{c}^{2}}} \right)+6\frac{\overrightarrow{v}\cdot \overrightarrow{R}}{c{{K}^{4}}} \right\} \\ 
\end{align}

{]}

Второе слагаемое выражения электрического поля через потенциалы
http://www.sciteclibrary.ru/cgi-bin/yabb2/YaBB.pl?num=1528093569/417\#417
{[}\overrightarrow{{{E}_{2}}}=-\frac{1}{c}\frac{\partial \overrightarrow{A}}{\partial t}=\frac{e}{{{K}^{2}}}\left[ \frac{\overrightarrow{v}}{c}\left\{ \frac{R}{K}\left( \frac{{{v}^{2}}}{{{c}^{2}}}-\frac{\overrightarrow{a}\cdot \overrightarrow{R}}{{{c}^{2}}}-1 \right)+1 \right\}-R\frac{\overrightarrow{a}}{{{c}^{2}}} \right]{]}
{[}\overrightarrow{{{E}_{2}}}=\frac{e}{{{K}^{2}}}\frac{\overrightarrow{v}}{c}\left\{
\frac{R}{K}\left(
\frac{{{v}^{2}}}{{{c}^{2}}}-\frac{\overrightarrow{a}\cdot \overrightarrow{R}}{{{c}^{2}}}-1
\right)+1
\right\}-\frac{e}{{{K}^{2}}}R\frac{\overrightarrow{a}}{{{c}^{2}}}{]}
{[}\overrightarrow{{{E}_{2}}}=\frac{eR}{{{K}^{3}}}\frac{\overrightarrow{v}}{c}\left(
\frac{{{v}^{2}}}{{{c}^{2}}}-\frac{\overrightarrow{a}\cdot \overrightarrow{R}}{{{c}^{2}}}-1
\right)+\frac{e}{{{K}^{2}}}\frac{\overrightarrow{v}}{c}-\frac{eR}{{{K}^{2}}}\frac{\overrightarrow{a}}{{{c}^{2}}}{]}
{[}\frac{1}{e}\frac{\partial \overrightarrow{{{E}_{2}}}}{\partial x}=\frac{\partial }{\partial x}\left[ \frac{R}{{{K}^{3}}}\frac{\overrightarrow{v}}{c}\left( \frac{{{v}^{2}}}{{{c}^{2}}}-\frac{\overrightarrow{a}\cdot \overrightarrow{R}}{{{c}^{2}}}-1 \right)+\frac{1}{{{K}^{2}}}\frac{\overrightarrow{v}}{c}-\frac{R}{{{K}^{2}}}\frac{\overrightarrow{a}}{{{c}^{2}}} \right]{]}
{[}\frac{1}{e}\frac{\partial \overrightarrow{{{E}_{2}}}}{\partial x}=\left(
\frac{R}{{{K}^{3}}}\frac{\overrightarrow{v}}{c}
\right)\frac{\partial }{\partial x}\left(
\frac{{{v}^{2}}}{{{c}^{2}}}-\frac{\overrightarrow{a}\cdot \overrightarrow{R}}{{{c}^{2}}}-1
\right)+\left(
\frac{{{v}^{2}}}{{{c}^{2}}}-\frac{\overrightarrow{a}\cdot \overrightarrow{R}}{{{c}^{2}}}-1
\right)\frac{\partial }{\partial x}\left(
\frac{R}{{{K}^{3}}}\frac{\overrightarrow{v}}{c}
\right)+\frac{\partial }{\partial x}\left(
\frac{1}{{{K}^{2}}}\frac{\overrightarrow{v}}{c}
\right)-\frac{\partial }{\partial x}\left(
\frac{R}{{{K}^{2}}}\frac{\overrightarrow{a}}{{{c}^{2}}} \right){]}
{[}\frac{\partial \overrightarrow{v}}{\partial x}=\frac{\partial \overrightarrow{v}}{\partial t'}\frac{\partial t'}{\partial x}=\overrightarrow{a}\left(
-\frac{1}{c}\frac{\partial R}{\partial x}
\right)=-\frac{\overrightarrow{a}}{c}\frac{{{R}_{x}}}{K}{]}
{[}\frac{\partial K}{\partial x}=\frac{{{R}_{x}}}{K}\left(
1+\frac{\overrightarrow{R}\cdot \overrightarrow{a}-{{v}^{2}}}{{{c}^{2}}}
\right)-\frac{{{v}_{x}}}{c}{]} {[}\frac{\partial }{\partial x}\left(
\frac{1}{{{K}^{2}}}\frac{\overrightarrow{v}}{c}
\right)=-\frac{{{R}_{x}}}{{{K}^{3}}}\frac{\overrightarrow{a}}{{{c}^{2}}}-\frac{2}{{{K}^{3}}}\frac{\overrightarrow{v}}{c}\left\{
\frac{{{R}_{x}}}{K}\left(
1+\frac{\overrightarrow{R}\cdot \overrightarrow{a}-{{v}^{2}}}{{{c}^{2}}}
\right)-\frac{{{v}_{x}}}{c} \right\}{]}
{[}\frac{\partial }{\partial x}\left(
1+\frac{\overrightarrow{a}\cdot \overrightarrow{R}}{{{c}^{2}}}-\frac{\overrightarrow{v}\cdot \overrightarrow{v}}{{{c}^{2}}}
\right)=\frac{1}{{{c}^{2}}}\left\{
-\frac{\overrightarrow{R}}{c}\frac{\partial \overrightarrow{a}}{\partial t'}\frac{{{R}_{x}}}{K}+\{\{a\}\_\{x\}\}+3\overrightarrow{v}\frac{\overrightarrow{a}}{c}\frac{{{R}_{x}}}{K}
\right\}{]} {[}\frac{\partial R}{\partial x}=\frac{{{R}_{x}}}{K}{]}
{[}\frac{\partial }{\partial x}\left(
\frac{R}{{{K}^{2}}}\frac{\overrightarrow{a}}{{{c}^{2}}}
\right)=\frac{\overrightarrow{a}}{{{c}^{2}}}\frac{\partial }{\partial x}\left(
\frac{R}{{{K}^{2}}}
\right)+\frac{R}{{{c}^{2}}{{K}^{2}}}\frac{\partial \overrightarrow{a}}{\partial x}=\frac{\overrightarrow{a}}{{{c}^{2}}}\frac{\partial }{\partial x}\left(
\frac{R}{{{K}^{2}}}
\right)+\frac{R}{{{c}^{2}}{{K}^{2}}}\frac{\partial \overrightarrow{a}}{\partial t'}\left(
-\frac{1}{c}\frac{\partial R}{\partial x} \right){]}
{[}\frac{\partial }{\partial x}\left(
\frac{R}{{{K}^{2}}}\frac{\overrightarrow{a}}{{{c}^{2}}}
\right)=\frac{\overrightarrow{a}}{{{c}^{2}}}\frac{\partial }{\partial x}\left(
\frac{R}{{{K}^{2}}}
\right)-\frac{R}{{{c}^{3}}{{K}^{2}}}\frac{\partial \overrightarrow{a}}{\partial t'}\frac{{{R}_{x}}}{K}{]}
{[}\frac{\partial }{\partial x}\left( \frac{R}{{{K}^{2}}}
\right)=\frac{1}{{{K}^{2}}}\frac{\partial R}{\partial x}+R\frac{\partial }{\partial x}\left(
\frac{1}{{{K}^{2}}}
\right)=\frac{1}{{{K}^{2}}}\frac{{{R}_{x}}}{K}-\frac{2R}{{{K}^{3}}}\left\{
\frac{{{R}_{x}}}{K}\left(
1+\frac{\overrightarrow{R}\cdot \overrightarrow{a}-{{v}^{2}}}{{{c}^{2}}}
\right)-\frac{{{v}_{x}}}{c} \right\}{]}
{[}\frac{\partial }{\partial x}\left(
\frac{R}{{{K}^{2}}}\frac{\overrightarrow{a}}{{{c}^{2}}}
\right)=\frac{\overrightarrow{a}}{{{c}^{2}}}\left[ \frac{1}{{{K}^{2}}}\frac{{{R}_{x}}}{K}-\frac{2R}{{{K}^{3}}}\left\{ \frac{{{R}_{x}}}{K}\left( 1+\frac{\overrightarrow{R}\cdot \overrightarrow{a}-{{v}^{2}}}{{{c}^{2}}} \right)-\frac{{{v}_{x}}}{c} \right\} \right]-\frac{R}{{{c}^{3}}{{K}^{2}}}\frac{\partial \overrightarrow{a}}{\partial t'}\frac{{{R}_{x}}}{K}{]}
{[}\frac{\partial }{\partial x}\left(
\frac{R}{{{K}^{2}}}\frac{\overrightarrow{a}}{{{c}^{2}}}
\right)=\frac{1}{{{K}^{3}}{{c}^{2}}}\left(
\overrightarrow{a}\left[ {{R}_{x}}-2R\left\{ \frac{{{R}_{x}}}{K}\left( 1+\frac{\overrightarrow{R}\cdot \overrightarrow{a}-{{v}^{2}}}{{{c}^{2}}} \right)-\frac{{{v}_{x}}}{c} \right\} \right]-\frac{R{{R}_{x}}}{c}\frac{\partial \overrightarrow{a}}{\partial t'}
\right){]}
\(\frac{}{{{}^{2}}}-\frac{}{/}\frac{}{{{}^{3}}}=\frac{{{}^{2}}}{{{}^{2}}}\)
{[}\frac{\partial }{\partial x}\left( \frac{R}{{{K}^{3}}}
\right)=\frac{1}{{{K}^{3}}}\frac{\partial R}{\partial x}+R\frac{\partial }{\partial x}\left(
\frac{1}{{{K}^{3}}}
\right)=\frac{1}{{{K}^{3}}}\frac{{{R}_{x}}}{K}-\frac{3R}{{{K}^{4}}}\left\{
\frac{{{R}_{x}}}{K}\left(
1+\frac{\overrightarrow{R}\cdot \overrightarrow{a}-{{v}^{2}}}{{{c}^{2}}}
\right)-\frac{{{v}_{x}}}{c} \right\}{]}
{[}\frac{\partial \overrightarrow{v}}{\partial x}=\frac{\partial \overrightarrow{v}}{\partial t'}\frac{\partial t'}{\partial x}=\overrightarrow{a}\left(
-\frac{1}{c}\frac{\partial R}{\partial x}
\right)=-\frac{\overrightarrow{a}}{c}\frac{{{R}_{x}}}{K}{]}
{[}\frac{\partial }{\partial x}\left(
\frac{R}{{{K}^{3}}}\frac{\overrightarrow{v}}{c}
\right)=\frac{\overrightarrow{v}}{c}\frac{\partial }{\partial x}\left(
\frac{R}{{{K}^{3}}}
\right)+\frac{R}{c{{K}^{3}}}\frac{\partial \overrightarrow{v}}{\partial t'}\frac{\partial t'}{\partial x}=\frac{\overrightarrow{v}}{c}\frac{\partial }{\partial x}\left(
\frac{R}{{{K}^{3}}}
\right)-\frac{R}{c{{K}^{3}}}\frac{\overrightarrow{a}}{c}\frac{{{R}_{x}}}{K}{]}
{[}\frac{\partial }{\partial x}\left(
\frac{R}{{{K}^{3}}}\frac{\overrightarrow{v}}{c}
\right)=\frac{\overrightarrow{v}}{c}\left[ \frac{{{R}_{x}}}{{{K}^{4}}}-\frac{3R}{{{K}^{4}}}\left\{ \frac{{{R}_{x}}}{K}\left( 1+\frac{\overrightarrow{R}\cdot \overrightarrow{a}-{{v}^{2}}}{{{c}^{2}}} \right)-\frac{{{v}_{x}}}{c} \right\} \right]-\frac{R{{R}_{x}}}{c{{K}^{4}}}\frac{\overrightarrow{a}}{c}{]}
{[}\frac{\partial }{\partial x}\left(
\frac{R}{{{K}^{3}}}\frac{\overrightarrow{v}}{c}
\right)=\frac{1}{{{K}^{4}}c}\left(
\overrightarrow{v}\left[ {{R}_{x}}-3R\left\{ \frac{{{R}_{x}}}{K}\left( 1+\frac{\overrightarrow{R}\cdot \overrightarrow{a}-{{v}^{2}}}{{{c}^{2}}} \right)-\frac{{{v}_{x}}}{c} \right\} \right]-\frac{R{{R}_{x}}}{c}\overrightarrow{a}
\right){]}
{[}\frac{1}{e}\frac{\partial \overrightarrow{{{E}_{2}}}}{\partial x}=\left(
\frac{R}{{{K}^{3}}}\frac{\overrightarrow{v}}{c}
\right)\frac{\partial }{\partial x}\left(
\frac{{{v}^{2}}}{{{c}^{2}}}-\frac{\overrightarrow{a}\cdot \overrightarrow{R}}{{{c}^{2}}}-1
\right)+\left(
\frac{{{v}^{2}}}{{{c}^{2}}}-\frac{\overrightarrow{a}\cdot \overrightarrow{R}}{{{c}^{2}}}-1
\right)\frac{\partial }{\partial x}\left(
\frac{R}{{{K}^{3}}}\frac{\overrightarrow{v}}{c}
\right)+\frac{\partial }{\partial x}\left(
\frac{1}{{{K}^{2}}}\frac{\overrightarrow{v}}{c}
\right)-\frac{\partial }{\partial x}\left(
\frac{R}{{{K}^{2}}}\frac{\overrightarrow{a}}{{{c}^{2}}} \right){]}

{[}

\begin{align}
  & \frac{1}{e}\frac{\partial \overrightarrow{{{E}_{2}}}}{\partial x}=\left( \frac{R}{{{K}^{3}}}\frac{\overrightarrow{v}}{c} \right)\frac{1}{{{c}^{2}}}\left\{ \frac{\overrightarrow{R}}{c}\frac{\partial \overrightarrow{a}}{\partial t'}\frac{{{R}_{x}}}{K}-{{a}_{x}}-3\overrightarrow{v}\frac{\overrightarrow{a}}{c}\frac{{{R}_{x}}}{K} \right\}+ \\ 
 & +\left( \frac{{{v}^{2}}}{{{c}^{2}}}-\frac{\overrightarrow{a}\cdot \overrightarrow{R}}{{{c}^{2}}}-1 \right)\frac{1}{{{K}^{4}}c}\left( \overrightarrow{v}\left[ {{R}_{x}}-3R\left\{ \frac{{{R}_{x}}}{K}\left( 1+\frac{\overrightarrow{R}\cdot \overrightarrow{a}-{{v}^{2}}}{{{c}^{2}}} \right)-\frac{{{v}_{x}}}{c} \right\} \right]-\frac{R{{R}_{x}}}{c}\overrightarrow{a} \right)- \\ 
 & -\frac{{{R}_{x}}}{{{K}^{3}}}\frac{\overrightarrow{a}}{{{c}^{2}}}-\frac{2}{{{K}^{3}}}\frac{\overrightarrow{v}}{c}\left\{ \frac{{{R}_{x}}}{K}\left( 1+\frac{\overrightarrow{R}\cdot \overrightarrow{a}-{{v}^{2}}}{{{c}^{2}}} \right)-\frac{{{v}_{x}}}{c} \right\}- \\ 
 & -\frac{1}{{{K}^{3}}{{c}^{2}}}\left( \overrightarrow{a}\left[ {{R}_{x}}-2R\left\{ \frac{{{R}_{x}}}{K}\left( 1+\frac{\overrightarrow{R}\cdot \overrightarrow{a}-{{v}^{2}}}{{{c}^{2}}} \right)-\frac{{{v}_{x}}}{c} \right\} \right]-\frac{R{{R}_{x}}}{c}\frac{\partial \overrightarrow{a}}{\partial t'} \right) \\ 
\end{align}

{]}

{[}

\begin{align}
  & \frac{1}{e}\frac{\partial \overrightarrow{{{E}_{2}}}}{\partial x}=\frac{R}{{{K}^{3}}}\frac{1}{{{c}^{2}}}\left\{ \left( \frac{\overrightarrow{R}}{c}\frac{\partial \overrightarrow{a}}{\partial t'} \right)\frac{{{R}_{x}}}{K}\frac{\overrightarrow{v}}{c}-{{a}_{x}}\frac{\overrightarrow{v}}{c}-3\left( \overrightarrow{v}\frac{\overrightarrow{a}}{c} \right)\frac{{{R}_{x}}}{K}\frac{\overrightarrow{v}}{c} \right\}- \\ 
 & -\left( 1+\frac{\overrightarrow{R}\cdot \overrightarrow{a}-{{v}^{2}}}{{{c}^{2}}} \right)\frac{1}{{{K}^{4}}}\left( {{R}_{x}}\frac{\overrightarrow{v}}{c}-3R\left\{ \frac{{{R}_{x}}}{K}\frac{\overrightarrow{v}}{c}\left( 1+\frac{\overrightarrow{R}\cdot \overrightarrow{a}-{{v}^{2}}}{{{c}^{2}}} \right)-\frac{{{v}_{x}}}{c}\frac{\overrightarrow{v}}{c} \right\}-\frac{R{{R}_{x}}}{{{c}^{2}}}\overrightarrow{a} \right)- \\ 
 & -\frac{{{R}_{x}}}{{{K}^{3}}}\frac{\overrightarrow{a}}{{{c}^{2}}}-\frac{2}{{{K}^{3}}}\left\{ \frac{{{R}_{x}}}{K}\frac{\overrightarrow{v}}{c}\left( 1+\frac{\overrightarrow{R}\cdot \overrightarrow{a}-{{v}^{2}}}{{{c}^{2}}} \right)-\frac{{{v}_{x}}}{c}\frac{\overrightarrow{v}}{c} \right\}- \\ 
 & -\frac{1}{{{K}^{3}}{{c}^{2}}}\left( {{R}_{x}}\overrightarrow{a}-2R\left\{ \frac{{{R}_{x}}}{K}\overrightarrow{a}\left( 1+\frac{\overrightarrow{R}\cdot \overrightarrow{a}-{{v}^{2}}}{{{c}^{2}}} \right)-\frac{{{v}_{x}}}{c}\overrightarrow{a} \right\}-\frac{R{{R}_{x}}}{c}\frac{\partial \overrightarrow{a}}{\partial t'} \right) \\ 
\end{align}

{]}

{[}

\begin{align}
  & \frac{1}{e}\frac{\partial \overrightarrow{{{E}_{2}}}}{\partial x}=\frac{R}{{{K}^{3}}}\frac{1}{{{c}^{2}}}\left\{ \left( \frac{\overrightarrow{R}}{c}\frac{\partial \overrightarrow{a}}{\partial t'} \right)\frac{{{R}_{x}}}{K}\frac{\overrightarrow{v}}{c}-{{a}_{x}}\frac{\overrightarrow{v}}{c}-3\left( \overrightarrow{v}\frac{\overrightarrow{a}}{c} \right)\frac{{{R}_{x}}}{K}\frac{\overrightarrow{v}}{c} \right\}- \\ 
 & -\left( 1+\frac{\overrightarrow{R}\cdot \overrightarrow{a}-{{v}^{2}}}{{{c}^{2}}} \right)\frac{1}{{{K}^{4}}}\left( {{R}_{x}}\frac{\overrightarrow{v}}{c}-3R\left\{ \frac{{{R}_{x}}}{K}\frac{\overrightarrow{v}}{c}\left( 1+\frac{\overrightarrow{R}\cdot \overrightarrow{a}-{{v}^{2}}}{{{c}^{2}}} \right)-\frac{{{v}_{x}}}{c}\frac{\overrightarrow{v}}{c} \right\}-\frac{R{{R}_{x}}}{{{c}^{2}}}\overrightarrow{a} \right)- \\ 
 & -\frac{{{R}_{x}}}{{{K}^{3}}}\frac{\overrightarrow{a}}{{{c}^{2}}}-\frac{2}{{{K}^{3}}}\left\{ \frac{{{R}_{x}}}{K}\frac{\overrightarrow{v}}{c}\left( 1+\frac{\overrightarrow{R}\cdot \overrightarrow{a}-{{v}^{2}}}{{{c}^{2}}} \right)-\frac{{{v}_{x}}}{c}\frac{\overrightarrow{v}}{c} \right\}- \\ 
 & -\frac{{{R}_{x}}\overrightarrow{a}}{{{K}^{3}}{{c}^{2}}}+\frac{R}{{{K}^{3}}{{c}^{2}}}\left\{ 2\frac{{{R}_{x}}}{K}\overrightarrow{a}\left( 1+\frac{\overrightarrow{R}\cdot \overrightarrow{a}-{{v}^{2}}}{{{c}^{2}}} \right)-2\frac{{{v}_{x}}}{c}\overrightarrow{a} \right\}+\frac{R}{{{K}^{3}}{{c}^{2}}}\frac{{{R}_{x}}}{c}\frac{\partial \overrightarrow{a}}{\partial t'} \\ 
\end{align}

{]}

{[}

\begin{align}
  & \frac{1}{e}\frac{\partial \overrightarrow{{{E}_{2}}}}{\partial x}=\frac{R}{{{K}^{3}}}\frac{1}{{{c}^{2}}}\left\{ \left( \frac{\overrightarrow{R}}{c}\frac{\partial \overrightarrow{a}}{\partial t'} \right)\frac{{{R}_{x}}}{K}\frac{\overrightarrow{v}}{c}+\frac{{{R}_{x}}}{c}\frac{\partial \overrightarrow{a}}{\partial t'}-{{a}_{x}}\frac{\overrightarrow{v}}{c}-2\frac{{{v}_{x}}}{c}\overrightarrow{a}-3\left( \overrightarrow{v}\frac{\overrightarrow{a}}{c} \right)\frac{{{R}_{x}}}{K}\frac{\overrightarrow{v}}{c} \right\}- \\ 
 & -\left( 1+\frac{\overrightarrow{R}\cdot \overrightarrow{a}-{{v}^{2}}}{{{c}^{2}}} \right)\frac{1}{{{K}^{4}}}\left( {{R}_{x}}\frac{\overrightarrow{v}}{c}-3R\left\{ \frac{{{R}_{x}}}{K}\frac{\overrightarrow{v}}{c}\left( 1+\frac{\overrightarrow{R}\cdot \overrightarrow{a}-{{v}^{2}}}{{{c}^{2}}} \right)-\frac{{{v}_{x}}}{c}\frac{\overrightarrow{v}}{c} \right\}-\frac{R{{R}_{x}}}{{{c}^{2}}}\overrightarrow{a} \right)- \\ 
 & -2\frac{{{R}_{x}}}{{{K}^{3}}}\frac{\overrightarrow{a}}{{{c}^{2}}}-\frac{2}{{{K}^{3}}}\left\{ \frac{{{R}_{x}}}{K}\frac{\overrightarrow{v}}{c}\left( 1+\frac{\overrightarrow{R}\cdot \overrightarrow{a}-{{v}^{2}}}{{{c}^{2}}} \right)-\frac{{{v}_{x}}}{c}\frac{\overrightarrow{v}}{c} \right\}+\frac{2R}{{{K}^{3}}{{c}^{2}}}\frac{{{R}_{x}}}{K}\overrightarrow{a}\left( 1+\frac{\overrightarrow{R}\cdot \overrightarrow{a}-{{v}^{2}}}{{{c}^{2}}} \right) \\ 
\end{align}

{]}

{[}

\begin{align}
  & \frac{1}{e}\frac{\partial \overrightarrow{{{E}_{2}}}}{\partial x}=\frac{R}{{{K}^{3}}}\frac{1}{{{c}^{2}}}\left\{ \left( \frac{\overrightarrow{R}}{c}\frac{\partial \overrightarrow{a}}{\partial t'} \right)\frac{{{R}_{x}}}{K}\frac{\overrightarrow{v}}{c}+\frac{{{R}_{x}}}{c}\frac{\partial \overrightarrow{a}}{\partial t'}-{{a}_{x}}\frac{\overrightarrow{v}}{c}-2\frac{{{v}_{x}}}{c}\overrightarrow{a}-3\left( \overrightarrow{v}\frac{\overrightarrow{a}}{c} \right)\frac{{{R}_{x}}}{K}\frac{\overrightarrow{v}}{c} \right\}- \\ 
 & -\left( 1+\frac{\overrightarrow{R}\cdot \overrightarrow{a}-{{v}^{2}}}{{{c}^{2}}} \right)\left( \frac{{{R}_{x}}}{{{K}^{4}}}\frac{\overrightarrow{v}}{c}-\frac{3R}{{{K}^{4}}}\left\{ \frac{{{R}_{x}}}{K}\frac{\overrightarrow{v}}{c}\left( 1+\frac{\overrightarrow{R}\cdot \overrightarrow{a}-{{v}^{2}}}{{{c}^{2}}} \right)-\frac{{{v}_{x}}}{c}\frac{\overrightarrow{v}}{c} \right\}-\frac{R{{R}_{x}}}{{{K}^{4}}{{c}^{2}}}\overrightarrow{a}-\frac{2R}{{{K}^{3}}{{c}^{2}}}\frac{{{R}_{x}}}{K}\overrightarrow{a} \right)- \\ 
 & -2\frac{{{R}_{x}}}{{{K}^{3}}}\frac{\overrightarrow{a}}{{{c}^{2}}}-2\frac{{{R}_{x}}}{{{K}^{4}}}\frac{\overrightarrow{v}}{c}\left( 1+\frac{\overrightarrow{R}\cdot \overrightarrow{a}-{{v}^{2}}}{{{c}^{2}}} \right)+\frac{2}{{{K}^{3}}}\frac{{{v}_{x}}}{c}\frac{\overrightarrow{v}}{c} \\ 
\end{align}

{]}

{[}

\begin{align}
  & \frac{1}{e}\frac{\partial \overrightarrow{{{E}_{2}}}}{\partial x}=\frac{R}{{{K}^{3}}}\frac{1}{{{c}^{2}}}\left\{ \left( \frac{\overrightarrow{R}}{c}\frac{\partial \overrightarrow{a}}{\partial t'} \right)\frac{{{R}_{x}}}{K}\frac{\overrightarrow{v}}{c}+\frac{{{R}_{x}}}{c}\frac{\partial \overrightarrow{a}}{\partial t'}-{{a}_{x}}\frac{\overrightarrow{v}}{c}-2\frac{{{v}_{x}}}{c}\overrightarrow{a}-3\left( \overrightarrow{v}\frac{\overrightarrow{a}}{c} \right)\frac{{{R}_{x}}}{K}\frac{\overrightarrow{v}}{c} \right\}- \\ 
 & -\left( 1+\frac{\overrightarrow{R}\cdot \overrightarrow{a}-{{v}^{2}}}{{{c}^{2}}} \right)\left( 3\frac{{{R}_{x}}}{{{K}^{4}}}\frac{\overrightarrow{v}}{c}-\frac{3R}{{{K}^{4}}}\left\{ \frac{{{R}_{x}}}{K}\frac{\overrightarrow{v}}{c}\left( 1+\frac{\overrightarrow{R}\cdot \overrightarrow{a}-{{v}^{2}}}{{{c}^{2}}} \right)-\frac{{{v}_{x}}}{c}\frac{\overrightarrow{v}}{c} \right\}-\frac{3R}{{{K}^{3}}{{c}^{2}}}\frac{{{R}_{x}}}{K}\overrightarrow{a} \right)-2\frac{{{R}_{x}}}{{{K}^{3}}}\frac{\overrightarrow{a}}{{{c}^{2}}}+\frac{2}{{{K}^{3}}}\frac{{{v}_{x}}}{c}\frac{\overrightarrow{v}}{c} \\ 
\end{align}

{]} Выражения производной второго слагаемого электрического поля по
остальным координатам \(y\) и \(z\) точки наблюдения легко получить
меняя в данном выражении \({{R}_{x}}\) \({{a}_{x}}\) \({{v}_{x}}\)
\(\overrightarrow{i}\) на \({{R}_{y}}\) \({{a}_{y}}\) \({{v}_{y}}\)
\(\overrightarrow{j}\)и на \({{R}_{z}}\) \({{a}_{z}}\) \({{v}_{z}}\)
\(\overrightarrow{k}\)соответственно Возьмём теперь производную \(x\)
компоненты второго слагаемого электрического поля по \(x\) координате
точки наблюдения {[}

\begin{align}
  & \frac{1}{e}{{\left( \frac{\partial {{E}_{x}}}{\partial x} \right)}_{2}}=\frac{R}{{{K}^{3}}}\frac{1}{{{c}^{2}}}\left\{ \left( \frac{\overrightarrow{R}}{c}\frac{\partial \overrightarrow{a}}{\partial t'} \right)\frac{{{R}_{x}}}{K}\frac{{{v}_{x}}}{c}+\frac{{{R}_{x}}}{c}\frac{\partial {{a}_{x}}}{\partial t'}-{{a}_{x}}\frac{{{v}_{x}}}{c}-2\frac{{{v}_{x}}}{c}{{a}_{x}}-3\left( \overrightarrow{v}\frac{\overrightarrow{a}}{c} \right)\frac{{{R}_{x}}}{K}\frac{{{v}_{x}}}{c} \right\}- \\ 
 & -\left( 1+\frac{\overrightarrow{R}\cdot \overrightarrow{a}-{{v}^{2}}}{{{c}^{2}}} \right)\left( 3\frac{{{R}_{x}}}{{{K}^{4}}}\frac{{{v}_{x}}}{c}-\frac{3R}{{{K}^{4}}}\left\{ \frac{{{R}_{x}}}{K}\frac{{{v}_{x}}}{c}\left( 1+\frac{\overrightarrow{R}\cdot \overrightarrow{a}-{{v}^{2}}}{{{c}^{2}}} \right)-\frac{{{v}_{x}}}{c}\frac{{{v}_{x}}}{c} \right\}-\frac{3R}{{{K}^{3}}{{c}^{2}}}\frac{{{R}_{x}}}{K}{{a}_{x}} \right)-2\frac{{{R}_{x}}}{{{K}^{3}}}\frac{{{a}_{x}}}{{{c}^{2}}}+\frac{2}{{{K}^{3}}}\frac{{{v}_{x}}}{c}\frac{{{x}_{x}}}{c} \\ 
\end{align}

{]}

{[}

\begin{align}
  & \frac{1}{e}\frac{\partial {{E}_{2}}_{_{x}}}{\partial x}=\frac{R}{{{K}^{3}}}\frac{1}{{{c}^{2}}}\left\{ \left( \frac{\overrightarrow{R}}{c}\frac{\partial \overrightarrow{a}}{\partial t'} \right)\frac{{{R}_{x}}}{K}\frac{{{v}_{x}}}{c}+\frac{{{R}_{x}}}{c}\frac{\partial {{a}_{x}}}{\partial t'}-3\frac{{{v}_{x}}}{c}{{a}_{x}}-3\left( \overrightarrow{v}\frac{\overrightarrow{a}}{c} \right)\frac{{{R}_{x}}}{K}\frac{{{v}_{x}}}{c} \right\}- \\ 
 & -\left( 1+\frac{\overrightarrow{R}\cdot \overrightarrow{a}-{{v}^{2}}}{{{c}^{2}}} \right)\left( 3\frac{{{R}_{x}}}{{{K}^{4}}}\frac{{{v}_{x}}}{c}-\frac{3R}{{{K}^{4}}}\left\{ \frac{{{R}_{x}}}{K}\frac{{{v}_{x}}}{c}\left( 1+\frac{\overrightarrow{R}\cdot \overrightarrow{a}-{{v}^{2}}}{{{c}^{2}}} \right)-\frac{{{v}_{x}}}{c}\frac{{{v}_{x}}}{c} \right\}-\frac{3R}{{{K}^{3}}{{c}^{2}}}\frac{{{R}_{x}}}{K}{{a}_{x}} \right)-2\frac{{{R}_{x}}}{{{K}^{3}}}\frac{{{a}_{x}}}{{{c}^{2}}}+\frac{2}{{{K}^{3}}}\frac{{{v}_{x}}}{c}\frac{{{x}_{x}}}{c} \\ 
\end{align}

{]}

{[}

\begin{align}
  & \frac{1}{e}\frac{\partial {{E}_{2}}_{_{x}}}{\partial x}=\frac{R}{{{K}^{3}}}\frac{1}{{{c}^{2}}}\left\{ \left( \frac{\overrightarrow{R}}{c}\frac{\partial \overrightarrow{a}}{\partial t'} \right)\frac{{{R}_{x}}}{K}\frac{{{v}_{x}}}{c}+\frac{{{R}_{x}}}{c}\frac{\partial {{a}_{x}}}{\partial t'}-3\frac{{{v}_{x}}}{c}{{a}_{x}}-3\left( \overrightarrow{v}\frac{\overrightarrow{a}}{c} \right)\frac{{{R}_{x}}}{K}\frac{{{v}_{x}}}{c} \right\}- \\ 
 & -\left( 1+\frac{\overrightarrow{R}\cdot \overrightarrow{a}-{{v}^{2}}}{{{c}^{2}}} \right)\frac{1}{{{K}^{3}}}\left( 3\frac{{{R}_{x}}}{K}\frac{{{v}_{x}}}{c}-\frac{3R}{K}\left\{ \frac{{{R}_{x}}}{K}\frac{{{v}_{x}}}{c}\left( 1+\frac{\overrightarrow{R}\cdot \overrightarrow{a}-{{v}^{2}}}{{{c}^{2}}} \right)-\frac{{{v}_{x}}}{c}\frac{{{v}_{x}}}{c} \right\}-\frac{3R}{K}\frac{{{R}_{x}}}{{{c}^{2}}}{{a}_{x}} \right)-2\frac{{{R}_{x}}}{{{K}^{3}}}\frac{{{a}_{x}}}{{{c}^{2}}}+\frac{2}{{{K}^{3}}}\frac{{{v}_{x}}}{c}\frac{{{x}_{x}}}{c} \\ 
\end{align}

{]} Выражения производных \(y\) и \(z\) компонент второго слагаемого
электрического поля по \(y\)и \(z\) координатам точки наблюдения легко
получить аналогичной заменой индексов в данном выражении Скалярно
складывая производные\(x\) , \(y\) и \(z\)компонент второго слагаемого
электрического поля по соответственно \(x\) , \(y\) и \(z\)координатам
точки наблюдения получаем выражение для дивергенции второго слагаемого
электрического поля {[}

\begin{align}
  & \frac{1}{e}{{\left( \frac{\partial {{E}_{x}}}{\partial x}+\frac{\partial {{E}_{y}}}{\partial y}+\frac{\partial {{E}_{z}}}{\partial z} \right)}_{2}}=\frac{R}{{{K}^{3}}}\frac{1}{{{c}^{2}}}\left\{ \left( \frac{\overrightarrow{R}}{c}\frac{\partial \overrightarrow{a}}{\partial t'} \right)\frac{\overrightarrow{R}}{K}\frac{\overrightarrow{v}}{c}+\frac{\overrightarrow{R}}{c}\frac{\partial \overrightarrow{a}}{\partial t'}-3\frac{\overrightarrow{v}}{c}\overrightarrow{a}-3\left( \overrightarrow{v}\frac{\overrightarrow{a}}{c} \right)\frac{\overrightarrow{R}}{K}\frac{\overrightarrow{v}}{c} \right\}- \\ 
 & -\left( 1+\frac{\overrightarrow{R}\cdot \overrightarrow{a}-{{v}^{2}}}{{{c}^{2}}} \right)\left( 3\frac{\overrightarrow{R}}{{{K}^{4}}}\frac{\overrightarrow{v}}{c}-\frac{3R}{{{K}^{4}}}\left\{ \frac{\overrightarrow{R}}{K}\frac{\overrightarrow{v}}{c}\left( 1+\frac{\overrightarrow{R}\cdot \overrightarrow{a}-{{v}^{2}}}{{{c}^{2}}} \right)-\frac{\overrightarrow{v}}{c}\frac{\overrightarrow{v}}{c} \right\}-\frac{3R\overrightarrow{R}}{{{K}^{4}}{{c}^{2}}}\overrightarrow{a} \right)-2\frac{\overrightarrow{R}}{{{K}^{3}}}\frac{\overrightarrow{a}}{{{c}^{2}}}+\frac{2}{{{K}^{3}}}\frac{\overrightarrow{v}}{c}\frac{\overrightarrow{v}}{c} \\ 
\end{align}

{]}

{[}

\begin{align}
  & \frac{1}{e}{{\left( \frac{\partial {{E}_{x}}}{\partial x}+\frac{\partial {{E}_{y}}}{\partial y}+\frac{\partial {{E}_{z}}}{\partial z} \right)}_{2}}=\frac{R}{{{K}^{3}}}\frac{1}{{{c}^{2}}}\left\{ \left( \frac{\overrightarrow{R}}{c}\frac{\partial \overrightarrow{a}}{\partial t'} \right)\left( \frac{\overrightarrow{R}}{K}\frac{\overrightarrow{v}}{c}+1 \right)-3\left( \overrightarrow{v}\frac{\overrightarrow{a}}{c} \right)\left( \frac{\overrightarrow{R}}{K}\frac{\overrightarrow{v}}{c}+1 \right) \right\}- \\ 
 & -\left( 1+\frac{\overrightarrow{R}\cdot \overrightarrow{a}-{{v}^{2}}}{{{c}^{2}}} \right)\left( 3\frac{\overrightarrow{R}}{{{K}^{4}}}\frac{\overrightarrow{v}}{c}-\frac{3R}{{{K}^{4}}}\left\{ \frac{\overrightarrow{R}}{K}\frac{\overrightarrow{v}}{c}\left( 1+\frac{\overrightarrow{R}\cdot \overrightarrow{a}-{{v}^{2}}}{{{c}^{2}}} \right)-\frac{\overrightarrow{v}}{c}\frac{\overrightarrow{v}}{c} \right\}-\frac{3R\overrightarrow{R}}{{{K}^{4}}{{c}^{2}}}\overrightarrow{a} \right)-2\frac{\overrightarrow{R}}{{{K}^{3}}}\frac{\overrightarrow{a}}{{{c}^{2}}}+\frac{2}{{{K}^{3}}}\frac{\overrightarrow{v}}{c}\frac{\overrightarrow{v}}{c} \\ 
\end{align}

{]}

{[}

\begin{align}
  & \frac{1}{e}{{\left( \frac{\partial {{E}_{x}}}{\partial x}+\frac{\partial {{E}_{y}}}{\partial y}+\frac{\partial {{E}_{z}}}{\partial z} \right)}_{2}}=\frac{R}{{{K}^{3}}}\frac{1}{{{c}^{2}}}\left\{ \left( \frac{\overrightarrow{R}}{c}\frac{\partial \overrightarrow{a}}{\partial t'} \right)\left( \frac{\overrightarrow{R}}{K}\frac{\overrightarrow{v}}{c}+1 \right)-3\left( \overrightarrow{v}\frac{\overrightarrow{a}}{c} \right)\left( \frac{\overrightarrow{R}}{K}\frac{\overrightarrow{v}}{c}+1 \right) \right\}- \\ 
 & -\left( 1+\frac{\overrightarrow{R}\cdot \overrightarrow{a}-{{v}^{2}}}{{{c}^{2}}} \right)\left( 3\frac{\overrightarrow{R}}{{{K}^{4}}}\frac{\overrightarrow{v}}{c}-\frac{3R}{{{K}^{4}}}\left\{ \frac{\overrightarrow{R}}{K}\frac{\overrightarrow{v}}{c}\left( 1+\frac{\overrightarrow{R}\cdot \overrightarrow{a}-{{v}^{2}}}{{{c}^{2}}} \right)-\frac{\overrightarrow{v}}{c}\frac{\overrightarrow{v}}{c} \right\}-\frac{3R\overrightarrow{R}}{{{K}^{4}}{{c}^{2}}}\overrightarrow{a} \right)-2\frac{\overrightarrow{R}}{{{K}^{3}}}\frac{\overrightarrow{a}}{{{c}^{2}}}+\frac{2}{{{K}^{3}}}\frac{\overrightarrow{v}}{c}\frac{\overrightarrow{v}}{c} \\ 
\end{align}

{]} Сумма двух компонент электрического поля
{[}\overrightarrow{E}=-\nabla \varphi -\frac{1}{c}\frac{\partial \overrightarrow{A}}{\partial t}=\frac{e}{{{K}^{3}}}\left[ \left( \overrightarrow{R}-R\frac{\overrightarrow{v}}{c} \right)\left( 1+\frac{\overrightarrow{a}\cdot \overrightarrow{R}}{{{c}^{2}}}-\frac{\overrightarrow{v}\cdot \overrightarrow{v}}{{{c}^{2}}} \right)-KR\frac{\overrightarrow{a}}{{{c}^{2}}} \right]{]}
{[}\overrightarrow{E}=\frac{e}{{{K}^{3}}}\left(
\overrightarrow{R}-R\frac{\overrightarrow{v}}{c} \right)\left(
1+\frac{\overrightarrow{a}\cdot \overrightarrow{R}}{{{c}^{2}}}-\frac{\overrightarrow{v}\cdot \overrightarrow{v}}{{{c}^{2}}}
\right)-\frac{e}{{{K}^{2}}}R\frac{\overrightarrow{a}}{{{c}^{2}}}{]}
{[}\frac{1}{e}\frac{\partial \overrightarrow{E}}{\partial x}=\frac{\partial }{\partial x}\left[ \frac{1}{{{K}^{3}}}\left( \overrightarrow{R}-R\frac{\overrightarrow{v}}{c} \right)\left( 1+\frac{\overrightarrow{a}\cdot \overrightarrow{R}}{{{c}^{2}}}-\frac{\overrightarrow{v}\cdot \overrightarrow{v}}{{{c}^{2}}} \right)-\frac{R}{{{K}^{2}}}\frac{\overrightarrow{a}}{{{c}^{2}}} \right]{]}
{[}\frac{1}{e}\frac{\partial \overrightarrow{E}}{\partial x}=\frac{1}{{{K}^{3}}}\left(
\overrightarrow{R}-R\frac{\overrightarrow{v}}{c}
\right)\frac{\partial }{\partial x}\left(
1+\frac{\overrightarrow{a}\cdot \overrightarrow{R}}{{{c}^{2}}}-\frac{\overrightarrow{v}\cdot \overrightarrow{v}}{{{c}^{2}}}
\right)+\left(
1+\frac{\overrightarrow{a}\cdot \overrightarrow{R}}{{{c}^{2}}}-\frac{\overrightarrow{v}\cdot \overrightarrow{v}}{{{c}^{2}}}
\right)\frac{\partial }{\partial x}\left\{ \frac{1}{{{K}^{3}}}\left(
\overrightarrow{R}-R\frac{\overrightarrow{v}}{c} \right)
\right\}-\frac{\partial }{\partial x}\left(
\frac{R}{{{K}^{2}}}\frac{\overrightarrow{a}}{{{c}^{2}}} \right){]}
{[}\frac{1}{e}\frac{\partial \overrightarrow{E}}{\partial x}=\frac{1}{{{K}^{3}}}\left(
\overrightarrow{R}-R\frac{\overrightarrow{v}}{c}
\right)\frac{\partial }{\partial x}\left(
1+\frac{\overrightarrow{a}\cdot \overrightarrow{R}}{{{c}^{2}}}-\frac{\overrightarrow{v}\cdot \overrightarrow{v}}{{{c}^{2}}}
\right)+\left(
1+\frac{\overrightarrow{a}\cdot \overrightarrow{R}}{{{c}^{2}}}-\frac{\overrightarrow{v}\cdot \overrightarrow{v}}{{{c}^{2}}}
\right)\frac{\partial }{\partial x}\left\{
\frac{\overrightarrow{R}}{{{K}^{3}}}-\frac{R}{{{K}^{3}}}\frac{\overrightarrow{v}}{c}
\right\}-\frac{\partial }{\partial x}\left(
\frac{R}{{{K}^{2}}}\frac{\overrightarrow{a}}{{{c}^{2}}} \right){]}..
{[}\frac{\partial }{\partial x}\left(
1+\frac{\overrightarrow{a}\cdot \overrightarrow{R}}{{{c}^{2}}}-\frac{\overrightarrow{v}\cdot \overrightarrow{v}}{{{c}^{2}}}
\right)=\frac{1}{{{c}^{2}}}\left\{
-\frac{\overrightarrow{R}}{c}\frac{\partial \overrightarrow{a}}{\partial t'}\frac{{{R}_{x}}}{K}+\{\{a\}\_\{x\}\}+3\overrightarrow{v}\frac{\overrightarrow{a}}{c}\frac{{{R}_{x}}}{K}
\right\}{]}. {[}\frac{\partial }{\partial x}\left(
\frac{\overrightarrow{R}}{{{K}^{3}}} \right)=\frac{1}{{{K}^{3}}}\left(
\overrightarrow{i}+\frac{\overrightarrow{v}}{c}\frac{{{R}_{x}}}{K}
\right)-\overrightarrow{R}\frac{3}{{{K}^{4}}}\left(
\frac{{{R}_{x}}}{K}+\left(
\frac{\overrightarrow{R}\cdot \overrightarrow{a}-{{v}^{2}}}{{{c}^{2}}}
\right)\frac{{{R}_{x}}}{K}-\frac{{{v}_{x}}}{c} \right){]}
{[}\frac{\partial }{\partial x}\left(
\frac{R}{{{K}^{3}}}\frac{\overrightarrow{v}}{c}
\right)=\frac{1}{{{K}^{4}}c}\left(
\overrightarrow{v}\left[ {{R}_{x}}-3R\left\{ \frac{{{R}_{x}}}{K}\left( 1+\frac{\overrightarrow{R}\cdot \overrightarrow{a}-{{v}^{2}}}{{{c}^{2}}} \right)-\frac{{{v}_{x}}}{c} \right\} \right]-\frac{R{{R}_{x}}}{c}\overrightarrow{a}
\right){]} {[}

\begin{align}
  & \frac{\partial }{\partial x}\left\{ \frac{\overrightarrow{R}}{{{K}^{3}}}-\frac{R}{{{K}^{3}}}\frac{\overrightarrow{v}}{c} \right\}=\frac{1}{{{K}^{3}}}\left( \overrightarrow{i}+\frac{\overrightarrow{v}}{c}\frac{{{R}_{x}}}{K} \right)-\overrightarrow{R}\frac{3}{{{K}^{4}}}\left( \frac{{{R}_{x}}}{K}+\left( \frac{\overrightarrow{R}\cdot \overrightarrow{a}-{{v}^{2}}}{{{c}^{2}}} \right)\frac{{{R}_{x}}}{K}-\frac{{{v}_{x}}}{c} \right)- \\ 
 & -\frac{1}{{{K}^{4}}c}\left( \overrightarrow{v}\left[ {{R}_{x}}-3R\left\{ \frac{{{R}_{x}}}{K}\left( 1+\frac{\overrightarrow{R}\cdot \overrightarrow{a}-{{v}^{2}}}{{{c}^{2}}} \right)-\frac{{{v}_{x}}}{c} \right\} \right]-\frac{R{{R}_{x}}}{c}\overrightarrow{a} \right)= \\ 
 &  \\ 
 & =\left( \overrightarrow{i}\frac{1}{{{K}^{3}}}+\frac{\overrightarrow{v}}{c}\frac{{{R}_{x}}}{{{K}^{4}}} \right)-\overrightarrow{R}\frac{3}{{{K}^{4}}}\left( \frac{{{R}_{x}}}{K}+\left( \frac{\overrightarrow{R}\cdot \overrightarrow{a}-{{v}^{2}}}{{{c}^{2}}} \right)\frac{{{R}_{x}}}{K}-\frac{{{v}_{x}}}{c} \right)- \\ 
 & -\left[ \frac{{{R}_{x}}}{{{K}^{4}}}\frac{\overrightarrow{v}}{c}-\frac{\overrightarrow{v}}{c}\frac{3R}{{{K}^{4}}}\left\{ \frac{{{R}_{x}}}{K}\left( 1+\frac{\overrightarrow{R}\cdot \overrightarrow{a}-{{v}^{2}}}{{{c}^{2}}} \right)-\frac{{{v}_{x}}}{c} \right\} \right]+\frac{R{{R}_{x}}}{{{K}^{4}}}\frac{\overrightarrow{a}}{{{c}^{2}}}= \\ 
 &  \\ 
 & =\left( \overrightarrow{i}\frac{1}{{{K}^{3}}} \right)-\overrightarrow{R}\frac{3}{{{K}^{4}}}\left( \frac{{{R}_{x}}}{K}\left( 1+\frac{\overrightarrow{R}\cdot \overrightarrow{a}-{{v}^{2}}}{{{c}^{2}}} \right)-\frac{{{v}_{x}}}{c} \right)+ \\ 
 & +\frac{\overrightarrow{v}}{c}\frac{3R}{{{K}^{4}}}\left\{ \frac{{{R}_{x}}}{K}\left( 1+\frac{\overrightarrow{R}\cdot \overrightarrow{a}-{{v}^{2}}}{{{c}^{2}}} \right)-\frac{{{v}_{x}}}{c} \right\}+\frac{R{{R}_{x}}}{{{K}^{4}}}\frac{\overrightarrow{a}}{{{c}^{2}}}= \\ 
 &  \\ 
 & =\left( \overrightarrow{i}\frac{1}{{{K}^{3}}} \right)+\left( \frac{\overrightarrow{v}}{c}\frac{3R}{{{K}^{4}}}-\overrightarrow{R}\frac{3}{{{K}^{4}}} \right)\left\{ \frac{{{R}_{x}}}{K}\left( 1+\frac{\overrightarrow{R}\cdot \overrightarrow{a}-{{v}^{2}}}{{{c}^{2}}} \right)-\frac{{{v}_{x}}}{c} \right\}+\frac{R{{R}_{x}}}{{{K}^{4}}}\frac{\overrightarrow{a}}{{{c}^{2}}} \\ 
\end{align}

{]} {[}\frac{\partial }{\partial x}\left\{
\frac{\overrightarrow{R}}{{{K}^{3}}}-\frac{R}{{{K}^{3}}}\frac{\overrightarrow{v}}{c}
\right\}=\frac{1}{{{K}^{3}}}\left[ \overrightarrow{i}+\left( \frac{\overrightarrow{v}}{c}\frac{3R}{K}-\overrightarrow{R}\frac{3}{K} \right)\left\{ \frac{{{R}_{x}}}{K}\left( 1+\frac{\overrightarrow{R}\cdot \overrightarrow{a}-{{v}^{2}}}{{{c}^{2}}} \right)-\frac{{{v}_{x}}}{c} \right\}+\frac{R{{R}_{x}}}{K}\frac{\overrightarrow{a}}{{{c}^{2}}} \right]{]}
{[}\frac{\partial }{\partial x}\left(
\frac{R}{{{K}^{2}}}\frac{\overrightarrow{a}}{{{c}^{2}}}
\right)=\frac{1}{{{K}^{3}}{{c}^{2}}}\left(
\overrightarrow{a}\left[ {{R}_{x}}-2R\left\{ \frac{{{R}_{x}}}{K}\left( 1+\frac{\overrightarrow{R}\cdot \overrightarrow{a}-{{v}^{2}}}{{{c}^{2}}} \right)-\frac{{{v}_{x}}}{c} \right\} \right]-\frac{R{{R}_{x}}}{c}\frac{\partial \overrightarrow{a}}{\partial t'}
\right){]} {[}

\begin{align}
  & \frac{1}{e}\frac{\partial \overrightarrow{E}}{\partial x}=\frac{1}{{{K}^{3}}}\left( \overrightarrow{R}-R\frac{\overrightarrow{v}}{c} \right)\frac{1}{{{c}^{2}}}\left\{ -\left( \frac{\overrightarrow{R}}{c}\frac{\partial \overrightarrow{a}}{\partial t'} \right)\frac{{{R}_{x}}}{K}+{{a}_{x}}+3\left( \overrightarrow{v}\frac{\overrightarrow{a}}{c} \right)\frac{{{R}_{x}}}{K} \right\}+ \\ 
 & +\left( 1+\frac{\overrightarrow{a}\cdot \overrightarrow{R}}{{{c}^{2}}}-\frac{\overrightarrow{v}\cdot \overrightarrow{v}}{{{c}^{2}}} \right)\frac{1}{{{K}^{3}}}\left[ \overrightarrow{i}+\left( \frac{\overrightarrow{v}}{c}\frac{3R}{K}-\overrightarrow{R}\frac{3}{K} \right)\left\{ \frac{{{R}_{x}}}{K}\left( 1+\frac{\overrightarrow{R}\cdot \overrightarrow{a}-{{v}^{2}}}{{{c}^{2}}} \right)-\frac{{{v}_{x}}}{c} \right\}+\frac{R{{R}_{x}}}{K}\frac{\overrightarrow{a}}{{{c}^{2}}} \right]- \\ 
 & -\frac{1}{{{K}^{3}}{{c}^{2}}}\left( \overrightarrow{a}\left[ {{R}_{x}}-2R\left\{ \frac{{{R}_{x}}}{K}\left( 1+\frac{\overrightarrow{R}\cdot \overrightarrow{a}-{{v}^{2}}}{{{c}^{2}}} \right)-\frac{{{v}_{x}}}{c} \right\} \right]-\frac{R{{R}_{x}}}{c}\frac{\partial \overrightarrow{a}}{\partial t'} \right) \\ 
\end{align}

{]} Выражения производной электрического поля по остальным координатам
\(y\) и \(z\) точки наблюдения легко получить меняя в данном выражении
\({{R}_{x}}\) \({{a}_{x}}\) \({{v}_{x}}\) \(\overrightarrow{i}\) на
\({{R}_{y}}\) \({{a}_{y}}\) \({{v}_{y}}\) \(\overrightarrow{j}\)и на
\({{R}_{z}}\) \({{a}_{z}}\) \({{v}_{z}}\)
\(\overrightarrow{k}\)соответственно Возьмём теперь производную
компоненты электрического поля по координате точки наблюдения {[}

\begin{align}
  & \frac{1}{e}\frac{\partial {{E}_{x}}}{\partial x}=\frac{1}{{{K}^{3}}}\frac{1}{{{c}^{2}}}\left( {{R}_{x}}-R\frac{{{v}_{x}}}{c} \right)\left\{ -\left( \frac{\overrightarrow{R}}{c}\frac{\partial \overrightarrow{a}}{\partial t'} \right)\frac{{{R}_{x}}}{K}+{{a}_{x}}+3\left( \overrightarrow{v}\frac{\overrightarrow{a}}{c} \right)\frac{{{R}_{x}}}{K} \right\}+ \\ 
 & +\left( 1+\frac{\overrightarrow{a}\cdot \overrightarrow{R}}{{{c}^{2}}}-\frac{\overrightarrow{v}\cdot \overrightarrow{v}}{{{c}^{2}}} \right)\frac{1}{{{K}^{3}}}\left[ 1+\left[ \left( \frac{{{v}_{x}}}{c}\frac{3R}{K}-{{R}_{x}}\frac{3}{K} \right)\left\{ \frac{{{R}_{x}}}{K}\left( 1+\frac{\overrightarrow{R}\cdot \overrightarrow{a}-{{v}^{2}}}{{{c}^{2}}} \right)-\frac{{{v}_{x}}}{c} \right\} \right]+\frac{R{{R}_{x}}}{K}\frac{{{a}_{x}}}{{{c}^{2}}} \right]- \\ 
 & -\frac{1}{{{K}^{3}}{{c}^{2}}}\left( {{a}_{x}}\left[ {{R}_{x}}-2R\left\{ \frac{{{R}_{x}}}{K}\left( 1+\frac{\overrightarrow{R}\cdot \overrightarrow{a}-{{v}^{2}}}{{{c}^{2}}} \right)-\frac{{{v}_{x}}}{c} \right\} \right]-\frac{R{{R}_{x}}}{c}\frac{\partial {{a}_{x}}}{\partial t'} \right) \\ 
\end{align}

{]}

{[}

\begin{align}
  & \frac{1}{e}\frac{\partial {{E}_{x}}}{\partial x}=\frac{1}{{{c}^{2}}}\frac{1}{{{K}^{3}}}\left[ {{R}_{x}}\left\{ -\left( \frac{\overrightarrow{R}}{c}\frac{\partial \overrightarrow{a}}{\partial t'} \right)\frac{{{R}_{x}}}{K}+{{a}_{x}}+3\left( \overrightarrow{v}\frac{\overrightarrow{a}}{c} \right)\frac{{{R}_{x}}}{K} \right\}-R\frac{{{v}_{x}}}{c}\left\{ -\left( \frac{\overrightarrow{R}}{c}\frac{\partial \overrightarrow{a}}{\partial t'} \right)\frac{{{R}_{x}}}{K}+{{a}_{x}}+3\left( \overrightarrow{v}\frac{\overrightarrow{a}}{c} \right)\frac{{{R}_{x}}}{K} \right\} \right]+ \\ 
 & +\left( 1+\frac{\overrightarrow{a}\cdot \overrightarrow{R}}{{{c}^{2}}}-\frac{\overrightarrow{v}\cdot \overrightarrow{v}}{{{c}^{2}}} \right)\frac{1}{{{K}^{3}}}\left[ 1+\left[ \left( \frac{{{v}_{x}}}{c}\frac{3R}{K}-{{R}_{x}}\frac{3}{K} \right)\left\{ \frac{{{R}_{x}}}{K}\left( 1+\frac{\overrightarrow{R}\cdot \overrightarrow{a}-{{v}^{2}}}{{{c}^{2}}} \right)-\frac{{{v}_{x}}}{c} \right\} \right]+\frac{R{{R}_{x}}}{K}\frac{{{a}_{x}}}{{{c}^{2}}} \right]- \\ 
 & -\frac{1}{{{K}^{3}}{{c}^{2}}}\left( {{a}_{x}}{{R}_{x}}-2R\left\{ \frac{{{a}_{x}}{{R}_{x}}}{K}\left( 1+\frac{\overrightarrow{R}\cdot \overrightarrow{a}-{{v}^{2}}}{{{c}^{2}}} \right)-\frac{{{a}_{x}}{{v}_{x}}}{c} \right\}-\frac{R{{R}_{x}}}{c}\frac{\partial {{a}_{x}}}{\partial t'} \right) \\ 
\end{align}

{]}

{[}

\begin{align}
  & \frac{1}{e}\frac{\partial {{E}_{x}}}{\partial x}=\frac{1}{{{c}^{2}}}\frac{1}{{{K}^{3}}}\left[ -\left( \frac{\overrightarrow{R}}{c}\frac{\partial \overrightarrow{a}}{\partial t'} \right)\frac{{{R}_{x}}{{R}_{x}}}{K}+{{a}_{x}}{{R}_{x}}+3\left( \overrightarrow{v}\frac{\overrightarrow{a}}{c} \right)\frac{{{R}_{x}}{{R}_{x}}}{K}+\left( \frac{\overrightarrow{R}}{c}\frac{\partial \overrightarrow{a}}{\partial t'} \right)\frac{R}{c}\frac{{{R}_{x}}{{v}_{x}}}{K}-\frac{R}{c}{{a}_{x}}{{v}_{x}}-3\frac{R}{c}\left( \overrightarrow{v}\frac{\overrightarrow{a}}{c} \right)\frac{{{R}_{x}}{{v}_{x}}}{K} \right]+ \\ 
 & +\left( 1+\frac{\overrightarrow{a}\cdot \overrightarrow{R}}{{{c}^{2}}}-\frac{\overrightarrow{v}\cdot \overrightarrow{v}}{{{c}^{2}}} \right)\frac{1}{{{K}^{3}}}\left[ 1+\left[ \left( \frac{{{v}_{x}}}{c}\frac{3R}{K}-{{R}_{x}}\frac{3}{K} \right)\left\{ \frac{{{R}_{x}}}{K}\left( 1+\frac{\overrightarrow{R}\cdot \overrightarrow{a}-{{v}^{2}}}{{{c}^{2}}} \right)-\frac{{{v}_{x}}}{c} \right\} \right]+\frac{R{{R}_{x}}}{K}\frac{{{a}_{x}}}{{{c}^{2}}} \right]- \\ 
 & -\frac{1}{{{K}^{3}}{{c}^{2}}}\left( {{a}_{x}}{{R}_{x}}-2R\left\{ \frac{{{a}_{x}}{{R}_{x}}}{K}\left( 1+\frac{\overrightarrow{R}\cdot \overrightarrow{a}-{{v}^{2}}}{{{c}^{2}}} \right)-\frac{{{a}_{x}}{{v}_{x}}}{c} \right\}-\frac{R{{R}_{x}}}{c}\frac{\partial {{a}_{x}}}{\partial t'} \right) \\ 
\end{align}

{]}

{[}

\begin{align}
  & \frac{1}{e}\frac{\partial {{E}_{x}}}{\partial x}=\frac{1}{{{c}^{2}}}\frac{1}{{{K}^{3}}}\left[ -\left( \frac{\overrightarrow{R}}{c}\frac{\partial \overrightarrow{a}}{\partial t'} \right)\frac{{{R}_{x}}{{R}_{x}}}{K}+3\left( \overrightarrow{v}\frac{\overrightarrow{a}}{c} \right)\frac{{{R}_{x}}{{R}_{x}}}{K}+\left( \frac{\overrightarrow{R}}{c}\frac{\partial \overrightarrow{a}}{\partial t'} \right)\frac{R}{c}\frac{{{R}_{x}}{{v}_{x}}}{K}-\frac{R}{c}{{a}_{x}}{{v}_{x}}-3\frac{R}{c}\left( \overrightarrow{v}\frac{\overrightarrow{a}}{c} \right)\frac{{{R}_{x}}{{v}_{x}}}{K} \right]+ \\ 
 & +\left( 1+\frac{\overrightarrow{a}\cdot \overrightarrow{R}}{{{c}^{2}}}-\frac{\overrightarrow{v}\cdot \overrightarrow{v}}{{{c}^{2}}} \right)\frac{1}{{{K}^{3}}}\left[ 1+\left[ \left( \frac{{{v}_{x}}}{c}\frac{3R}{K}-{{R}_{x}}\frac{3}{K} \right)\left\{ \frac{{{R}_{x}}}{K}\left( 1+\frac{\overrightarrow{R}\cdot \overrightarrow{a}-{{v}^{2}}}{{{c}^{2}}} \right)-\frac{{{v}_{x}}}{c} \right\} \right]+\frac{R{{R}_{x}}}{K}\frac{{{a}_{x}}}{{{c}^{2}}} \right]+ \\ 
 & +\frac{1}{{{K}^{3}}{{c}^{2}}}\left( 2R\left\{ \frac{{{a}_{x}}{{R}_{x}}}{K}\left( 1+\frac{\overrightarrow{R}\cdot \overrightarrow{a}-{{v}^{2}}}{{{c}^{2}}} \right)-\frac{{{a}_{x}}{{v}_{x}}}{c} \right\}+\frac{R{{R}_{x}}}{c}\frac{\partial {{a}_{x}}}{\partial t'} \right) \\ 
\end{align}

{]}

{[}

\begin{align}
  & \frac{1}{e}\frac{\partial {{E}_{x}}}{\partial x}=\frac{1}{{{c}^{2}}}\frac{1}{{{K}^{3}}}\left[ -\left( \frac{\overrightarrow{R}}{c}\frac{\partial \overrightarrow{a}}{\partial t'} \right)\frac{{{R}_{x}}{{R}_{x}}}{K}+3\left( \overrightarrow{v}\frac{\overrightarrow{a}}{c} \right)\frac{{{R}_{x}}{{R}_{x}}}{K}+\left( \frac{\overrightarrow{R}}{c}\frac{\partial \overrightarrow{a}}{\partial t'} \right)\frac{R}{c}\frac{{{R}_{x}}{{v}_{x}}}{K}-\frac{R}{c}{{a}_{x}}{{v}_{x}}-3\frac{R}{c}\left( \overrightarrow{v}\frac{\overrightarrow{a}}{c} \right)\frac{{{R}_{x}}{{v}_{x}}}{K} \right]+ \\ 
 & +\left( 1+\frac{\overrightarrow{a}\cdot \overrightarrow{R}}{{{c}^{2}}}-\frac{\overrightarrow{v}\cdot \overrightarrow{v}}{{{c}^{2}}} \right)\frac{1}{{{K}^{3}}}\left[ 1+\frac{{{v}_{x}}}{c}\frac{3R}{K}\left\{ \frac{{{R}_{x}}}{K}\left( 1+\frac{\overrightarrow{R}\cdot \overrightarrow{a}-{{v}^{2}}}{{{c}^{2}}} \right)-\frac{{{v}_{x}}}{c} \right\}-{{R}_{x}}\frac{3}{K}\left\{ \frac{{{R}_{x}}}{K}\left( 1+\frac{\overrightarrow{R}\cdot \overrightarrow{a}-{{v}^{2}}}{{{c}^{2}}} \right)-\frac{{{v}_{x}}}{c} \right\}+\frac{R{{R}_{x}}}{K}\frac{{{a}_{x}}}{{{c}^{2}}} \right]+ \\ 
 & +\frac{1}{{{K}^{3}}{{c}^{2}}}\left( 2R\frac{{{a}_{x}}{{R}_{x}}}{K}\left( 1+\frac{\overrightarrow{R}\cdot \overrightarrow{a}-{{v}^{2}}}{{{c}^{2}}} \right)-\frac{2R}{c}{{a}_{x}}{{v}_{x}}+\frac{R{{R}_{x}}}{c}\frac{\partial {{a}_{x}}}{\partial t'} \right) \\ 
\end{align}

{]}

{[}

\begin{align}
  & \frac{1}{e}\frac{\partial {{E}_{x}}}{\partial x}=\frac{1}{{{c}^{2}}}\frac{1}{{{K}^{3}}}\left[ -\left( \frac{\overrightarrow{R}}{c}\frac{\partial \overrightarrow{a}}{\partial t'} \right)\frac{{{R}_{x}}{{R}_{x}}}{K}+3\left( \overrightarrow{v}\frac{\overrightarrow{a}}{c} \right)\frac{{{R}_{x}}{{R}_{x}}}{K}+\left( \frac{\overrightarrow{R}}{c}\frac{\partial \overrightarrow{a}}{\partial t'} \right)\frac{R}{c}\frac{{{R}_{x}}{{v}_{x}}}{K}-\frac{3R}{c}{{a}_{x}}{{v}_{x}}-3\frac{R}{c}\left( \overrightarrow{v}\frac{\overrightarrow{a}}{c} \right)\frac{{{R}_{x}}{{v}_{x}}}{K} \right]+ \\ 
 & +\left( 1+\frac{\overrightarrow{a}\cdot \overrightarrow{R}}{{{c}^{2}}}-\frac{\overrightarrow{v}\cdot \overrightarrow{v}}{{{c}^{2}}} \right)\frac{1}{{{K}^{3}}}\left[ 1+\frac{3R}{K}\left\{ \frac{{{v}_{x}}}{c}\frac{{{R}_{x}}}{K}\left( 1+\frac{\overrightarrow{R}\cdot \overrightarrow{a}-{{v}^{2}}}{{{c}^{2}}} \right)-\frac{{{v}_{x}}}{c}\frac{{{v}_{x}}}{c} \right\}-\frac{3}{K}\left\{ \frac{{{R}_{x}}{{R}_{x}}}{K}\left( 1+\frac{\overrightarrow{R}\cdot \overrightarrow{a}-{{v}^{2}}}{{{c}^{2}}} \right)-\frac{{{v}_{x}}{{R}_{x}}}{c} \right\}+\frac{R{{R}_{x}}}{K}\frac{{{a}_{x}}}{{{c}^{2}}} \right]+ \\ 
 & +\frac{1}{{{K}^{3}}{{c}^{2}}}\left( 2R\frac{{{a}_{x}}{{R}_{x}}}{K}\left( 1+\frac{\overrightarrow{R}\cdot \overrightarrow{a}-{{v}^{2}}}{{{c}^{2}}} \right)+\frac{R{{R}_{x}}}{c}\frac{\partial {{a}_{x}}}{\partial t'} \right) \\ 
\end{align}

{]}

{[}

\begin{align}
  & \frac{1}{e}\frac{\partial {{E}_{x}}}{\partial x}=\frac{1}{{{c}^{2}}}\frac{1}{{{K}^{3}}}\left[ -\left( \frac{\overrightarrow{R}}{c}\frac{\partial \overrightarrow{a}}{\partial t'} \right)\frac{{{R}_{x}}{{R}_{x}}}{K}+3\left( \overrightarrow{v}\frac{\overrightarrow{a}}{c} \right)\frac{{{R}_{x}}{{R}_{x}}}{K}+\left( \frac{\overrightarrow{R}}{c}\frac{\partial \overrightarrow{a}}{\partial t'} \right)\frac{R}{c}\frac{{{R}_{x}}{{v}_{x}}}{K}-\frac{3R}{c}{{a}_{x}}{{v}_{x}}-3\frac{R}{c}\left( \overrightarrow{v}\frac{\overrightarrow{a}}{c} \right)\frac{{{R}_{x}}{{v}_{x}}}{K} \right]+ \\ 
 & +\left( 1+\frac{\overrightarrow{a}\cdot \overrightarrow{R}}{{{c}^{2}}}-\frac{\overrightarrow{v}\cdot \overrightarrow{v}}{{{c}^{2}}} \right)\frac{1}{{{K}^{3}}}\left[ 1+\frac{2R}{{{c}^{2}}}\frac{{{a}_{x}}{{R}_{x}}}{K}+\frac{3R}{K}\left\{ \frac{{{v}_{x}}}{c}\frac{{{R}_{x}}}{K}\left( 1+\frac{\overrightarrow{R}\cdot \overrightarrow{a}-{{v}^{2}}}{{{c}^{2}}} \right)-\frac{{{v}_{x}}}{c}\frac{{{v}_{x}}}{c} \right\}-\frac{3}{K}\left\{ \frac{{{R}_{x}}{{R}_{x}}}{K}\left( 1+\frac{\overrightarrow{R}\cdot \overrightarrow{a}-{{v}^{2}}}{{{c}^{2}}} \right)-\frac{{{v}_{x}}{{R}_{x}}}{c} \right\}+\frac{R{{R}_{x}}}{K}\frac{{{a}_{x}}}{{{c}^{2}}} \right]+ \\ 
 & +\frac{1}{{{K}^{3}}{{c}^{2}}}\left( \frac{R{{R}_{x}}}{c}\frac{\partial {{a}_{x}}}{\partial t'} \right) \\ 
\end{align}

{]}

{[}

\begin{align}
  & \frac{1}{e}\frac{\partial {{E}_{x}}}{\partial x}=\frac{1}{{{c}^{2}}}\frac{1}{{{K}^{3}}}\left[ -\left( \frac{\overrightarrow{R}}{c}\frac{\partial \overrightarrow{a}}{\partial t'} \right)\frac{{{R}_{x}}{{R}_{x}}}{K}+3\left( \overrightarrow{v}\frac{\overrightarrow{a}}{c} \right)\frac{{{R}_{x}}{{R}_{x}}}{K}+\left( \frac{\overrightarrow{R}}{c}\frac{\partial \overrightarrow{a}}{\partial t'} \right)\frac{R}{c}\frac{{{R}_{x}}{{v}_{x}}}{K}-\frac{3R}{c}{{a}_{x}}{{v}_{x}}-3\frac{R}{c}\left( \overrightarrow{v}\frac{\overrightarrow{a}}{c} \right)\frac{{{R}_{x}}{{v}_{x}}}{K}+\frac{R{{R}_{x}}}{c}\frac{\partial {{a}_{x}}}{\partial t'} \right]+ \\ 
 & +\left( 1+\frac{\overrightarrow{a}\cdot \overrightarrow{R}}{{{c}^{2}}}-\frac{\overrightarrow{v}\cdot \overrightarrow{v}}{{{c}^{2}}} \right)\frac{1}{{{K}^{3}}}\left[ 1+\frac{2R}{{{c}^{2}}}\frac{{{a}_{x}}{{R}_{x}}}{K}+\frac{3R}{K}\left\{ \frac{{{v}_{x}}}{c}\frac{{{R}_{x}}}{K}\left( 1+\frac{\overrightarrow{R}\cdot \overrightarrow{a}-{{v}^{2}}}{{{c}^{2}}} \right)-\frac{{{v}_{x}}}{c}\frac{{{v}_{x}}}{c} \right\}-\frac{3}{K}\left\{ \frac{{{R}_{x}}{{R}_{x}}}{K}\left( 1+\frac{\overrightarrow{R}\cdot \overrightarrow{a}-{{v}^{2}}}{{{c}^{2}}} \right)-\frac{{{v}_{x}}{{R}_{x}}}{c} \right\}+\frac{R{{R}_{x}}}{K}\frac{{{a}_{x}}}{{{c}^{2}}} \right] \\ 
\end{align}

{]} Выражения производных \(y\) и \(z\) компонент электрического поля по
\(y\)и \(z\) координатам точки наблюдения легко получить аналогичной
заменой индексов в данном выражении Скалярно складывая производные\(x\)
, \(y\) и \(z\)компонент электрического поля по соответственно \(x\) ,
\(y\) и \(z\)координатам точки наблюдения получаем выражение для
дивергенции электрического поля {[}

\begin{align}
  & \frac{1}{e}\left( \frac{\partial {{E}_{x}}}{\partial x}+\frac{\partial {{E}_{y}}}{\partial y}+\frac{\partial {{E}_{z}}}{\partial z} \right)=\frac{1}{{{c}^{2}}}\frac{1}{{{K}^{3}}}\left[ -\left( \frac{\overrightarrow{R}}{c}\frac{\partial \overrightarrow{a}}{\partial t'} \right)\frac{{{R}^{2}}}{K}+3\left( \overrightarrow{v}\frac{\overrightarrow{a}}{c} \right)\frac{{{R}^{2}}}{K}+\left( \frac{\overrightarrow{R}}{c}\frac{\partial \overrightarrow{a}}{\partial t'} \right)\frac{R}{c}\frac{\overrightarrow{R}\overrightarrow{v}}{K}-\frac{3R}{c}\overrightarrow{a}\overrightarrow{v}-3\frac{R}{c}\left( \overrightarrow{v}\frac{\overrightarrow{a}}{c} \right)\frac{\overrightarrow{R}\overrightarrow{v}}{K}+\frac{R\overrightarrow{R}}{c}\frac{\partial \overrightarrow{a}}{\partial t'} \right]+ \\ 
 & +\left( 1+\frac{\overrightarrow{a}\cdot \overrightarrow{R}}{{{c}^{2}}}-\frac{\overrightarrow{v}\cdot \overrightarrow{v}}{{{c}^{2}}} \right)\frac{1}{{{K}^{3}}}\left[ 3+\frac{2R}{{{c}^{2}}}\frac{\overrightarrow{a}\overrightarrow{R}}{K}+\frac{3R}{K}\left\{ \frac{\overrightarrow{v}}{c}\frac{\overrightarrow{R}}{K}\left( 1+\frac{\overrightarrow{R}\cdot \overrightarrow{a}-{{v}^{2}}}{{{c}^{2}}} \right)-\frac{\overrightarrow{v}}{c}\frac{\overrightarrow{v}}{c} \right\}-\frac{3}{K}\left\{ \frac{{{R}^{2}}}{K}\left( 1+\frac{\overrightarrow{R}\cdot \overrightarrow{a}-{{v}^{2}}}{{{c}^{2}}} \right)-\frac{\overrightarrow{v}\overrightarrow{R}}{c} \right\}+\frac{R\overrightarrow{R}}{K}\frac{\overrightarrow{a}}{{{c}^{2}}} \right] \\ 
 &  \\ 
\end{align}

{]}

{[}

\begin{align}
  & \frac{1}{e}\left( \frac{\partial {{E}_{x}}}{\partial x}+\frac{\partial {{E}_{y}}}{\partial y}+\frac{\partial {{E}_{z}}}{\partial z} \right)=\frac{1}{{{c}^{2}}}\frac{1}{{{K}^{3}}}\left[ -\left( \frac{\overrightarrow{R}}{c}\frac{\partial \overrightarrow{a}}{\partial t'} \right)\frac{{{R}^{2}}}{K}+3\left( \overrightarrow{v}\frac{\overrightarrow{a}}{c} \right)\frac{{{R}^{2}}}{K}+\left( \frac{\overrightarrow{R}}{c}\frac{\partial \overrightarrow{a}}{\partial t'} \right)\frac{R}{c}\frac{\overrightarrow{R}\overrightarrow{v}}{K}-\frac{3R}{c}\overrightarrow{a}\overrightarrow{v}-3\frac{R}{c}\left( \overrightarrow{v}\frac{\overrightarrow{a}}{c} \right)\frac{\overrightarrow{R}\overrightarrow{v}}{K}+\frac{R\overrightarrow{R}}{c}\frac{\partial \overrightarrow{a}}{\partial t'} \right]+ \\ 
 & +\left( 1+\frac{\overrightarrow{a}\cdot \overrightarrow{R}}{{{c}^{2}}}-\frac{\overrightarrow{v}\cdot \overrightarrow{v}}{{{c}^{2}}} \right)\frac{1}{{{K}^{3}}}\left[ 3+\frac{2R}{{{c}^{2}}}\frac{\overrightarrow{a}\overrightarrow{R}}{K}+\frac{R\overrightarrow{R}}{K}\frac{\overrightarrow{a}}{{{c}^{2}}}+\frac{3R}{K}\left\{ \frac{\overrightarrow{v}}{c}\frac{\overrightarrow{R}}{K}\left( 1+\frac{\overrightarrow{R}\cdot \overrightarrow{a}-{{v}^{2}}}{{{c}^{2}}} \right)-\frac{\overrightarrow{v}}{c}\frac{\overrightarrow{v}}{c} \right\}-\frac{3}{K}\left\{ \frac{{{R}^{2}}}{K}\left( 1+\frac{\overrightarrow{R}\cdot \overrightarrow{a}-{{v}^{2}}}{{{c}^{2}}} \right)-\frac{\overrightarrow{v}\overrightarrow{R}}{c} \right\} \right] \\ 
 &  \\ 
\end{align}

{]}

{[}

\begin{align}
  & \frac{1}{e}\left( \frac{\partial {{E}_{x}}}{\partial x}+\frac{\partial {{E}_{y}}}{\partial y}+\frac{\partial {{E}_{z}}}{\partial z} \right)=\frac{1}{{{c}^{2}}}\frac{1}{{{K}^{3}}}\left[ \left( \frac{\overrightarrow{R}}{c}\frac{\partial \overrightarrow{a}}{\partial t'} \right)\frac{R}{c}\frac{\overrightarrow{R}\overrightarrow{v}}{K}-\left( \frac{\overrightarrow{R}}{c}\frac{\partial \overrightarrow{a}}{\partial t'} \right)\frac{{{R}^{2}}}{K}+\frac{R\overrightarrow{R}}{c}\frac{\partial \overrightarrow{a}}{\partial t'}+3\left( \overrightarrow{v}\frac{\overrightarrow{a}}{c} \right)\frac{{{R}^{2}}}{K}-\frac{3R}{c}\overrightarrow{a}\overrightarrow{v}-3\frac{R}{c}\left( \overrightarrow{v}\frac{\overrightarrow{a}}{c} \right)\frac{\overrightarrow{R}\overrightarrow{v}}{K} \right]+ \\ 
 & +\left( 1+\frac{\overrightarrow{a}\cdot \overrightarrow{R}}{{{c}^{2}}}-\frac{\overrightarrow{v}\cdot \overrightarrow{v}}{{{c}^{2}}} \right)\frac{1}{{{K}^{3}}}\left[ 3+\frac{2R}{{{c}^{2}}}\frac{\overrightarrow{a}\overrightarrow{R}}{K}+\frac{R\overrightarrow{R}}{K}\frac{\overrightarrow{a}}{{{c}^{2}}}+\frac{3R}{K}\frac{\overrightarrow{v}}{c}\frac{\overrightarrow{R}}{K}\left( 1+\frac{\overrightarrow{R}\cdot \overrightarrow{a}-{{v}^{2}}}{{{c}^{2}}} \right)-\frac{3R}{K}\frac{\overrightarrow{v}}{c}\frac{\overrightarrow{v}}{c}-\frac{3}{K}\frac{{{R}^{2}}}{K}\left( 1+\frac{\overrightarrow{R}\cdot \overrightarrow{a}-{{v}^{2}}}{{{c}^{2}}} \right)+\frac{3}{K}\frac{\overrightarrow{v}\overrightarrow{R}}{c} \right] \\ 
 &  \\ 
\end{align}

{]}

{[}

\begin{align}
  & \frac{1}{e}\left( \frac{\partial {{E}_{x}}}{\partial x}+\frac{\partial {{E}_{y}}}{\partial y}+\frac{\partial {{E}_{z}}}{\partial z} \right)=\frac{1}{{{c}^{2}}}\frac{1}{{{K}^{3}}}\left[ \left( \frac{\overrightarrow{R}}{c}\frac{\partial \overrightarrow{a}}{\partial t'} \right)\left\{ \frac{R}{c}\frac{\overrightarrow{R}\overrightarrow{v}}{K}-\frac{{{R}^{2}}}{K}+R \right\}+3\left( \frac{\overrightarrow{a}\overrightarrow{v}}{c} \right)\left\{ \frac{{{R}^{2}}}{K}-R-\frac{R}{c}\frac{\overrightarrow{R}\overrightarrow{v}}{K} \right\} \right]+ \\ 
 & +\left( 1+\frac{\overrightarrow{a}\cdot \overrightarrow{R}}{{{c}^{2}}}-\frac{\overrightarrow{v}\cdot \overrightarrow{v}}{{{c}^{2}}} \right)\frac{1}{{{K}^{3}}}\left[ 3+3\frac{R}{{{c}^{2}}}\frac{\overrightarrow{a}\overrightarrow{R}}{K}+\frac{3R}{K}\frac{\overrightarrow{v}}{c}\frac{\overrightarrow{R}}{K}\left( 1+\frac{\overrightarrow{R}\cdot \overrightarrow{a}-{{v}^{2}}}{{{c}^{2}}} \right)-\frac{3R}{K}\frac{\overrightarrow{v}}{c}\frac{\overrightarrow{v}}{c}-\frac{3}{K}\frac{{{R}^{2}}}{K}\left( 1+\frac{\overrightarrow{R}\cdot \overrightarrow{a}-{{v}^{2}}}{{{c}^{2}}} \right)+\frac{3}{K}\frac{\overrightarrow{v}\overrightarrow{R}}{c} \right] \\ 
 &  \\ 
\end{align}

{]} {[}\left\{
\frac{R}{c}\frac{\overrightarrow{R}\overrightarrow{v}}{K}-\frac{{{R}^{2}}}{K}+R
\right\}=\left\{
\frac{R\left( \frac{\overrightarrow{v}\cdot \overrightarrow{R}}{c} \right)}{R-\frac{\overrightarrow{v}\cdot \overrightarrow{R}}{c}}-\frac{{{R}^{2}}}{R-\frac{\overrightarrow{v}\cdot \overrightarrow{R}}{c}}+\frac{R\left( R-\frac{\overrightarrow{v}\cdot \overrightarrow{R}}{c} \right)}{R-\frac{\overrightarrow{v}\cdot \overrightarrow{R}}{c}}
\right\}=\frac{R\left( \frac{\overrightarrow{v}\cdot \overrightarrow{R}}{c} \right)-{{R}^{2}}+R\left( R-\frac{\overrightarrow{v}\cdot \overrightarrow{R}}{c} \right)}{R-\frac{\overrightarrow{v}\cdot \overrightarrow{R}}{c}}=0{]}
{[}

\begin{align}
  & \frac{1}{e}\left( \frac{\partial {{E}_{x}}}{\partial x}+\frac{\partial {{E}_{y}}}{\partial y}+\frac{\partial {{E}_{z}}}{\partial z} \right)=\left( 1+\frac{\overrightarrow{a}\cdot \overrightarrow{R}}{{{c}^{2}}}-\frac{\overrightarrow{v}\cdot \overrightarrow{v}}{{{c}^{2}}} \right)\frac{1}{{{K}^{3}}} \\ 
 & \left[ 3+3\frac{R}{{{c}^{2}}}\frac{\overrightarrow{a}\overrightarrow{R}}{K}-\frac{3R}{K}\frac{\overrightarrow{v}}{c}\frac{\overrightarrow{v}}{c}+\frac{3R}{K}\frac{\overrightarrow{v}}{c}\frac{\overrightarrow{R}}{K}\left( 1+\frac{\overrightarrow{R}\cdot \overrightarrow{a}-{{v}^{2}}}{{{c}^{2}}} \right)-\frac{3}{K}\frac{{{R}^{2}}}{K}\left( 1+\frac{\overrightarrow{R}\cdot \overrightarrow{a}-{{v}^{2}}}{{{c}^{2}}} \right)+\frac{3}{K}\frac{\overrightarrow{v}\overrightarrow{R}}{c} \right] \\ 
\end{align}

{]}

{[}

\begin{align}
  & \frac{1}{e}\left( \frac{\partial {{E}_{x}}}{\partial x}+\frac{\partial {{E}_{y}}}{\partial y}+\frac{\partial {{E}_{z}}}{\partial z} \right)=\frac{1}{{{K}^{3}}}\left( 1+\frac{\overrightarrow{a}\cdot \overrightarrow{R}}{{{c}^{2}}}-\frac{\overrightarrow{v}\cdot \overrightarrow{v}}{{{c}^{2}}} \right) \\ 
 & \left[ 3+3\frac{R}{K}\frac{\overrightarrow{a}\cdot \overrightarrow{R}}{{{c}^{2}}}-\frac{3R}{K}\frac{{{v}^{2}}}{{{c}^{2}}}+3\left\{ \frac{R}{{{K}^{2}}}\left( \frac{\overrightarrow{v}\cdot \overrightarrow{R}}{c} \right)-\frac{{{R}^{2}}}{{{K}^{2}}} \right\}\left( 1+\frac{\overrightarrow{R}\cdot \overrightarrow{a}-{{v}^{2}}}{{{c}^{2}}} \right)+\frac{3}{K}\left( \frac{\overrightarrow{v}\cdot \overrightarrow{R}}{c} \right) \right] \\ 
\end{align}

{]} {[}\left\{ \frac{R}{{{K}^{2}}}\left(
\frac{\overrightarrow{v}\cdot \overrightarrow{R}}{c}
\right)-\frac{{{R}^{2}}}{{{K}^{2}}} \right\}=\frac{R}{K}\left\{
\frac{1}{K}\left( \frac{\overrightarrow{v}\cdot \overrightarrow{R}}{c}
\right)-\frac{R}{K} \right\}=\frac{R}{K}\left\{
\frac{\frac{\overrightarrow{v}\cdot \overrightarrow{R}}{c}}{R-\frac{\overrightarrow{v}\cdot \overrightarrow{R}}{c}}-\frac{R}{R-\frac{\overrightarrow{v}\cdot \overrightarrow{R}}{c}}
\right\}=\frac{R}{K}\left\{
\frac{\frac{\overrightarrow{v}\cdot \overrightarrow{R}}{c}-R}{R-\frac{\overrightarrow{v}\cdot \overrightarrow{R}}{c}}
\right\}=-\frac{R}{K}{]} Итак, дивергенция электрического поля равна
{[}\frac{1}{e}\left(
\frac{\partial {{E}_{x}}}{\partial x}+\frac{\partial {{E}_{y}}}{\partial y}+\frac{\partial {{E}_{z}}}{\partial z}
\right)=\frac{1}{{{K}^{3}}}\left(
1+\frac{\overrightarrow{a}\cdot \overrightarrow{R}}{{{c}^{2}}}-\frac{\overrightarrow{v}\cdot \overrightarrow{v}}{{{c}^{2}}}
\right)\left[ 3+3\frac{R}{K}\frac{\overrightarrow{a}\cdot \overrightarrow{R}}{{{c}^{2}}}-\frac{3R}{K}\frac{{{v}^{2}}}{{{c}^{2}}}-3\frac{R}{K}\left( 1+\frac{\overrightarrow{R}\cdot \overrightarrow{a}-{{v}^{2}}}{{{c}^{2}}} \right)+\frac{3}{K}\left( \frac{\overrightarrow{v}\cdot \overrightarrow{R}}{c} \right) \right]{]}
При этом дивергенция первого слагаемого электрического поля {[}

\begin{align}
  & \frac{1}{e}{{\left( \frac{\partial {{E}_{x}}}{\partial x}+\frac{\partial {{E}_{y}}}{\partial y}+\frac{\partial {{E}_{z}}}{\partial z} \right)}_{1}}=-\frac{{{R}^{2}}}{{{K}^{4}}{{c}^{2}}}\left( \frac{\overrightarrow{R}}{c}\frac{\partial \overrightarrow{a}}{\partial t'} \right)+\frac{{{R}^{2}}}{{{K}^{4}}{{c}^{2}}}\left( \frac{3\ \overrightarrow{v}\cdot \overrightarrow{a}}{c} \right)-\frac{2}{{{K}^{3}}}+ \\ 
 & +\left( 1+\frac{\overrightarrow{a}\cdot \overrightarrow{R}}{{{c}^{2}}}-\frac{\overrightarrow{v}\cdot \overrightarrow{v}}{{{c}^{2}}} \right)\cdot \left\{ \frac{5}{{{K}^{3}}}-\frac{3{{R}^{2}}}{{{K}^{5}}}\left( 1+\frac{\overrightarrow{R}\cdot \overrightarrow{a}-{{v}^{2}}}{{{c}^{2}}} \right)+6\frac{\overrightarrow{v}\cdot \overrightarrow{R}}{c{{K}^{4}}} \right\} \\ 
\end{align}

{]} И дивергенция второго слагаемого {[}

\begin{align}
  & \frac{1}{e}{{\left( \frac{\partial {{E}_{x}}}{\partial x}+\frac{\partial {{E}_{y}}}{\partial y}+\frac{\partial {{E}_{z}}}{\partial z} \right)}_{2}}=\frac{R}{{{K}^{3}}}\frac{1}{{{c}^{2}}}\left\{ \left( \frac{\overrightarrow{R}}{c}\frac{\partial \overrightarrow{a}}{\partial t'} \right)\left( \frac{\overrightarrow{R}}{K}\frac{\overrightarrow{v}}{c}+1 \right)-3\left( \overrightarrow{v}\frac{\overrightarrow{a}}{c} \right)\left( \frac{\overrightarrow{R}}{K}\frac{\overrightarrow{v}}{c}+1 \right) \right\}- \\ 
 & -\left( 1+\frac{\overrightarrow{R}\cdot \overrightarrow{a}-{{v}^{2}}}{{{c}^{2}}} \right)\left( 3\frac{\overrightarrow{R}}{{{K}^{4}}}\frac{\overrightarrow{v}}{c}-\frac{3R}{{{K}^{4}}}\left\{ \frac{\overrightarrow{R}}{K}\frac{\overrightarrow{v}}{c}\left( 1+\frac{\overrightarrow{R}\cdot \overrightarrow{a}-{{v}^{2}}}{{{c}^{2}}} \right)-\frac{\overrightarrow{v}}{c}\frac{\overrightarrow{v}}{c} \right\}-\frac{3R\overrightarrow{R}}{{{K}^{4}}{{c}^{2}}}\overrightarrow{a} \right)-2\frac{\overrightarrow{R}}{{{K}^{3}}}\frac{\overrightarrow{a}}{{{c}^{2}}}+\frac{2}{{{K}^{3}}}\frac{\overrightarrow{v}}{c}\frac{\overrightarrow{v}}{c} \\ 
\end{align}

{]} Для проверки сумма дивергенций обоих слагаемых электрического поля
{[}

\begin{align}
  & \frac{1}{e}\left( \frac{\partial {{E}_{x}}}{\partial x}+\frac{\partial {{E}_{y}}}{\partial y}+\frac{\partial {{E}_{z}}}{\partial z} \right)=-\frac{R}{{{K}^{3}}{{c}^{2}}}\frac{R}{K}\left( \frac{\overrightarrow{R}}{c}\frac{\partial \overrightarrow{a}}{\partial t'} \right)+\frac{R}{{{K}^{3}}{{c}^{2}}}\frac{R}{K}\left( \frac{3\ \overrightarrow{v}\cdot \overrightarrow{a}}{c} \right)-\frac{2}{{{K}^{3}}}+ \\ 
 & +\frac{R}{{{K}^{3}}}\frac{1}{{{c}^{2}}}\left\{ \left( \frac{\overrightarrow{R}}{c}\frac{\partial \overrightarrow{a}}{\partial t'} \right)\left( \frac{\overrightarrow{R}}{K}\frac{\overrightarrow{v}}{c}+1 \right)-3\left( \overrightarrow{v}\frac{\overrightarrow{a}}{c} \right)\left( \frac{\overrightarrow{R}}{K}\frac{\overrightarrow{v}}{c}+1 \right) \right\}+ \\ 
 & +\left( 1+\frac{\overrightarrow{a}\cdot \overrightarrow{R}}{{{c}^{2}}}-\frac{\overrightarrow{v}\cdot \overrightarrow{v}}{{{c}^{2}}} \right)\cdot \left\{ \frac{5}{{{K}^{3}}}-\frac{3{{R}^{2}}}{{{K}^{5}}}\left( 1+\frac{\overrightarrow{R}\cdot \overrightarrow{a}-{{v}^{2}}}{{{c}^{2}}} \right)+6\frac{\overrightarrow{v}\cdot \overrightarrow{R}}{c{{K}^{4}}} \right\}- \\ 
 & -\left( 1+\frac{\overrightarrow{R}\cdot \overrightarrow{a}-{{v}^{2}}}{{{c}^{2}}} \right)\left( 3\frac{\overrightarrow{R}}{{{K}^{4}}}\frac{\overrightarrow{v}}{c}-\frac{3R}{{{K}^{4}}}\left\{ \frac{\overrightarrow{R}}{K}\frac{\overrightarrow{v}}{c}\left( 1+\frac{\overrightarrow{R}\cdot \overrightarrow{a}-{{v}^{2}}}{{{c}^{2}}} \right)-\frac{\overrightarrow{v}}{c}\frac{\overrightarrow{v}}{c} \right\}-\frac{3R\overrightarrow{R}}{{{K}^{4}}{{c}^{2}}}\overrightarrow{a} \right)-2\frac{\overrightarrow{R}}{{{K}^{3}}}\frac{\overrightarrow{a}}{{{c}^{2}}}+\frac{2}{{{K}^{3}}}\frac{\overrightarrow{v}}{c}\frac{\overrightarrow{v}}{c} \\ 
\end{align}

{]}

{[}

\begin{align}
  & \frac{1}{e}\left( \frac{\partial {{E}_{x}}}{\partial x}+\frac{\partial {{E}_{y}}}{\partial y}+\frac{\partial {{E}_{z}}}{\partial z} \right)=-\frac{2}{{{K}^{3}}}-2\frac{\overrightarrow{R}}{{{K}^{3}}}\frac{\overrightarrow{a}}{{{c}^{2}}}+\frac{2}{{{K}^{3}}}\frac{\overrightarrow{v}}{c}\frac{\overrightarrow{v}}{c}+ \\ 
 & +\frac{R}{{{K}^{3}}}\frac{1}{{{c}^{2}}}\left\{ \left( \frac{\overrightarrow{R}}{c}\frac{\partial \overrightarrow{a}}{\partial t'} \right)\left( \frac{\overrightarrow{R}}{K}\frac{\overrightarrow{v}}{c}-\frac{R}{K}+1 \right)-3\left( \overrightarrow{v}\frac{\overrightarrow{a}}{c} \right)\left( \frac{\overrightarrow{R}}{K}\frac{\overrightarrow{v}}{c}-\frac{R}{K}+1 \right) \right\}+ \\ 
 & +\left( 1+\frac{\overrightarrow{a}\cdot \overrightarrow{R}}{{{c}^{2}}}-\frac{\overrightarrow{v}\cdot \overrightarrow{v}}{{{c}^{2}}} \right)\cdot \left\{ \frac{5}{{{K}^{3}}}-\frac{3{{R}^{2}}}{{{K}^{5}}}\left( 1+\frac{\overrightarrow{R}\cdot \overrightarrow{a}-{{v}^{2}}}{{{c}^{2}}} \right)+6\frac{\overrightarrow{v}\cdot \overrightarrow{R}}{c{{K}^{4}}}-3\frac{\overrightarrow{R}}{{{K}^{4}}}\frac{\overrightarrow{v}}{c}+\frac{3R}{{{K}^{4}}}\left\{ \frac{\overrightarrow{R}}{K}\frac{\overrightarrow{v}}{c}\left( 1+\frac{\overrightarrow{R}\cdot \overrightarrow{a}-{{v}^{2}}}{{{c}^{2}}} \right)-\frac{\overrightarrow{v}}{c}\frac{\overrightarrow{v}}{c} \right\}+\frac{3R\overrightarrow{R}}{{{K}^{4}}{{c}^{2}}}\overrightarrow{a} \right\} \\ 
 &  \\ 
\end{align}

{]} {[}\left\{
\frac{1}{c}\frac{\overrightarrow{R}\overrightarrow{v}}{K}-\frac{R}{K}+1
\right\}=0{]} {[}

\begin{align}
  & \frac{1}{e}\left( \frac{\partial {{E}_{x}}}{\partial x}+\frac{\partial {{E}_{y}}}{\partial y}+\frac{\partial {{E}_{z}}}{\partial z} \right)=-\frac{2}{{{K}^{3}}}\left( 1+\frac{\overrightarrow{a}\cdot \overrightarrow{R}}{{{c}^{2}}}-\frac{\overrightarrow{v}\cdot \overrightarrow{v}}{{{c}^{2}}} \right)+ \\ 
 & +\left( 1+\frac{\overrightarrow{a}\cdot \overrightarrow{R}}{{{c}^{2}}}-\frac{\overrightarrow{v}\cdot \overrightarrow{v}}{{{c}^{2}}} \right)\frac{1}{{{K}^{3}}}\left\{ 5-\frac{3{{R}^{2}}}{{{K}^{2}}}\left( 1+\frac{\overrightarrow{R}\cdot \overrightarrow{a}-{{v}^{2}}}{{{c}^{2}}} \right)+3\frac{\overrightarrow{v}\cdot \overrightarrow{R}}{cK}+\frac{3R}{K}\left\{ \frac{\overrightarrow{R}}{K}\frac{\overrightarrow{v}}{c}\left( 1+\frac{\overrightarrow{R}\cdot \overrightarrow{a}-{{v}^{2}}}{{{c}^{2}}} \right)-\frac{\overrightarrow{v}}{c}\frac{\overrightarrow{v}}{c} \right\}+\frac{3R\overrightarrow{R}}{K{{c}^{2}}}\overrightarrow{a} \right\} \\ 
\end{align}

{]}

{[}\frac{1}{e}\left(
\frac{\partial {{E}_{x}}}{\partial x}+\frac{\partial {{E}_{y}}}{\partial y}+\frac{\partial {{E}_{z}}}{\partial z}
\right)=\left(
1+\frac{\overrightarrow{a}\cdot \overrightarrow{R}}{{{c}^{2}}}-\frac{\overrightarrow{v}\cdot \overrightarrow{v}}{{{c}^{2}}}
\right)\frac{1}{{{K}^{3}}}\left\{ 3-\frac{3{{R}^{2}}}{{{K}^{2}}}\left(
1+\frac{\overrightarrow{R}\cdot \overrightarrow{a}-{{v}^{2}}}{{{c}^{2}}}
\right)+3\frac{\overrightarrow{v}\cdot \overrightarrow{R}}{cK}+\frac{3R}{K}\left\{
\frac{\overrightarrow{R}}{K}\frac{\overrightarrow{v}}{c}\left(
1+\frac{\overrightarrow{R}\cdot \overrightarrow{a}-{{v}^{2}}}{{{c}^{2}}}
\right)-\frac{\overrightarrow{v}}{c}\frac{\overrightarrow{v}}{c}
\right\}+\frac{3R\overrightarrow{R}}{K{{c}^{2}}}\overrightarrow{a}
\right\}{]}

{[}\frac{1}{e}\left(
\frac{\partial {{E}_{x}}}{\partial x}+\frac{\partial {{E}_{y}}}{\partial y}+\frac{\partial {{E}_{z}}}{\partial z}
\right)=\left(
1+\frac{\overrightarrow{a}\cdot \overrightarrow{R}}{{{c}^{2}}}-\frac{\overrightarrow{v}\cdot \overrightarrow{v}}{{{c}^{2}}}
\right)\frac{1}{{{K}^{3}}}\left\{ 3-\frac{3{{R}^{2}}}{{{K}^{2}}}\left(
1+\frac{\overrightarrow{R}\cdot \overrightarrow{a}-{{v}^{2}}}{{{c}^{2}}}
\right)+3\frac{\overrightarrow{v}\cdot \overrightarrow{R}}{cK}+\frac{3R}{K}\frac{\overrightarrow{R}}{K}\frac{\overrightarrow{v}}{c}\left(
1+\frac{\overrightarrow{R}\cdot \overrightarrow{a}-{{v}^{2}}}{{{c}^{2}}}
\right)-\frac{3R}{K}\frac{\overrightarrow{v}}{c}\frac{\overrightarrow{v}}{c}+3\frac{R}{K}\frac{\overrightarrow{R}\cdot \overrightarrow{a}}{{{c}^{2}}}
\right\}{]} {[}\frac{1}{e}\left(
\frac{\partial {{E}_{x}}}{\partial x}+\frac{\partial {{E}_{y}}}{\partial y}+\frac{\partial {{E}_{z}}}{\partial z}
\right)=\left(
1+\frac{\overrightarrow{a}\cdot \overrightarrow{R}}{{{c}^{2}}}-\frac{\overrightarrow{v}\cdot \overrightarrow{v}}{{{c}^{2}}}
\right)\frac{1}{{{K}^{3}}}\left\{ 3-\frac{3R}{K}\frac{R}{K}\left(
1+\frac{\overrightarrow{R}\cdot \overrightarrow{a}-{{v}^{2}}}{{{c}^{2}}}
\right)+\frac{3R}{K}\frac{\overrightarrow{R}}{K}\frac{\overrightarrow{v}}{c}\left(
1+\frac{\overrightarrow{R}\cdot \overrightarrow{a}-{{v}^{2}}}{{{c}^{2}}}
\right)+3\frac{\overrightarrow{v}\cdot \overrightarrow{R}}{cK}-\frac{3R}{K}\frac{\overrightarrow{v}}{c}\frac{\overrightarrow{v}}{c}+3\frac{R}{K}\frac{\overrightarrow{R}\cdot \overrightarrow{a}}{{{c}^{2}}}
\right\}{]} {[}\frac{1}{e}\left(
\frac{\partial {{E}_{x}}}{\partial x}+\frac{\partial {{E}_{y}}}{\partial y}+\frac{\partial {{E}_{z}}}{\partial z}
\right)=\left(
1+\frac{\overrightarrow{a}\cdot \overrightarrow{R}}{{{c}^{2}}}-\frac{\overrightarrow{v}\cdot \overrightarrow{v}}{{{c}^{2}}}
\right)\frac{1}{{{K}^{3}}}\left\{ 3+\frac{3R}{K}\left(
\frac{\overrightarrow{R}}{K}\frac{\overrightarrow{v}}{c}-\frac{R}{K}
\right)\left(
1+\frac{\overrightarrow{R}\cdot \overrightarrow{a}-{{v}^{2}}}{{{c}^{2}}}
\right)+3\frac{\overrightarrow{v}\cdot \overrightarrow{R}}{cK}-\frac{3R}{K}\frac{\overrightarrow{v}}{c}\frac{\overrightarrow{v}}{c}+3\frac{R}{K}\frac{\overrightarrow{R}\cdot \overrightarrow{a}}{{{c}^{2}}}
\right\}{]} {[}\left\{ \frac{1}{K}\left(
\frac{\overrightarrow{v}\cdot \overrightarrow{R}}{c} \right)-\frac{R}{K}
\right\}=-1{]} {[}\frac{1}{e}\left(
\frac{\partial {{E}_{x}}}{\partial x}+\frac{\partial {{E}_{y}}}{\partial y}+\frac{\partial {{E}_{z}}}{\partial z}
\right)=\left(
1+\frac{\overrightarrow{a}\cdot \overrightarrow{R}}{{{c}^{2}}}-\frac{\overrightarrow{v}\cdot \overrightarrow{v}}{{{c}^{2}}}
\right)\frac{1}{{{K}^{3}}}\left\{ 3-\frac{3R}{K}\left(
1+\frac{\overrightarrow{R}\cdot \overrightarrow{a}-{{v}^{2}}}{{{c}^{2}}}
\right)+3\frac{\overrightarrow{v}\cdot \overrightarrow{R}}{cK}-\frac{3R}{K}\frac{\overrightarrow{v}}{c}\frac{\overrightarrow{v}}{c}+3\frac{R}{K}\frac{\overrightarrow{R}\cdot \overrightarrow{a}}{{{c}^{2}}}
\right\}{]}

{[}\frac{1}{e}\left(
\frac{\partial {{E}_{x}}}{\partial x}+\frac{\partial {{E}_{y}}}{\partial y}+\frac{\partial {{E}_{z}}}{\partial z}
\right)=\frac{1}{{{K}^{3}}}\left(
1+\frac{\overrightarrow{a}\cdot \overrightarrow{R}}{{{c}^{2}}}-\frac{\overrightarrow{v}\cdot \overrightarrow{v}}{{{c}^{2}}}
\right)\left[ 3+3\frac{R}{K}\frac{\overrightarrow{a}\cdot \overrightarrow{R}}{{{c}^{2}}}-\frac{3R}{K}\frac{{{v}^{2}}}{{{c}^{2}}}-3\frac{R}{K}\left( 1+\frac{\overrightarrow{R}\cdot \overrightarrow{a}-{{v}^{2}}}{{{c}^{2}}} \right)+\frac{3}{K}\left( \frac{\overrightarrow{v}\cdot \overrightarrow{R}}{c} \right) \right]{]}

Сокращаем

{[}\frac{1}{e}\left(
\frac{\partial {{E}_{x}}}{\partial x}+\frac{\partial {{E}_{y}}}{\partial y}+\frac{\partial {{E}_{z}}}{\partial z}
\right)=\frac{3}{{{K}^{3}}}\left(
1+\frac{\overrightarrow{a}\cdot \overrightarrow{R}}{{{c}^{2}}}-\frac{\overrightarrow{v}\cdot \overrightarrow{v}}{{{c}^{2}}}
\right)\left[ 1+\frac{R}{K}\frac{\overrightarrow{a}\cdot \overrightarrow{R}}{{{c}^{2}}}-\frac{R}{K}\frac{{{v}^{2}}}{{{c}^{2}}}-\frac{R}{K}\left( 1+\frac{\overrightarrow{R}\cdot \overrightarrow{a}-{{v}^{2}}}{{{c}^{2}}} \right)+\frac{1}{K}\left( \frac{\overrightarrow{v}\cdot \overrightarrow{R}}{c} \right) \right]{]}

{[}\frac{1}{e}\left(
\frac{\partial {{E}_{x}}}{\partial x}+\frac{\partial {{E}_{y}}}{\partial y}+\frac{\partial {{E}_{z}}}{\partial z}
\right)=\frac{3}{{{K}^{3}}}\left(
1+\frac{\overrightarrow{a}\cdot \overrightarrow{R}}{{{c}^{2}}}-\frac{\overrightarrow{v}\cdot \overrightarrow{v}}{{{c}^{2}}}
\right)\left[ 1+\frac{R}{K}\frac{\overrightarrow{a}\cdot \overrightarrow{R}}{{{c}^{2}}}-\frac{R}{K}\frac{{{v}^{2}}}{{{c}^{2}}}-\frac{R}{K}-\frac{R}{K}\frac{\overrightarrow{R}\cdot \overrightarrow{a}-{{v}^{2}}}{{{c}^{2}}}+\frac{1}{K}\left( \frac{\overrightarrow{v}\cdot \overrightarrow{R}}{c} \right) \right]{]}

{[}\frac{1}{e}\left(
\frac{\partial {{E}_{x}}}{\partial x}+\frac{\partial {{E}_{y}}}{\partial y}+\frac{\partial {{E}_{z}}}{\partial z}
\right)=\frac{3}{{{K}^{3}}}\left(
1+\frac{\overrightarrow{a}\cdot \overrightarrow{R}}{{{c}^{2}}}-\frac{\overrightarrow{v}\cdot \overrightarrow{v}}{{{c}^{2}}}
\right)\left[ 1+\frac{R}{K}\frac{\overrightarrow{a}\cdot \overrightarrow{R}}{{{c}^{2}}}-\frac{R}{K}\frac{{{v}^{2}}}{{{c}^{2}}}-\frac{R}{K}\frac{\overrightarrow{R}\cdot \overrightarrow{a}-{{v}^{2}}}{{{c}^{2}}}+\frac{1}{K}\left( \frac{\overrightarrow{v}\cdot \overrightarrow{R}}{c} \right)-\frac{R}{K} \right]{]}
{[}\left\{ \frac{1}{K}\left(
\frac{\overrightarrow{v}\cdot \overrightarrow{R}}{c} \right)-\frac{R}{K}
\right\}=-1{]}

{[}\frac{1}{e}\left(
\frac{\partial {{E}_{x}}}{\partial x}+\frac{\partial {{E}_{y}}}{\partial y}+\frac{\partial {{E}_{z}}}{\partial z}
\right)=\frac{3}{{{K}^{3}}}\left(
1+\frac{\overrightarrow{a}\cdot \overrightarrow{R}}{{{c}^{2}}}-\frac{\overrightarrow{v}\cdot \overrightarrow{v}}{{{c}^{2}}}
\right)\left[ \frac{R}{K}\frac{\overrightarrow{a}\cdot \overrightarrow{R}}{{{c}^{2}}}-\frac{R}{K}\frac{{{v}^{2}}}{{{c}^{2}}}-\frac{R}{K}\frac{\overrightarrow{R}\cdot \overrightarrow{a}-{{v}^{2}}}{{{c}^{2}}} \right]{]}

{[}\frac{1}{e}\left(
\frac{\partial {{E}_{x}}}{\partial x}+\frac{\partial {{E}_{y}}}{\partial y}+\frac{\partial {{E}_{z}}}{\partial z}
\right)=\frac{3}{{{K}^{3}}}\frac{R}{K}\left(
1+\frac{\overrightarrow{a}\cdot \overrightarrow{R}}{{{c}^{2}}}-\frac{\overrightarrow{v}\cdot \overrightarrow{v}}{{{c}^{2}}}
\right)\left[ \frac{\overrightarrow{a}\cdot \overrightarrow{R}}{{{c}^{2}}}-\frac{{{v}^{2}}}{{{c}^{2}}}-\frac{\overrightarrow{R}\cdot \overrightarrow{a}-{{v}^{2}}}{{{c}^{2}}} \right]{]}

В обозначениях программы Дробышева \(k=\frac{K}{R}\) \(K=k\cdot R\)
\(c=1\) {[}\frac{1}{e}\left(
\frac{\partial {{E}_{x}}}{\partial x}+\frac{\partial {{E}_{y}}}{\partial y}+\frac{\partial {{E}_{z}}}{\partial z}
\right)=\frac{1}{{{\left( k\cdot R \right)}^{3}}}\left(
1+\overrightarrow{a}\cdot \overrightarrow{R}-\overrightarrow{v}\cdot \overrightarrow{v}
\right)\left[ 3+\frac{3}{k}\overrightarrow{a}\cdot \overrightarrow{R}-\frac{3}{k}{{v}^{2}}-\frac{3}{k}\left( 1+\overrightarrow{R}\cdot \overrightarrow{a}-{{v}^{2}} \right)+\frac{3}{k\cdot R}\left( \overrightarrow{v}\cdot \overrightarrow{R} \right) \right]{]}
{[}\frac{1}{e}\left(
\frac{\partial {{E}_{x}}}{\partial x}+\frac{\partial {{E}_{y}}}{\partial y}+\frac{\partial {{E}_{z}}}{\partial z}
\right)=\frac{3}{{{\left( k\cdot R \right)}^{3}}}\left(
1+\overrightarrow{a}\cdot \overrightarrow{R}-\overrightarrow{v}\cdot \overrightarrow{v}
\right)\left[ 1+\frac{1}{k}\left( \overrightarrow{a}\cdot \overrightarrow{R}-{{v}^{2}} \right)-\frac{1}{k}\left( 1+\overrightarrow{R}\cdot \overrightarrow{a}-{{v}^{2}} \right)+\frac{1}{k\cdot R}\left( \overrightarrow{v}\cdot \overrightarrow{R} \right) \right]{]}
{[}\frac{1}{e}\left(
\frac{\partial {{E}_{x}}}{\partial x}+\frac{\partial {{E}_{y}}}{\partial y}+\frac{\partial {{E}_{z}}}{\partial z}
\right)=\frac{3}{{{\left( k\cdot R \right)}^{3}}}\left(
1+\overrightarrow{a}\cdot \overrightarrow{R}-\overrightarrow{v}\cdot \overrightarrow{v}
\right)\left[ 1-\frac{1}{k}+\frac{1}{k\cdot R}\left( \overrightarrow{v}\cdot \overrightarrow{R} \right) \right]{]}
Таким образом возникновение потенциальных ям скалярного потенциала

действительно не связано с дополнительными зарядами, материализующимися
"из воздуха", и уравнение \(\nabla\cdot\vec{E}=0\) для вакуума
выполняется. Однако для самого точечного заряда выражение дивергенции
электрического поля приходит к неопределенности вида ноль делить на
ноль. Поэтому представляет интерес произвести вычисления дивергенции
поля для неточечного заряда, имеющего заданное распределение объёмной
плотности заряда в пространстве. Интересно показать распределение
отношения радиуса Лиенара Вихерта к запаздывающему радиусу для той же
задачи

Расчёт суммарного электрического поля в тот же момент времени в целом
повторяет результат, полученный Дробышевым и Меандром в той же теме

Однако не были показаны слагаемые выражения электрического поля через
потенциалы. Первое {[}\overrightarrow{{{E}_{1}}}=-\nabla \varphi {]}

И второе
{[}\overrightarrow{{{E}_{2}}}=-\frac{1}{c}\frac{\partial \overrightarrow{A}}{\partial t}{]}

Дивергенция кулоновского поля покоящегося точечного заряда Возникает
вопрос: из электродинамики известно, что дивергенция электрического поля
покоящегося в вакууме заряда равна нулю. Как же получается, что при его
движении дивергенция поля в вакууме становится отличной от нуля? Кроме
того, почему на графике дивергенции на далёких расстояниях появились
концентрические окружности, указывающие на убывание дивергенции поля с
расстоянием от заряда? Электрическое поле в электростатике находится как
минус градиент скалярного потенциала
{[}\overrightarrow{E}=-\nabla \varphi =-\nabla \frac{e}{R}=\frac{e}{{{R}^{2}}}\frac{\overrightarrow{R}}{R}{]}
Найдем производную электрического поля по координате точки наблюдения
{[}\frac{1}{e}\frac{\partial \overrightarrow{E}}{\partial x}=\frac{\partial }{\partial x}\left[ \frac{\overrightarrow{R}}{{{R}^{3}}} \right]=\frac{\partial }{\partial x}\left(
\frac{\overrightarrow{R}}{{{R}^{3}}} \right){]}
{[}\frac{\partial }{\partial x}\left(
\frac{\overrightarrow{R}}{{{R}^{3}}}
\right)=\frac{1}{{{R}^{3}}}\frac{\partial \overrightarrow{R}}{\partial x}+\overrightarrow{R}\frac{\partial }{\partial x}\left(
\frac{1}{{{R}^{3}}}
\right)=\frac{1}{{{R}^{3}}}\frac{\partial \overrightarrow{R}}{\partial x}-\overrightarrow{R}\frac{3}{{{R}^{4}}}\frac{\partial R}{\partial x}{]}
Производная \(R\) по координате точки наблюдения
{[}\frac{\partial R}{\partial x}=\frac{\partial }{\partial x}{{\left( \overrightarrow{R}\cdot \overrightarrow{R} \right)}^{{}^{1}/{}_{2}}}=\frac{1}{2}{{\left( \overrightarrow{R}\cdot \overrightarrow{R} \right)}^{-{}^{1}/{}_{2}}}\frac{\partial }{\partial x}\left(
\overrightarrow{R}\cdot \overrightarrow{R} \right)=\frac{1}{2R}\left(
\overrightarrow{R}\cdot \frac{\partial \overrightarrow{R}}{\partial x}+\frac{\partial \overrightarrow{R}}{\partial x}\cdot \overrightarrow{R}
\right)=\frac{\overrightarrow{R}}{R}\cdot \frac{\partial \overrightarrow{R}}{\partial x}{]}
Поскольку в декартовой системе координат
{[}\overrightarrow{R}=\overrightarrow{i}~\{\{R\}\emph{\{x\}\}+\overrightarrow{j}~\{\{R\}}\{y\}\}+\overrightarrow{k}~\{\{R\}\emph{\{z\}\}=\overrightarrow{i}\left(
x-\{\{x\}}\{q\}\} \right)+\overrightarrow{j}\left(
y-\{\{y\}\emph{\{q\}\} \right)+\overrightarrow{k}\left(
z-\{\{z\}}\{q\}\} \right){]} постольку
{[}\frac{\partial \overrightarrow{R}}{\partial x}=

\begin{matrix}
   \overrightarrow{i}\ \frac{\partial {{R}_{x}}}{\partial x}  \\
   \overrightarrow{j}\ \frac{\partial {{R}_{y}}}{\partial x}  \\
   \overrightarrow{k}\ \frac{\partial {{R}_{z}}}{\partial x}  \\
\end{matrix}

=

\begin{matrix}
   \overrightarrow{i}\ \frac{\partial \left( x-{{x}_{q}} \right)}{\partial x}  \\
   \overrightarrow{j}\ \frac{\partial \left( y-{{y}_{q}} \right)}{\partial x}  \\
   \overrightarrow{k}\ \frac{\partial \left( z-{{z}_{q}} \right)}{\partial x}  \\
\end{matrix}

=\overrightarrow{i}{]} и
{[}\frac{\partial R}{\partial x}=\frac{\overrightarrow{R}}{R}\cdot \frac{\partial \overrightarrow{R}}{\partial x}=\frac{\overrightarrow{R}}{R}\cdot \overrightarrow{i}=\frac{{{R}_{x}}}{R}{]}
Теперь
{[}\frac{1}{e}\frac{\partial \overrightarrow{E}}{\partial x}=\frac{\partial }{\partial x}\left(
\frac{\overrightarrow{R}}{{{R}^{3}}}
\right)=\frac{1}{{{R}^{3}}}\frac{\partial \overrightarrow{R}}{\partial x}-\overrightarrow{R}\frac{3}{{{R}^{4}}}\frac{\partial R}{\partial x}=\frac{\overrightarrow{i}}{{{R}^{3}}}-\overrightarrow{R}\frac{3}{{{R}^{4}}}\frac{{{R}_{x}}}{R}{]}
Производная компоненты поля по координате точки наблюдения
{[}\frac{1}{e}\frac{\partial {{E}_{x}}}{\partial x}=\frac{\partial }{\partial x}\left(
\frac{{{R}_{x}}}{{{R}^{3}}}
\right)=\frac{1}{{{R}^{3}}}\frac{\partial {{R}_{x}}}{\partial x}-\{\{R\}\emph{\{x\}\}\frac{3}{{{R}^{4}}}\frac{\partial R}{\partial x}=\frac{1}{{{R}^{3}}}-\{\{R\}}\{x\}\}\frac{3}{{{R}^{4}}}\frac{{{R}_{x}}}{R}=\frac{1}{{{R}^{3}}}\left(
1-3\frac{{{R}_{x}}^{2}}{{{R}^{2}}} \right){]} Теперь полезно сравнить
выражение, полученное для производной компоненты поля исходя из
потенциалов ЛВ {[}

\begin{align}
  & \frac{1}{e}\frac{\partial {{E}_{x}}}{\partial x}=\frac{1}{{{c}^{2}}}\frac{1}{{{K}^{3}}}\left[ -\left( \frac{\overrightarrow{R}}{c}\frac{\partial \overrightarrow{a}}{\partial t'} \right)\frac{{{R}_{x}}{{R}_{x}}}{K}+3\left( \overrightarrow{v}\frac{\overrightarrow{a}}{c} \right)\frac{{{R}_{x}}{{R}_{x}}}{K}+\left( \frac{\overrightarrow{R}}{c}\frac{\partial \overrightarrow{a}}{\partial t'} \right)\frac{R}{c}\frac{{{R}_{x}}{{v}_{x}}}{K}-\frac{3R}{c}{{a}_{x}}{{v}_{x}}-3\frac{R}{c}\left( \overrightarrow{v}\frac{\overrightarrow{a}}{c} \right)\frac{{{R}_{x}}{{v}_{x}}}{K}+\frac{R{{R}_{x}}}{c}\frac{\partial {{a}_{x}}}{\partial t'} \right]+ \\ 
 & +\left( 1+\frac{\overrightarrow{a}\cdot \overrightarrow{R}}{{{c}^{2}}}-\frac{\overrightarrow{v}\cdot \overrightarrow{v}}{{{c}^{2}}} \right)\frac{1}{{{K}^{3}}}\left[ 1+\frac{2R}{{{c}^{2}}}\frac{{{a}_{x}}{{R}_{x}}}{K}+\frac{3R}{K}\left\{ \frac{{{v}_{x}}}{c}\frac{{{R}_{x}}}{K}\left( 1+\frac{\overrightarrow{R}\cdot \overrightarrow{a}-{{v}^{2}}}{{{c}^{2}}} \right)-\frac{{{v}_{x}}}{c}\frac{{{v}_{x}}}{c} \right\}-\frac{3}{K}\left\{ \frac{{{R}_{x}}{{R}_{x}}}{K}\left( 1+\frac{\overrightarrow{R}\cdot \overrightarrow{a}-{{v}^{2}}}{{{c}^{2}}} \right)-\frac{{{v}_{x}}{{R}_{x}}}{c} \right\}+\frac{R{{R}_{x}}}{K}\frac{{{a}_{x}}}{{{c}^{2}}} \right] \\ 
\end{align}

{]} Подставив в него нулевые скорость и ускорение
{[}\frac{1}{e}\frac{\partial {{E}_{x}}}{\partial x}=\frac{1}{{{K}^{3}}}\left[ 1-\frac{3}{K}\left\{ \frac{{{R}_{x}}{{R}_{x}}}{K} \right\} \right]{]}
Дивергенция {[}\frac{1}{e}\left(
\frac{\partial {{E}_{x}}}{\partial x}+\frac{\partial {{E}_{y}}}{\partial y}+\frac{\partial {{E}_{z}}}{\partial z}
\right)=\frac{1}{{{R}^{3}}}\left\{ \left(
1-3\frac{{{R}_{x}}^{2}}{{{R}^{2}}} \right)+\left(
1-3\frac{{{R}_{y}}^{2}}{{{R}^{2}}} \right)+\left(
1-3\frac{{{R}_{z}}^{2}}{{{R}^{2}}} \right)
\right\}=\frac{1}{{{R}^{3}}}\left\{ 3-3\frac{{{R}^{2}}}{{{R}^{2}}}
\right\}=0{]} Все ли в порядке с теоремой Гаусса?

Сумма двух компонент электрического поля
{[}\overrightarrow{E}=-\nabla \varphi -\frac{1}{c}\frac{\partial \overrightarrow{A}}{\partial t}=\frac{e}{{{K}^{3}}}\left[ \left( \overrightarrow{R}-R\frac{\overrightarrow{v}}{c} \right)\left( 1+\frac{\overrightarrow{a}\cdot \overrightarrow{R}}{{{c}^{2}}}-\frac{\overrightarrow{v}\cdot \overrightarrow{v}}{{{c}^{2}}} \right)-KR\frac{\overrightarrow{a}}{{{c}^{2}}} \right]{]}
{[}\overrightarrow{E}=\frac{e}{{{K}^{3}}}\left(
\overrightarrow{R}-R\frac{\overrightarrow{v}}{c} \right)\left(
1+\frac{\overrightarrow{a}\cdot \overrightarrow{R}}{{{c}^{2}}}-\frac{\overrightarrow{v}\cdot \overrightarrow{v}}{{{c}^{2}}}
\right)-\frac{e}{{{K}^{2}}}R\frac{\overrightarrow{a}}{{{c}^{2}}}{]}
Пусть заряд в запаздывающий момент времени \({{t}_{0}}'\) находится в
начале координат. Вообразим сферу радиуса \({{R}_{0}}\) с центром в
начале координат и вычислим нормальную компоненту электрического поля в
момент {[}\{\{t\}\emph{\{0\}\}=\{\{t\}}\{0\}\}'+\frac{{{R}_{0}}}{c}{]}
на поверхности этой сферы. Для этого вектор \(\overrightarrow{E}\)
необходимо скалярно умножить на вектор нормали
\(\overrightarrow{n}=\frac{\overrightarrow{{{R}_{0}}}}{{{R}_{0}}}\) ,
где \(\overrightarrow{{{R}_{0}}}\) есть вектор проведённый из начала
координат в точку поверхности сферы. При данных условиях, поскольку для
всей поверхности сферы запаздывающий момент одинаков и заряд в этот
запаздывающий момент находится в начале координат, то \(R={{R}_{0}}\) и
\(\overrightarrow{R}=\overrightarrow{{{R}_{0}}}\). Отсюда
\({{E}_{n}}=\overrightarrow{E}\frac{\overrightarrow{R}}{R}=\frac{e}{{{K}^{3}}}\left( \overrightarrow{R}\frac{\overrightarrow{R}}{R}-R\frac{\overrightarrow{v}}{c}\frac{\overrightarrow{R}}{R} \right)\left( 1+\frac{\overrightarrow{a}\cdot \overrightarrow{R}}{{{c}^{2}}}-\frac{\overrightarrow{v}\cdot \overrightarrow{v}}{{{c}^{2}}} \right)-\frac{e}{{{K}^{2}}}R\frac{\overrightarrow{a}}{{{c}^{2}}}\frac{\overrightarrow{R}}{R}\)
\({{E}_{n}}=\frac{e}{{{K}^{3}}}\left( R-\frac{\overrightarrow{v}\cdot \overrightarrow{R}}{c} \right)\left( 1+\frac{\overrightarrow{a}\cdot \overrightarrow{R}}{{{c}^{2}}}-\frac{\overrightarrow{v}\cdot \overrightarrow{v}}{{{c}^{2}}} \right)-\frac{e}{{{K}^{2}}}\frac{\overrightarrow{a}\cdot \overrightarrow{R}}{{{c}^{2}}}\)
\({{E}_{n}}=\frac{e}{{{K}^{2}}}\left( 1+\frac{\overrightarrow{a}\cdot \overrightarrow{R}}{{{c}^{2}}}-\frac{\overrightarrow{v}\cdot \overrightarrow{v}}{{{c}^{2}}} \right)-\frac{e}{{{K}^{2}}}\frac{\overrightarrow{a}\cdot \overrightarrow{R}}{{{c}^{2}}}\)
\({{E}_{n}}=\frac{e}{{{K}^{2}}}\left( 1-\frac{\overrightarrow{v}\cdot \overrightarrow{v}}{{{c}^{2}}} \right)=\frac{e}{{{\left( R-\frac{\overrightarrow{v}\cdot \overrightarrow{R}}{c} \right)}^{2}}}\left( 1-\frac{\overrightarrow{v}\cdot \overrightarrow{v}}{{{c}^{2}}} \right)\)
Интересно, что ускорение заряда не входит в формулу для нормальной
компоненты электрического поля. Нормальная компонента электрического
поля зависит только от скорости заряда в момент \({{t}_{0}}'\). Выберем
направление скорости заряда в этот момент в качестве оси и введём
сферическую систему координат. В этой системе
\(\overrightarrow{v}\cdot \overrightarrow{R}=vR\cos \left( \theta \right)\)
\({{E}_{n}}=e\frac{\left( 1-\frac{{{v}^{2}}}{{{c}^{2}}} \right)}{{{\left( R-\frac{vR\cos \left( \theta \right)}{c} \right)}^{2}}}=\frac{e}{{{R}^{2}}}\frac{\left( 1-\frac{{{v}^{2}}}{{{c}^{2}}} \right)}{{{\left( 1-\frac{v}{c}\cos \left( \theta \right) \right)}^{2}}}\)
Интеграл по поверхности сферы {[}\iint{{{E}_{n}}{{R}^{2}}}\sin \left(
\theta  \right)d\theta d\varphi =e\int\limits\emph{\{\theta =0\}\textsuperscript{\{\pi \}\{\int\limits\emph{\{\varphi =0\}\^{}\{2\pi \}\{\frac{\left( 1-\frac{{{v}^{2}}}{{{c}^{2}}} \right)}{{{\left( 1-\frac{v}{c}\cos \left( \theta  \right) \right)}^{2}}}\}\}\sin \left(
\theta  \right)d\theta d\varphi =2\pi e\int\limits}\{\theta =0\}}\{\pi \}\{\frac{\left( 1-\frac{{{v}^{2}}}{{{c}^{2}}} \right)}{{{\left( 1-\frac{v}{c}\cos \left( \theta  \right) \right)}^{2}}}\}\sin \left(
\theta  \right)d\theta =4\pi e{]} В данном рассмотренном выше случае
соответствует теореме Гаусса. Но возникает вопрос, будет соответствовать
ли теореме Гаусса этот интеграл по той же поверхности, но в другие
моменты времени? Для простоты рассмотрим практически важный случай,
когда векторы ускорения и скорости заряда сонаправлены. Уравнение
движения заряда для этого случая
{[}\{\{z\}}\{q\}\}=\{\{z\}\emph{\{0\}\}+\{\{v\}}\{0\}\}\left(
t'-\{\{t\}\emph{\{0\}\}'
\right)+\frac{{{a}_{0}}{{\left( t'-{{t}_{0}}' \right)}^{2}}}{2}{]}
{[}\{\{v\}}\{q\}\}=\{\{v\}\emph{\{0\}\}+\{\{a\}}\{0\}\}\left(
t'-\{\{t\}\emph{\{0\}\}' \right){]} Расчёт запаздывающего момента \(t'\)
производится с помощью разрешения уравнения вида (это уравнение записано
с помощью теоремы косинусов)
{[}\{\{R\}\textsuperscript{\{2\}\}=\{\{c\}}\{2\}\}\{\{\left( t-t'
\right)\}\^{}\{2\}\}=\{\{R\}}\{0\}\}\textsuperscript{\{2\}-2\{\{R\}\emph{\{0\}\}\{\{z\}}\{q\}\}\cos \theta +\{\{z\}\emph{\{q\}\}\^{}\{2\}{]}
Нормальная компонента электрического поля
\({{E}_{n}}=\overrightarrow{E}\frac{\overrightarrow{{{R}_{0}}}}{{{R}_{0}}}=\frac{e}{{{K}^{3}}}\left( \overrightarrow{R}\frac{\overrightarrow{{{R}_{0}}}}{{{R}_{0}}}-R\frac{\overrightarrow{v}}{c}\frac{\overrightarrow{{{R}_{0}}}}{{{R}_{0}}} \right)\left( 1+\frac{\overrightarrow{a}\cdot \overrightarrow{R}}{{{c}^{2}}}-\frac{\overrightarrow{v}\cdot \overrightarrow{v}}{{{c}^{2}}} \right)-\frac{e}{{{K}^{2}}}R\frac{\overrightarrow{a}}{{{c}^{2}}}\frac{\overrightarrow{{{R}_{0}}}}{{{R}_{0}}}\)
Для упрощения аналитического вывода формулы запаздывающего момента можно
применить два приближения: приближение малых скоростей \(v\to 0\) и
приближение малых ускорений \(a\to 0\) В приближении малых скоростей при
\({{z}_{0}}=0\) уравнение для вычисления запаздывающего момента
принимает вид {[}\{\{c\}\^{}\{2\}\}\{\{\left( t-t'
\right)\}\^{}\{2\}\}=\{\{R\}}\{0\}\}}\{2\}-2\{\{R\}\emph{\{0\}\}\frac{{{a}_{0}}{{\left( t'-{{t}_{0}}' \right)}^{2}}}{2}\cos \theta +\{\{\left(
\frac{{{a}_{0}}{{\left( t'-{{t}_{0}}' \right)}^{2}}}{2}
\right)\}\^{}\{2\}\}{]} А в приближении малых ускорений при
\({{z}_{0}}=0\) {[}\{\{c\}\^{}\{2\}\}\{\{\left( t-t'
\right)\}\^{}\{2\}\}=\{\{R\}}\{0\}\}\^{}\{2\}-2\{\{R\}\emph{\{0\}\}\{\{v\}}\{0\}\}\left(
t'-\{\{t\}\emph{\{0\}\}' \right)\cos \theta +\{\{\left(
\{\{v\}}\{0\}\}\left( t'-\{\{t\}\emph{\{0\}\}' \right)
\right)\}\^{}\{2\}\}{]} В приближении малых ускорений аналитически
решать все же проще
{[}t'=\frac{-{{R}_{0}}{{v}_{0}}cos\left( \theta  \right)+{{c}^{2}}t-{{t}_{0}}'{{v}_{0}}^{2}-\sqrt{cos{{\left( \theta  \right)}^{2}}{{R}_{0}}^{2}{{v}_{0}}^{2}-2cos\left( \theta  \right){{R}_{0}}{{c}^{2}}t{{v}_{0}}+2cos\left( \theta  \right){{R}_{0}}{{c}^{2}}{{t}_{0}}'{{v}_{0}}+{{c}^{2}}{{t}^{2}}{{v}_{0}}^{2}-2{{c}^{2}}t\cdot {{t}_{0}}'{{v}_{0}}^{2}+{{c}^{2}}{{\left( {{t}_{0}}' \right)}^{2}}{{v}_{0}}^{2}+{{R}_{0}}^{2}{{c}^{2}}-{{R}_{0}}^{2}{{v}_{0}}^{2}}}{({{c}^{2}}-{{v}_{0}}^{2})}{]}
{[}t'=\frac{-{{R}_{0}}{{v}_{0}}cos\left( \theta  \right)+{{c}^{2}}t-{{t}_{0}}'{{v}_{0}}^{2}-\sqrt{cos{{\left( \theta  \right)}^{2}}{{R}_{0}}^{2}{{v}_{0}}^{2}-2cos\left( \theta  \right){{R}_{0}}{{v}_{0}}\left( t-{{t}_{0}}' \right){{c}^{2}}+{{\left( ct{{v}_{0}} \right)}^{2}}-2{{c}^{2}}t\left( {{t}_{0}}' \right){{v}_{0}}^{2}+{{\left( c\left( {{t}_{0}}' \right){{v}_{0}} \right)}^{2}}+{{R}_{0}}^{2}\left( {{c}^{2}}-{{v}_{0}}^{2} \right)}}{({{c}^{2}}-{{v}_{0}}^{2})}{]}
{[}t'=\frac{-{{R}_{0}}{{v}_{0}}cos\left( \theta  \right)+{{c}^{2}}t-{{t}_{0}}'{{v}_{0}}^{2}-\sqrt{{{\left( cos\left( \theta  \right){{R}_{0}}{{v}_{0}} \right)}^{2}}-2cos\left( \theta  \right){{R}_{0}}{{v}_{0}}\left( t-{{t}_{0}}' \right){{c}^{2}}+{{c}^{2}}{{v}_{0}}^{2}{{\left( t-{{t}_{0}}' \right)}^{2}}+{{R}_{0}}^{2}\left( {{c}^{2}}-{{v}_{0}}^{2} \right)}}{({{c}^{2}}-{{v}_{0}}^{2})}{]}
Откуда {[}R=c\left( t-t' \right)=c\left(
t-\frac{-{{R}_{0}}{{v}_{0}}cos\left( \theta  \right)+{{c}^{2}}t-{{t}_{0}}'{{v}_{0}}^{2}-\sqrt{{{\left( cos\left( \theta  \right){{R}_{0}}{{v}_{0}} \right)}^{2}}-2cos\left( \theta  \right){{R}_{0}}{{v}_{0}}\left( t-{{t}_{0}}' \right){{c}^{2}}+{{c}^{2}}{{v}_{0}}^{2}{{\left( t-{{t}_{0}}' \right)}^{2}}+{{R}_{0}}^{2}\left( {{c}^{2}}-{{v}_{0}}^{2} \right)}}{({{c}^{2}}-{{v}_{0}}^{2})}
\right){]} Нормальная компонента электрического поля в приближении малых
ускорений \(a\to 0\)
\({{E}_{n}}=\frac{e}{{{K}^{3}}}\left( \overrightarrow{R}\frac{\overrightarrow{{{R}_{0}}}}{{{R}_{0}}}-R\frac{\overrightarrow{v}}{c}\frac{\overrightarrow{{{R}_{0}}}}{{{R}_{0}}} \right)\left( 1-\frac{\overrightarrow{v}\cdot \overrightarrow{v}}{{{c}^{2}}} \right)\)
Теперь необходимо определить радиус векторы. В декартовой системе
{[}\overrightarrow{R}=\overrightarrow{i}\left( x
\right)+\overrightarrow{j}\left( y \right)+\overrightarrow{k}\left(
z-\{\{z\}}\{q\}\} \right){]}
{[}\overrightarrow{{{R}_{0}}}=\overrightarrow{i}\left( x
\right)+\overrightarrow{j}\left( y \right)+\overrightarrow{k}\left( z
\right){]} где
\(x={{R}_{0}}\sin \left( \theta \right)\cos \left( \varphi \right)\)
\(y={{R}_{0}}\sin \left( \theta \right)\sin \left( \varphi \right)\)
\(z={{R}_{0}}\cos \left( \theta \right)\) а
{[}\{\{z\}\emph{\{q\}\}=\{\{z\}}\{0\}\}+\{\{v\}\emph{\{0\}\}\left(
t'-\{\{t\}}\{0\}\}' \right){]} Следовательно, с использованием
сферических координат эти векторы запишутся
{[}\overrightarrow{R}=\overrightarrow{i}{{R}_{0}}\sin \left(
\theta  \right)\cos \left(
\varphi  \right)+\overrightarrow{j}{{R}_{0}}\sin \left(
\theta  \right)\sin \left( \varphi  \right)+\overrightarrow{k}\left(
\{\{R\}\emph{\{0\}\}\cos \left( \theta  \right)-\{\{z\}}\{q\}\}
\right){]}
{[}\overrightarrow{{{R}_{0}}}=\overrightarrow{i}{{R}_{0}}\sin \left(
\theta  \right)\cos \left(
\varphi  \right)+\overrightarrow{j}{{R}_{0}}\sin \left(
\theta  \right)\sin \left(
\varphi  \right)+\overrightarrow{k}{{R}_{0}}\cos \left(
\theta  \right){]} Скалярные произведения
{[}\overrightarrow{R}\frac{\overrightarrow{{{R}_{0}}}}{{{R}_{0}}}=\{\{R\}\emph{\{0\}\}\sin {{\left( \theta  \right)}^{2}}\cos {{\left( \varphi  \right)}^{2}}+\{\{R\}}\{0\}\}\sin {{\left( \theta  \right)}^{2}}\sin {{\left( \varphi  \right)}^{2}}+\{\{R\}\emph{\{0\}\}\cos {{\left( \theta  \right)}^{2}}-\{\{z\}}\{q\}\}\cos \left(
\theta  \right){]}
{[}\overrightarrow{R}\frac{\overrightarrow{{{R}_{0}}}}{{{R}_{0}}}=\{\{R\}\emph{\{0\}\}\sin {{\left( \theta  \right)}^{2}}+\{\{R\}}\{0\}\}\cos {{\left( \theta  \right)}^{2}}-\{\{z\}\emph{\{q\}\}\cos \left(
\theta  \right){]}
{[}\overrightarrow{R}\frac{\overrightarrow{{{R}_{0}}}}{{{R}_{0}}}=\{\{R\}}\{0\}\}-\{\{z\}\emph{\{q\}\}\cos \left(
\theta  \right){]} Второе скалярное произведение, учитывая что
{[}\overrightarrow{v}=\{\{v\}}\{q\}\}\overrightarrow{k}{]}
\(R\frac{\overrightarrow{v}}{c}\frac{\overrightarrow{{{R}_{0}}}}{{{R}_{0}}}=R\frac{{{v}_{q}}}{c}\cos \left( \theta \right)=c\left( t-t' \right)\frac{{{v}_{q}}}{c}\cos \left( \theta \right)=\left( t-t' \right){{v}_{q}}\cos \left( \theta \right)\)
Сумма скалярных произведений {[}

\begin{align}
  & \left( \overrightarrow{R}\frac{\overrightarrow{{{R}_{0}}}}{{{R}_{0}}}-R\frac{\overrightarrow{v}}{c}\frac{\overrightarrow{{{R}_{0}}}}{{{R}_{0}}} \right)=\left\{ {{R}_{0}}-{{z}_{q}}\cos \left( \theta  \right) \right\}-\left\{ \left( t-t' \right){{v}_{q}}\cos \left( \theta  \right) \right\}= \\ 
 & ={{R}_{0}}-\left( {{z}_{q}}\cos \left( \theta  \right)+\left( t-t' \right){{v}_{q}}\cos \left( \theta  \right) \right)= \\ 
 & ={{R}_{0}}-\left( {{z}_{q}}+\left( t-t' \right){{v}_{q}} \right)\cos \left( \theta  \right) \\ 
\end{align}

{]} радиус Лиенара Вихерта
\(\overrightarrow{v}\cdot \overrightarrow{R}={{v}_{q}}\left( {{R}_{0}}\cos \left( \theta \right)-{{z}_{q}} \right)\)
{[}K=R-\frac{{{v}_{q}}\left( {{R}_{0}}\cos \left( \theta  \right)-{{z}_{q}} \right)}{c}=c\left(
t-t'
\right)-\frac{{{v}_{q}}\left( {{R}_{0}}\cos \left( \theta  \right)-{{z}_{q}} \right)}{c}{]}
Нормальная компонента поля
\({{E}_{n}}=\frac{e}{{{K}^{3}}}\left( \overrightarrow{R}\frac{\overrightarrow{{{R}_{0}}}}{{{R}_{0}}}-R\frac{\overrightarrow{v}}{c}\frac{\overrightarrow{{{R}_{0}}}}{{{R}_{0}}} \right)\left( 1-\frac{\overrightarrow{v}\cdot \overrightarrow{v}}{{{c}^{2}}} \right)\)
{[}\{\{E\}\emph{\{n\}\}=\frac{e}{{{K}^{3}}}\left(
\{\{R\}}\{0\}\}-\{\{z\}\emph{\{q\}\}\cos \left( \theta  \right)-\left(
t-t' \right)\{\{v\}}\{q\}\}\cos \left( \theta  \right) \right)\left(
1-\frac{{{v}_{q}}^{2}}{{{c}^{2}}} \right){]}
{[}\{\{E\}\emph{\{n\}\}=\frac{e}{{{K}^{3}}}\left( \{\{R\}}\{0\}\}-\left(
\{\{z\}\emph{\{q\}\}+\left( t-t' \right)\{\{v\}}\{q\}\}
\right)\cos \left( \theta  \right) \right)\left(
1-\frac{{{v}_{q}}^{2}}{{{c}^{2}}} \right){]} Интеграл по поверхности
сферы {[}\iint{{{E}_{n}}{{R}^{2}}}\sin \left(
\theta  \right)d\theta d\varphi =2\pi \int\limits\emph{\{\theta =0\}\textsuperscript{\{\pi \}\{\{\{E\}\emph{\{n\}\}\{\{R\}}\{0\}\}}\{2\}\}\sin \left(
\theta  \right)d\theta =2\pi e\int\limits}\{\theta =0\}\^{}\{\pi \}\{\frac{{{R}_{0}}-\left( {{z}_{q}}+\left( t-t' \right){{v}_{q}} \right)\cos \left( \theta  \right)}{{{K}^{3}}}\left(
1-\frac{{{v}_{0}}^{2}}{{{c}^{2}}}
\right)\{\{R\}\_\{0\}\}\^{}\{2\}\}\sin \left( \theta  \right){]}
Численный расчёт этого интеграла при z\_\_0 := 0.5 R0 := 2 tau\_\_0 = 0,
v\_\_0 = 0.9, c = 1

Следовательно, для равномерно движущегося заряда теорема Гаусса
выполняется в случае использования потенциалов Лиенара Вихерта

Более общий случай (векторы ускорения и скорости заряда сонаправлены).
{[}z\left( t',\{\{v\}\emph{\{0\}\},\{\{a\}}\{0\}\}
\right)=\{\{z\}\emph{\{0\}\}+\{\{v\}}\{0\}\}\left(
t'-\{\{t\}\emph{\{0\}\}'
\right)+\frac{{{a}_{0}}{{\left( t'-{{t}_{0}}' \right)}^{2}}}{2}{]}
{[}v\left( t',\{\{v\}}\{0\}\},\{\{a\}\emph{\{0\}\}
\right)=\{\{v\}}\{0\}\}+\{\{a\}\emph{\{0\}\}\left( t'-\{\{t\}}\{0\}\}'
\right){]} Расчёт запаздывающего момента \(t'\) производится с помощью
разрешения уравнения вида (это уравнение записано с помощью теоремы
косинусов) {[}\{\{R\}\textsuperscript{\{2\}\}=\{\{c\}}\{2\}\}\{\{\left(
t-t'
\right)\}\^{}\{2\}\}=\{\{R\}\emph{\{0\}\}\^{}\{2\}-2\{\{R\}}\{0\}\}z\left(
t',\{\{v\}\emph{\{0\}\},\{\{a\}}\{0\}\} \right)\cos \theta +z\{\{\left(
t',\{\{v\}\emph{\{0\}\},\{\{a\}}\{0\}\} \right)\}\^{}\{2\}\}{]}
Нормальная компонента электрического поля
\({{E}_{n}}=\overrightarrow{E}\frac{\overrightarrow{{{R}_{0}}}}{{{R}_{0}}}=\frac{e}{{{K}^{3}}}\left( \overrightarrow{R}\frac{\overrightarrow{{{R}_{0}}}}{{{R}_{0}}}-R\frac{\overrightarrow{v}}{c}\frac{\overrightarrow{{{R}_{0}}}}{{{R}_{0}}} \right)\left( 1+\frac{\overrightarrow{a}\cdot \overrightarrow{R}}{{{c}^{2}}}-\frac{\overrightarrow{v}\cdot \overrightarrow{v}}{{{c}^{2}}} \right)-\frac{e}{{{K}^{2}}}R\frac{\overrightarrow{a}}{{{c}^{2}}}\frac{\overrightarrow{{{R}_{0}}}}{{{R}_{0}}}\)
с использованием сферических координат векторы запишутся
{[}\overrightarrow{R}=\overrightarrow{i}{{R}_{0}}\sin \left(
\theta  \right)\cos \left(
\varphi  \right)+\overrightarrow{j}{{R}_{0}}\sin \left(
\theta  \right)\sin \left( \varphi  \right)+\overrightarrow{k}\left(
\{\{R\}\emph{\{0\}\}\cos \left( \theta  \right)-\{\{z\}}\{q\}\}\left( t'
\right) \right){]}
{[}\overrightarrow{{{R}_{0}}}=\overrightarrow{i}{{R}_{0}}\sin \left(
\theta  \right)\cos \left(
\varphi  \right)+\overrightarrow{j}{{R}_{0}}\sin \left(
\theta  \right)\sin \left(
\varphi  \right)+\overrightarrow{k}{{R}_{0}}\cos \left(
\theta  \right){]} Скалярные произведения

{[}\overrightarrow{R}\frac{\overrightarrow{{{R}_{0}}}}{{{R}_{0}}}=\{\{R\}\emph{\{0\}\}\sin {{\left( \theta  \right)}^{2}}\cos {{\left( \varphi  \right)}^{2}}+\{\{R\}}\{0\}\}\sin {{\left( \theta  \right)}^{2}}\sin {{\left( \varphi  \right)}^{2}}+\{\{R\}\emph{\{0\}\}\cos {{\left( \theta  \right)}^{2}}-\{\{z\}}\{q\}\}\left(
t' \right)\cos \left( \theta  \right){]}
{[}\overrightarrow{R}\frac{\overrightarrow{{{R}_{0}}}}{{{R}_{0}}}=\{\{R\}\emph{\{0\}\}\sin {{\left( \theta  \right)}^{2}}+\{\{R\}}\{0\}\}\cos {{\left( \theta  \right)}^{2}}-\{\{z\}\emph{\{q\}\}\left(
t' \right)\cos \left( \theta  \right){]}
{[}\overrightarrow{R}\frac{\overrightarrow{{{R}_{0}}}}{{{R}_{0}}}=\{\{R\}}\{0\}\}-\{\{z\}\emph{\{q\}\}\left(
t' \right)\cos \left( \theta  \right){]} Второе скалярное произведение,
учитывая что {[}\overrightarrow{v}=\{\{v\}}\{q\}\}\left( t'
\right)\overrightarrow{k}{]}
{[}R\frac{\overrightarrow{v}}{c}\frac{\overrightarrow{{{R}_{0}}}}{{{R}_{0}}}=R\frac{{{v}_{q}}\left( t' \right)}{c}\cos \left(
\theta  \right)=c\left( t-t'
\right)\frac{{{v}_{q}}\left( t' \right)}{c}\cos \left(
\theta  \right)=\left( t-t' \right)\{\{v\}\_\{q\}\}\left( t'
\right)\cos \left( \theta  \right){]} Сумма скалярных произведений
\(\left( \overrightarrow{R}\frac{\overrightarrow{{{R}_{0}}}}{{{R}_{0}}}-R\frac{\overrightarrow{v}}{c}\frac{\overrightarrow{{{R}_{0}}}}{{{R}_{0}}} \right)={{R}_{0}}-{{z}_{q}}\left( t' \right)\cos \left( \theta \right)-\left( t-t' \right){{v}_{q}}\left( t' \right)\cos \left( \theta \right)\)
Ещё
\(\frac{\overrightarrow{a}\cdot \overrightarrow{R}}{{{c}^{2}}}={{a}_{q}}\left( t' \right)\frac{\left( {{R}_{0}}\cos \left( \theta \right)-{{z}_{q}}\left( t' \right) \right)}{{{c}^{2}}}\)
И
\(\frac{e}{{{K}^{2}}}R\frac{\overrightarrow{a}}{{{c}^{2}}}\frac{\overrightarrow{{{R}_{0}}}}{{{R}_{0}}}=\frac{e}{{{K}^{2}}}R\frac{{{a}_{q}}\left( t' \right)}{{{c}^{2}}}\frac{{{R}_{0}}\cos \left( \theta \right)}{{{R}_{0}}}=\frac{e}{{{K}^{2}}}c\left( t-t' \right)\frac{{{a}_{q}}\left( t' \right)}{{{c}^{2}}}\cos \left( \theta \right)=\frac{e}{{{K}^{2}}}\left( t-t' \right)\frac{{{a}_{q}}\left( t' \right)}{c}\cos \left( \theta \right)\)
Теперь нормальная компонента поля
\({{E}_{n}}=e\frac{{{R}_{0}}-{{z}_{q}}\left( t' \right)\cos \left( \theta \right)-\left( t-t' \right){{v}_{q}}\left( t' \right)\cos \left( \theta \right)}{{{K}^{3}}}\left( 1+{{a}_{q}}\left( t' \right)\frac{\left( {{R}_{0}}\cos \left( \theta \right)-{{z}_{q}}\left( t' \right) \right)}{{{c}^{2}}}-\frac{{{v}_{q}}{{\left( t' \right)}^{2}}}{{{c}^{2}}} \right)-\frac{e}{{{K}^{2}}}\left( t-t' \right)\frac{{{a}_{q}}\left( t' \right)}{c}\cos \left( \theta \right)\)
радиус Лиенара Вихерта
\(\overrightarrow{v}\cdot \overrightarrow{R}={{v}_{q}}\left( t' \right)\left( {{R}_{0}}\cos \left( \theta \right)-{{z}_{q}}\left( t' \right) \right)\)
{[}

\begin{align}
  & K=R-\frac{{{v}_{q}}\left( t' \right)\left( {{R}_{0}}\cos \left( \theta  \right)-{{z}_{q}}\left( t' \right) \right)}{c}= \\ 
 & =c\left( t-t' \right)-\frac{{{v}_{q}}\left( t' \right)\left( {{R}_{0}}\cos \left( \theta  \right)-{{z}_{q}}\left( t' \right) \right)}{c} \\ 
\end{align}

{]}

О применимости калибровки Лоренца к электродинамике Николаева А.Ю.
Дроздов 632Equation Section (Next) Одним из возможных объяснений ЭМИ
сферически симметричного взрыва является индукция Николаева, записанная
им {[}CITATION ГНи \l 1049 {]} в следующем виде
{[}div~E=\frac{1}{c}\frac{\partial {{H}_{\parallel }}}{\partial t}=-\frac{1}{c}\frac{\partial }{\partial t}div~\overrightarrow{{{A}_{H}}}{]}
64264* MERGEFORMAT (.) Но для совпадения предсказаний теории с
результатами эксперимента необходимо дополнить данное Николаевым
уравнение продольной электрической индукции множителем
\({{\mu }_{\parallel }}\) который по аналогии с параметром \$\mu \$ из
классической электродинамики можно было бы назвать продольной магнитной
проницаемостью вещества или вакуума. Причем, продольная магнитная
проницаемость вакуума \({{\mu }_{\parallel }}\) должна отличаться от
обычной классической магнитной проницаемости вакуума по меньшей мере
знаком. С точки зрения электродинамической теории полученные результаты
указывают на необходимость введения в четвертое уравнение Максвелла
дополнительного слагаемого. Действительно, уравнение
\$div~\varepsilon \overrightarrow{E}=4\pi \rho \$известное из курса
электростатики при применении теоремы Гаусса к рассматриваемой задаче
центрально-симметричного взрыва изначально электро-нейтральной плазмы
даёт в итоге нулевую радиальную компоненту электрического поля на
поверхности большого радиуса {[}\{\{R\}\emph{\{O\}\}{]}. Это уравнение
должно быть модифицировано следующим образом
\(div\ \varepsilon \overrightarrow{E}=4\pi \rho +4\pi {{\rho }_{}}=4\pi \rho -\frac{\varepsilon {{\mu }_{\parallel }}}{c}\frac{\partial \ }{\partial t}div\overrightarrow{{{A}_{H}}}\)
65265* MERGEFORMAT (.) Итак, четвёртое уравнение Максвелла для
удовлетворительного объяснения явлений ЭМИ ядерных и иных
центрально-симметричных взрывов плазмы должно иметь вид
{[}divD=4\pi \rho +\frac{\varepsilon {{\mu }_{\parallel }}}{c}\frac{\partial {{H}_{\parallel }}}{\partial t}=4\pi \rho +\frac{\varepsilon }{c}\frac{\partial {{B}_{\parallel }}}{\partial t}{]}
66266* MERGEFORMAT (.) где
\({{B}_{\parallel }}={{\mu }_{\parallel }}{{H}_{\parallel }}\) Наряду с
индукцией Николаева в электродинамике широко известна индукция Фарадея,
выражающаяся во втором уравнении Максвелла
{[}rot~E=-\frac{1}{c}\frac{\partial {{B}_{\bot }}}{\partial t}=-\frac{\mu }{c}\frac{\partial {{H}_{\bot }}}{\partial t}{]}
67267* MERGEFORMAT (.) Выражение электрического поля через векторный и
скалярные потенциалы, удовлетворяющее второму уравнению Максвелла, а
значит и явлению Фарадеевской индукции имеет вид
{[}E=-\frac{\mu }{c}\frac{\partial {{A}_{H}}}{\partial t}-grad\varphi {]}
68268* MERGEFORMAT (.) Взяв дивергенцию этого выражения получаем
{[}div~E=-\frac{\mu }{c}div\frac{\partial {{A}_{H}}}{\partial t}-div~grad\varphi {]}
69269* MERGEFORMAT (.) Откуда
{[}div~D=\frac{\varepsilon \mu }{c}\frac{\partial {{H}_{\parallel }}}{\partial t}-\varepsilon {{\nabla }^{2}}\varphi {]}
70270* MERGEFORMAT (.) Приравняв полученное выражение к дополненному
индукцией Николаева четвёртому уравнению Максвелла 266, получаем
волновое уравнение
{[}\frac{\varepsilon \left( \mu -{{\mu }_{\parallel }} \right)}{c}\frac{\partial {{H}_{\parallel }}}{\partial t}-\varepsilon {{\nabla }^{2}}\varphi =4\pi \rho {]}
71271* MERGEFORMAT (.) Из полученного уравнения видно, что если бы
продольная и обычная магнитная проницаемость не отличались бы друг от
друга знаком, то скалярный потенциал не зависел бы от времени, что
означало бы невозможность получения волнового уравнения без махинаций с
калибровками. Исходя из вышеизложенного, мы можем дополнить рассуждения
Е.И. Тамма, которые он приводит при выводе понятия ток смещения, «исходя
из убеждения в правильности уравнения неразрывности»
{[}div~j=-\frac{\partial }{\partial t}\rho {]} 72272* MERGEFORMAT (.)
Запишем первое уравнение Максвелла в виде
{[}rot~H=\frac{4\pi }{c}{{j}_{}}=\frac{4\pi }{c}\left( j+\{\{j\}}\{\}\}
\right){]} 73273* MERGEFORMAT (.) Возьмём от него дивергенцию
{[}0=div~rot~H=\frac{4\pi }{c}div~\{\{j\}\emph{\{\}\}=\frac{4\pi }{c}div~j+\frac{4\pi }{c}div~\{\{j\}}\{\}\}{]}
74274* MERGEFORMAT (.) Далее
{[}div~\{\{j\}\emph{\{\}\}=-div~j=\frac{\partial }{\partial t}\rho {]}
75275* MERGEFORMAT (.) Но исходя из дополненного индукцией Николаева
четвёртого уравнения Максвелла 266 {[}\rho =\frac{1}{4\pi }\left(
div~D-\frac{\varepsilon {{\mu }_{\parallel }}}{c}\frac{\partial {{H}_{\parallel }}}{\partial t}
\right){]} 76276* MERGEFORMAT (.) Откуда для тока смещения запишем
{[}div~\{\{j\}}\{\}\}=\frac{1}{4\pi }\frac{\partial }{\partial t}\left(
div~D-\frac{\varepsilon {{\mu }_{\parallel }}}{c}\frac{\partial {{H}_{\parallel }}}{\partial t}
\right){]} 77277* MERGEFORMAT (.) Выразим \(div\ D\), используя
уравнение 270
{[}div~\{\{j\}\emph{\{\}\}=\frac{1}{4\pi }\frac{\partial }{\partial t}\left(
\frac{\varepsilon \mu }{c}\frac{\partial {{H}_{\parallel }}}{\partial t}-\varepsilon {{\nabla }^{2}}\varphi -\frac{\varepsilon {{\mu }_{\parallel }}}{c}\frac{\partial {{H}_{\parallel }}}{\partial t}
\right)=\frac{\varepsilon \left( \mu -{{\mu }_{\parallel }} \right)}{4\pi c}\frac{{{\partial }^{2}}{{H}_{\parallel }}}{\partial {{t}^{2}}}-\frac{\varepsilon }{4\pi }{{\nabla }^{2}}\frac{\partial }{\partial t}\varphi {]}
78278* MERGEFORMAT (.) Это выражение может быть удовлетворено, если
положить для тока смещения
{[}\{\{j\}}\{\}\}=-\frac{\varepsilon \left( \mu -{{\mu }_{\parallel }} \right)}{4\pi c}\frac{{{\partial }^{2}}{{}_{H}}}{\partial {{t}^{2}}}-\frac{\varepsilon }{4\pi }\nabla \frac{\partial }{\partial t}\varphi +c~rot~I{]}
79279* MERGEFORMAT (.) Подставляя это выражение в первое уравнение
Максвелла 273 получаем
{[}rot~H=\frac{4\pi }{c}j-\frac{\varepsilon \left( \mu -{{\mu }_{\parallel }} \right)}{{{c}^{2}}}\frac{{{\partial }^{2}}{{}_{H}}}{\partial {{t}^{2}}}-\frac{\varepsilon }{c}\nabla \frac{\partial }{\partial t}\varphi +4\pi ~rot~I{]}
80280* MERGEFORMAT (.) Составим из первого уравнения Максвелла 280
волновое уравнение
{[}rot~rot~\{\{A\}\emph{\{H\}\}=grad~div~\{\{A\}}\{H\}\}-\{\{\nabla \}\^{}\{2\}\}\{\{A\}\emph{\{H\}\}=\frac{4\pi }{c}j-\frac{\varepsilon \left( \mu -{{\mu }_{\parallel }} \right)}{{{c}^{2}}}\frac{{{\partial }^{2}}{{}_{H}}}{\partial {{t}^{2}}}-\frac{\varepsilon }{c}\nabla \frac{\partial }{\partial t}\varphi +4\pi ~rot~I{]}
81281* MERGEFORMAT (.) Таким образом, вместе с
{[}\frac{\varepsilon \left( \mu -{{\mu }_{\parallel }} \right)}{c}\frac{\partial }{\partial t}div~\{\{A\}}\{H\}\}+\varepsilon {{\nabla }^{2}}\varphi =-4\pi \rho {]}
82282* MERGEFORMAT (.) Получаем систему волновых уравнений Применив
калибровку Лоренца
{[}div~\{\{A\}\emph{\{H\}\}=-\frac{\varepsilon }{c}\frac{\partial \varphi }{\partial t}{]}
83283* MERGEFORMAT (.) Получаем
{[}\{\{\nabla \}\^{}\{2\}\}\varphi -\frac{\varepsilon \left( \mu -{{\mu }_{\parallel }} \right)}{{{c}^{2}}}\frac{{{\partial }^{2}}\varphi }{\partial {{t}^{2}}}=-\frac{4\pi \rho }{\varepsilon }{]}
84284* MERGEFORMAT (.)
{[}\{\{\nabla \}\^{}\{2\}\}\{\{A\}}\{H\}\}-\frac{\varepsilon \left( \mu -{{\mu }_{\parallel }} \right)}{{{c}^{2}}}\frac{{{\partial }^{2}}{{A}_{H}}}{\partial {{t}^{2}}}=-\frac{4\pi j}{c}-4\pi ~rot~I{]}
85285* MERGEFORMAT (.) Полученные волновые уравнения отличаются от
классических волновых уравнений {[} CITATION Там57 \l 1049 {]} учетом
того обстоятельства, что магнитная проницаемость вакуума, проявляющаяся
в явлениях индукции Фарадея и индукции Николаева должна отличаться. Надо
сказать, что в данной работе благодаря применению двух различающихся
коэффициентов магнитной проницаемости при получении волновых уравнений
Даламбера не было необходимости жонглировать калибровками, в отличие от
выкладок Тамма Кроме того: формулы 14 и 15 приводятся в курсах
электродинамики как решения уравнений Даламбера. Однако в рамках
настоящей работы эти уравнения могут быть получены как решения уравнений
284 и 285 поэтому магнитная проницаемость в формуле 15 должна быть взята
как разность {[}\left( \mu -\{\{\mu \}\_\{\parallel \}\}
\right){]}магнитных проницаемостей из законов индукции Фарадея и
индукции Николаева. Итак, в данной работе мне удалось избавиться от
необходимости жонглирования калибровками, которая имеет место в
классической электродинамике, при выводе волнового уравнения Даламбера.
Но является ли мой вывод полностью непротиворечивым? Оказывается, нет.
Если внимательно присмотреться к уравнениям 265 и 268 то можно увидеть,
что в этих выражениях по разному выглядит составляющая электрического
поля связанная с явлением электромагнитной индукции. В индукции
Николаева
\(\overrightarrow{E}=-\frac{{{\mu }_{\parallel }}}{c}\frac{\partial }{\partial t}\overrightarrow{{{A}_{H}}}\)причём
\({{\mu }_{\parallel }}<0\)а в индукции Фарадея
\(\overrightarrow{E}=-\frac{\mu }{c}\frac{\partial }{\partial t}\overrightarrow{{{A}_{H}}}\).
Пытаясь привести эти два выражения к общему виду можно предположить, что
магнитная проницаемость вакуума является знакопеременной функцией угла
между переменным током -- источником изменяющегося во времени векторного
потенциала и радиус вектором, проведённым к точке измерения значения
электромагнитной индукции. Однако в рамках представленного материала
более физически обоснованной представляется гипотеза, что на самом деле
причинной явления электромагнитной индукции является не изменяющийся во
времени векторный потенциал, а изменяющийся во времени ток смещения.

Список литературы Apfelbaum, E. a. (2009). The predictions of the
critical point parameters for Al, Cu and W found from the correspondence
between the critical point and unit compressibility line (Zeno line)
positions. Chemical Physics Letters(467), pp. 318--322. Retrieved from
https://www.sciencedirect.com/science/article/abs/pii/S0009261408015790
Baker, D. A. (1964). Second-Order Electric Field due to a Conducting
Curent. American Journal of Physics, 32(2), 153-157. Edwards, W. F.
(1976). Continuing investigation into possible electric arising from
steady conduction current. Phys. Rev. D, 14(4), 922-\/-938.
doi:10.1103/PhysRevD.14.922 Mende, F. F. (1993). Conference ``Physics in
Ukraine''. Experimental corroboration and theoretical interpretation of
dependence of charge value on DC flow velocity through superconductors.
Kiev. Roser, W. (1962). Second-Order Electric Field due to a Conducting
Curent. American Journal of Physics, 30(7), 509-511. Адамьян Ю.Э., В. В.
(1999). Нагрев и ускорение плазмы при взрыве проводника в вакууме в
сильном магнитном поле. Журнал технической физики, 69(5). Добровольская
А.С., К. Н. (2015). Применение дискретной кинетической модели для
системы уравнений Эйлера - Пуассона к задачам плазменной аэродинамики.
Физико-химическая кинетика в газовой динамике, 16(3). Получено из
http://chemphys.edu.ru/issues/2015-16-3/articles/571/ Зиновьев, В.
(1989). Теплофизические свойства металлов при высоких температурах.
«Металлургия». Лифшиц, Л. Л. (1973). Теоретическая физика. В Теория поля
(Т. 2). Москва. Менде, Ф. Ф. (2013). Электрический импульс космического
термоядерного взрыва. Инженерная физика, 16-24. Менде, Ф. Ф. (2015).
Является ли заряд инвариантом скорости? Получено из
http://fmnauka.narod.ru/javljaetsja\_li\_zarjad\_invariantom\_skorosti.pdf
Николаев, Г. (2004). Электродинамика физического вакуума. Томск: НТЛ.
Сасакин, М. (2016). Разлёт в вакуум и взаимодействие с преградой
мелкодисперсных частиц при электрическом взрыве проводника.
Санкт-Петербург. Тамм, И. (1957). Основы теории электричества. Москва.
Чиркин, В. (1968). Теплофизические свойства материалов ядерной техники.
Справочник. М: Атомиздат.


    % Add a bibliography block to the postdoc
    
    
    
\end{document}
